\documentclass[9pt,twoside]{extarticle}
\usepackage[a4paper,top=14mm,bottom=14mm,left=19mm,right=19mm,headheight=12pt,headsep=3.5mm,footskip=7mm,includeheadfoot]{geometry}
\usepackage{fontspec}
\defaultfontfeatures{Ligatures=TeX,Scale=MatchLowercase}
\setmainfont{Libertinus Serif}
\setsansfont{Libertinus Sans}
\usepackage{mathtools}
\usepackage{amssymb}
\usepackage{array}
\usepackage{unicode-math}
\setmathfont{Libertinus Math}
\usepackage[english]{babel}
\usepackage{microtype}
\usepackage{adjustbox}
\usepackage{fancyhdr}
\usepackage{ragged2e}
\usepackage[hidelinks,unicode]{hyperref}
\usepackage{bookmark}
\hypersetup{pdftitle={Representations of Reductive Groups over Finite Fields},pdfauthor={P. Deligne and G. Lusztig},pdfsubject={Source-language critical edition},pdfkeywords={Deligne, Lusztig, reductive groups, finite fields, editable LaTeX}}
\pagestyle{fancy}
\fancyhf{}
\fancyhead[LE,RO]{\footnotesize\sffamily Reductive groups over finite fields - English source edition}
\fancyhead[LO,RE]{\footnotesize\sffamily Printed page \thepage}
\fancyfoot[C]{\footnotesize\thepage}
\renewcommand{\headrulewidth}{0.25pt}
\renewcommand{\footrulewidth}{0pt}
\setlength{\parindent}{1em}
\setlength{\parskip}{0.11em}
\setlength{\abovedisplayskip}{2.8pt plus 1pt minus 1pt}
\setlength{\belowdisplayskip}{2.8pt plus 1pt minus 1pt}
\setlength{\abovedisplayshortskip}{1.8pt}
\setlength{\belowdisplayshortskip}{1.8pt}
\setlength{\jot}{1.2pt}
\emergencystretch=1.5em
\newcommand{\EditionHeading}[1]{\par\medskip{\noindent\centering\large\bfseries #1\par}\smallskip}
\newcommand{\EditionChapter}[1]{\par\medskip{\noindent\large\bfseries #1\par}\smallskip}
\newcommand{\EditionPage}[3]{%
  \ifnum#1>2\clearpage\fi
  \pdfbookmark[1]{Printed page #2}{en-folio-#2}%
  \hypertarget{en-physical-#1}{}%
  \noindent\begin{adjustbox}{max totalsize={\textwidth}{0.918\textheight},center}
  \begin{minipage}[t]{\textwidth}
  \fontsize{9.4}{11.1}\selectfont\justifying
  #3
  \end{minipage}
  \end{adjustbox}%
}
\begin{document}
\setcounter{page}{103}
\pdfbookmark[0]{Representations of Reductive Groups over Finite Fields}{d027-title-en}

\EditionPage{2}{103}{%
\begin{center}\footnotesize Annals of Mathematics, 103 (1976), 103-161\end{center}

\begin{center}\fontsize{14.5}{16.5}\selectfont\bfseries Representations of reductive groups over finite fields\end{center}

\begin{center}\large By P. DELIGNE and G. LUSZTIG\end{center}

\EditionHeading{Introduction}

Let us consider a connected, reductive algebraic group \(G\), defined over a finite field \(\mathbb{F}_q\), with Frobenius map \(F\). We shall be concerned with the representation theory of the finite group \(G^F\), over fields of characteristic 0.\par

In 1968, Macdonald conjectured, on the basis of the character tables known at the time (\(GL_n\), \(Sp_4\)), that there should be a well defined correspondence which, to any \(F\)-stable maximal torus \(T\) of \(G\) and a character \(\theta\) of \(T^F\) in general position, associates an irreducible representation of \(G^F\); moreover, if \(T\) modulo the centre of \(G\) is anisotropic over \(\mathbb{F}_q\), the corresponding representation of \(G^F\) should be cuspidal (see Seminar on algebraic groups and related finite groups, by A. Borel et al., Lecture Notes in Mathematics, 131, pp. 117 and 101). In this paper we prove Macdonald's conjecture. More precisely, for \(T\) as above and \(\theta\) an arbitrary character of \(T^F\) we construct virtual representations \(R_T^\theta\) which have all the required properties.\par

These are defined in Chapter 1 as the alternating sum of the cohomology with compact support of the variety of Borel subgroups of \(G\) which are in a fixed relative position with their \(F\)-transform, with coefficients in certain \(G^F\)-equivariant locally constant \(l\)-adic sheaves of rank one. This generalizes a construction made by Drinfeld for \(SL_2\) (see Ch. 2).\par

In Chapter 4 we prove a character formula for \(R_T^\theta\), based on the fixed point formula of Chapter 3. This character formula is in terms of certain “Green functions” on the unipotent elements; in Chapters 6, 7, 8 we prove that these Green functions satisfy all the axioms predicted by Springer, Kilmoyer and Macdonald ([9], [12]).\par

By 6.8, \(\pm R_T^\theta\) is irreducible if \(\theta\) is in general position and the vanishing theorem (9.9) gives an explicit model for it provided that \(q\) is not too small (if \(G\) is a classical group or \(G_2\) any \(q\) will do; in the general case \(q \geq 30\) is sufficient).\par

In Chapter 10 we study the irreducible components of the Gelfand-Graev representation of \(G^F\), assuming that the centre of \(G\) is connected. The proof uses the results of Chapter 5 together with the disjointness theorem (6.2).\par
}

\EditionPage{3}{104}{%
Finally, in Chapter 11 we discuss the case of the Suzuki and Ree groups.\par

It would be very desirable to find formulas for the Green functions more explicit than (4.1.2). Such formulas are known for \(GL_n\) (Green [4]), \(Sp_4\) (Srinivasan [13]) and \(G_2\) (Chang, Ree [2]). Very recently, Kazhdan has proved, using results of Springer, that the Green functions can be expressed as exponential sums on the Lie algebra (see [12]) provided that the characteristic is good and \(q\) is not too small.\par

Some of the results in this paper were announced in [7], [8].\par

The second author would like to thank the I.H.E.S. for its hospitality during part of the time of preparation of this paper.\par

\EditionHeading{TABLE OF CONTENTS}

Conventions and standard notations ........................ 104 Chapter 1. Some basic definitions ............................... 105 2. Examples in the classical groups ..................... 116 3. A fixed point formula ................................ 118 4. The character formula ................................ 123 5. Characters of tori ................................... 126 6. Intertwining numbers ................................. 135 7. Computations on semisimple elements .................. 140 8. Induced and cuspidal representations ................. 146 9. A vanishing theorem .................................. 147 10. A decomposition of the Gelfand-Graev representation .. 154 11. Suzuki and Ree groups ................................ 160\par

Conventions and standard notations\par

0.1. In this paper, \(k\) is always an algebraically closed field and \(p\) is its characteristic exponent. The following assumptions and definitions will be in force in Chapters 4 to 8, in Chapter 10, and in parts of Chapters 1 and 9.\par

(0.1.1) \(p\) is a prime, and \(k\) is an algebraic closure of the prime field \(\mathbb{F}_p\) of characteristic \(p\). For \(q\) a power of \(p\), \(\mathbb{F}_q\) is the subfield with \(q\) elements of \(k\). \(G\) is a reductive algebraic group over \(k\), obtained by extension of scalars from \(G_0\) over \(\mathbb{F}_q\). We denote by \(F\) the corresponding Frobenius endomorphism \(F:G\to G\) (as well as the Frobenius endomorphism of any \(k\)-scheme defined over \(\mathbb{F}_q\)).\par

0.2. \(l\) is a prime different from \(p\), and \(\overline{\mathbb{Q}}_l\) is an algebraic closure of the \(l\)-adic field. When there is no ambiguity on \(l\), we will write simply \(H_c^i(X)\) (resp. \(H^i(X)\)) for \(H_c^i(X,\overline{\mathbb{Q}}_l)\) (resp. \(H^i(X,\overline{\mathbb{Q}}_l)\)) (\(l\)-adic cohomology of \(X\), a scheme over \(k\)). The groups \(\mathbb{Z}/l^n(1)\) are the groups \(\mu_{l^n}(k)\) of \(l^n\)-roots of unity in \(k\); they form a projective system, with transition maps \(x\mapsto x^{l^{n-m}}\); the projective limit is \(\mathbb{Z}_l(1)\). One defines \(\mathbb{Z}_l(n)=\mathbb{Z}_l(1)^{\otimes n}\), and \(\mathbb{Z}_l(-n)\) the dual of \(\mathbb{Z}_l(n)\). The symbol \((n)\) is for a Tate twist: tensoring with \(\mathbb{Z}_l(n)\). We will\par
}

\EditionPage{4}{105}{%
make an accidental use of the similarly defined group \(\widehat{\mathbb{Z}}_{p'}(1)=\varprojlim_{(n,p)=1}\mu_n(k)\).\par

0.3. For \(H\) a finite group, \(\mathcal{R}(H)\) is the Grothendieck group of the finite dimensional representations of \(H\) over \(\overline{\mathbb{Q}}_l\). If \(f\) and \(f'\) are class functions on \(H\), with values in a cyclotomic field (often \(\subset\overline{\mathbb{Q}}_l\)), the inner product \(\langle f,f'\rangle_H\) (or simply \(\langle f,f'\rangle\)) is \[
\frac{1}{|H|}\sum_{x\in H} f(x)\,\overline{f'(x)},
\] where the bar is the automorphism inducing \(\zeta\mapsto\zeta^{-1}\) on the roots of unity. The same notation \(\langle\ ,\ \rangle\) will be used for the inner product of elements of \(\mathcal{R}(H)\otimes\mathbb{Q}\), identified with their characters.\par

0.4. The length function on a Coxeter group \(W\) (with canonical generators \(s_1,\ldots,s_n\)) will be denoted by \(l(\ )\). A reduced expression for \(w\in W\) is a decomposition \(w=s_{i_1}\cdots s_{i_k}\) with \(l(w)=k\).\par

0.5. Miscellaneous:\par

\(\operatorname{pr}_i\): \(i\)th projection (as in \(\operatorname{pr}_1:X\times X\to X\));\par

\(X^T\): the fixed point subscheme of the endomorphism \(T\) of the scheme \(X\) (for instance: \(G^F=G_0(\mathbb{F}_q)\));\par

\(p'\) (in index): away from \(p\) (as in \(\widehat{\mathbb{Z}}_{p'}\), completion away from \(p\)) or the subgroup of elements of order prime to \(p\) (as in \((\mathbb{Q}/\mathbb{Z})_{p'}\));\par

\(\operatorname{ad}g\): the inner automorphism \(x\mapsto gxg^{-1}\), or maps deduced from it;\par

\(1\), or \(e\): the identity element of a group;\par

\(\operatorname{Tr}(f,V^*)\), for \(f\) an endomorphism of a graded vector space, is \[
\sum(-1)^i\operatorname{Tr}(f,V^i);
\]\par

reductive: reductive groups are meant to be connected and smooth.\par

The Jordan decomposition of an element \(g\in G\) (\(G\) as in (0.1.1)) is \(g=su\) where \(s\) is a semisimple element and \(u\) is a unipotent element commuting with \(s\).\par

0.6. We will often identify a scheme over \(k\) with the set of its \(k\)-rational points. This should cause no confusion.\par

\EditionChapter{1. Some basic definitions}

1.1. Suppose that in some category we are given a family \((X_i)_{i\in I}\) of objects and a compatible system of isomorphisms \(\varphi_{ji}:X_i\xrightarrow{\sim}X_j\). This is as good as giving a single object \(X\), the “common value” or “projective limit” of the family. This projective limit is provided with isomorphisms \(\sigma_i:X\xrightarrow{\sim}X_i\) such that \(\varphi_{ji}\sigma_i=\sigma_j\). We will use such a construction to define the maximal torus \(\mathbf{T}\) and the Weyl group \(\mathbf{W}\) of a connected reductive algebraic group \(G\) over \(k\).\par

As index set \(I\), we take the set of pairs \((T,B)\) consisting of a maximal torus \(T\) and a Borel subgroup \(B\) containing \(T\). For \(i\in I\), \(i=(T,B)\), we take\par
}

\EditionPage{5}{106}{%
\(T_i=T,\quad W_i=N(T)/T.\)\par

The isomorphism \(\varphi_{ji}\) is the isomorphism induced by \(\operatorname{ad}g\) where \(g\) is any element of \(G\) conjugating \(i\) into \(j\); these elements \(g\) form a single right \(T_i\)-coset, so that \(\varphi_{ji}\) is independent of the choice of \(g\).\par

One similarly defines the root system of \(\mathbf{T}\), its set of simple roots, the action of \(\mathbf{W}\) on \(\mathbf{T}\) and the fundamental reflections in \(\mathbf{W}\).\par

Let \(F:G\to G\) be an isogeny. For any \(i\in I\), \(i=(T,B)\), \(F(i)=(F(T),F(B))\) is again in \(I\); \(F\) induces an isogeny \(F:T_i\to T_{F(i)}\) and an isomorphism \(F:W_i\to W_{F(i)}\). The corresponding endomorphism (resp. automorphism) \(\sigma_{F(i)}^{-1}F\sigma_i\) of the torus \(\mathbf{T}\) (resp. of the Weyl groups \(\mathbf{W}\)) is independent of the choice of \(i\); we say that it is induced by \(F\).\par

1.2. Let \(X\) (or \(X_G\)) be the set of all Borel subgroups of \(G\). The group \(G\) acts on \(X\) by conjugation, and \(X\) is a smooth projective homogeneous space of \(G\). For each Borel subgroup \(B\) of \(G\), \(B\) is the stabilizer of the corresponding point of \(X\), hence there is a natural isomorphism \(G/B\to X:g\mapsto gBg^{-1}\).\par

The set of orbits of \(G\) in \(X\times X\) can be identified with the Weyl group \(\mathbf{W}\) of \(G\) as follows: for any \(i=(T,B)\) as in (1.1), use the composite bijection \[
\mathbf{W}\xrightarrow[\sigma_i]{\sim}N(T)/T\xrightarrow{\sim}B\backslash G/B\xrightarrow[(e,g)]{\sim}G\backslash(G/B\times G/B)\xrightarrow{\sim}G\backslash(X\times X);
\] this is independent of the choice of \((T,B)\). We denote by \(O(w)\) the orbit corresponding to \(w\in\mathbf{W}\); this is the orbit of \((B,\dot{w}B\dot{w}^{-1})\), where \(\dot{w}\in N(T)\) represents \(w\). We shall say that two Borel subgroups \(B'\), \(B''\) of \(G\) are in relative position \(w\), \(w\in\mathbf{W}\), if and only if \((B',B'')\in O(w)\). In diagrams, we will picture this by \[
B'\xrightarrow{w}B''.
\]\par

The basic properties of Bruhat decomposition can be expressed as follows:\par

(a) If \(w=w_1w_2\), with \(l(w)=l(w_1)+l(w_2)\) then\par

(a₁) \((B',B'')\in O(w_1)\) and \((B'',B''')\in O(w_2)\)  \(\Rightarrow (B',B''')\in O(w)\);\par

(a₂) if \((B',B''')\in O(w)\), there is one and only one \(B''\) such that \((B',B'')\in O(w_1)\) and \((B'',B''')\in O(w_2)\).\par

On the scheme level: \[
O(w_1)\times_X O(w_2)\xrightarrow{\sim}O(w).
\]\par

(b) Let \(s\) be an elementary reflection and let \(B\) be a Borel subgroup.\par

(b₁) \(P=B\cup BsB\) is a (minimal parabolic) subgroup, i.e., if \((B',B'')\in O(s)\) and \((B'',B''')\in O(s)\), then either \((B',B''')\in O(s)\) or \(B'=B'''\).\par

(b₂) The quotient \(\overline{L}\) of \(L=P/U_P\) (\(U_P\) = unipotent radical of \(P\)) by its centre is isomorphic to \(PGL(2)\). The space \(X_{\overline{L}}\) is hence a projective line. The inverse image map \(X_{\overline{L}}\to X\) has as image the set of Borel subgroups \(B'\) in relative position \(e\) or \(s\) with \(B\).\par

(b₃) Via either projection, \(\overline{O(s)}=O(s)\cup O(e)\subset X\times X\) is hence a fibre\par
}

\EditionPage{6}{107}{%
space over \(X\), with fibre \(P_1\). It is provided with the section \(O(e)\); this reduces its structural group from the projective to the affine group. The complement \(O(s)\) of that section is a fibre space with fibre the affine line.\par

1.3. The assumptions (0.1.1) are in force in the rest of this chapter. The scheme \(X\) is hence defined over \(\mathbb{F}_q\) and provided with a Frobenius map \(F:X\to X\).\par

\noindent\textbf{DEFINITION 1.4.} For \(w\) in the Weyl group \(\mathbf{W}\) of \(G\), \(X(w)\subset X\) is the locally closed subscheme of \(X\) consisting of all Borel subgroups \(B\) of \(G\) such that \(B\) and \(F(B)\) are in relative position \(w\).\par

One can also regard \(X(w)\) as the intersection, in \(X\times X\), of \(O(w)\) with the graph of Frobenius. It is easily checked that this intersection is transverse. The orbit \(O(w)\) being smooth of dimension \(\dim(X)+l(w)\), it follows that \(X(w)\) is smooth and purely of dimension \(l(w)\). The subscheme \(X(w)\) of \(X\) is \(G^F\)-stable. Hence, for each prime number \(l\ne p\), \(G^F\) acts on the \(l\)-adic cohomology with compact supports of \(X(w)\).\par

\noindent\textbf{DEFINITION 1.5.} \(R^1(w)\) is the virtual representation \[
\sum(-1)^i H_c^i(X(w),\overline{\mathbb{Q}}_l)
\] of \(G^F\) (an element of the Grothendieck group of representations of \(G^F\) over \(\overline{\mathbb{Q}}_l\)).\par

For \(w=e\), \(X(w)\) is of dimension 0; it is the set of rational Borel subgroups, and \(R^1(w)\) is induced by the unit representation of \(B^F\), where \(B\) is an \(F\)-stable Borel subgroup.\par

It will follow from (3.3) that the character of \(R^1(w)\) has integral values, independent of \(l\). This will justify omitting \(l\) from the notation; \(R^1(w)\) could be also regarded as a complex virtual representation of \(G^F\).\par

An element of the Weyl group \(\mathbf{W}\) is said to be \(F\)-conjugate to \(w\in\mathbf{W}\), if it is of the form \(w_1wF(w_1)^{-1}\), for some \(w_1\in\mathbf{W}\). We denote by \(W_F^{\sim}\) the set of \(F\)-conjugacy classes in \(\mathbf{W}\). We shall recall in (1.14), that the \(G^F\)-conjugacy classes of \(F\)-stable (i.e., \(\mathbb{F}_q\)-rational) maximal tori in \(G\) are parametrized by \(W_F^{\sim}\).\par

\noindent\textbf{THEOREM 1.6.} \(R^1(w)\) depends only on the \(F\)-conjugacy class of \(w\).\par

Let \(w\) and \(w'\) be \(F\)-conjugate.\par

Case 1. \(w=w_1w_2\), \(w'=w_2F(w_1)\) and \(l(w)=l(w_1)+l(w_2)=l(w_2)+l(F(w_1))=l(w')\) (0.4). For \(B\in X(w)\), there is a unique Borel subgroup \(\sigma B\) such that \((B,\sigma B)\in O(w_1)\) and \((\sigma B,FB)\in O(w_2)\); we have the diagram of relative positions:\par
}

\EditionPage{7}{108}{%
\[
\begin{array}{ccc}
B & \xrightarrow{w_1w_2} & F(B)\\
{\scriptstyle w_1}\downarrow & \nearrow_{\scriptstyle w_2} & \downarrow{\scriptstyle F(w_1)}\\
\sigma B & \xrightarrow{w_2F(w_1)} & F(\sigma B)
\end{array}
\]

(the bottom one because \(l(w_2F(w_1))=l(w_2)+l(F(w_1))\)), and \(\sigma B\in X(w')\). By the same argument applied to \(w_2F(w_1)\) and \(F(w_1)F(w_2)=F(w)\), we get a map \(\tau:X(w')\to X(Fw)\). The diagram \[
\begin{array}{ccc}
X(w) & \xrightarrow{\sigma} & X(w')\\
{\scriptstyle F}\downarrow & \swarrow_{\scriptstyle\tau} & \downarrow{\scriptstyle F}\\
X(Fw) & \xrightarrow{\sigma^{(q)}} & X(Fw')
\end{array}
\] is commutative. The vertical maps induce equivalences of étale sites, hence so do \(\tau\) and \(\sigma\) [SGA 1, IX, 4.10]. The resulting isomorphism \[
H_c^*(X(w'))\xrightarrow{\sim}H_c^*(X(w))
\] is \(G^F\)-equivariant, whence (1.6) in this case.\par

Case 2. For some fundamental reflection \(s\), we have \(w'=swF(s)\) and \(l(w')=l(w)+2\). For \(B\in X(w')\), we can find two Borel subgroups \(\gamma B\), \(\delta B\) such that \((B,\gamma B)\in O(s)\), \((\gamma B,\delta B)\in O(w)\), \((\delta B,FB)\in O(F(s))\); moreover, \(\gamma B\) and \(\delta B\) are uniquely determined by these requirements. We define a partition \(X(w')=X_1\cup X_2\) by \[
X_1=\{B\in X(w')\mid \delta B=F(\gamma B)\},\qquad
X_2=\{B\in X(w')\mid \delta B\ne F(\gamma B)\}.
\]\par

Note that \(X_1\) is a closed subscheme of \(X(w')\), while \(X_2\) is an open one. We have \(\gamma:X_1\to X(w)\) and, for \(B'\in X(w)\), \(\gamma^{-1}(B')\) is the set of all Borel subgroups \(B\) such that \((B,B')\in O(s)\); hence \(\gamma^{-1}(B')\) is an affine line over \(k\), and \(X_1\) is (via \(\gamma\)) an affine line bundle over \(X(w)\). It follows that \(\gamma\) induces an isomorphism \[
\tag{1.6.1}
H_c^i(X_1)\xrightarrow{\sim}H_c^{i-2}(X(w))(-1),\qquad i\ge 0.
\]\par

If \(B\in X_2\), we have \[
(\delta B,FB)\in O(F(s)),\quad (FB,F(\gamma B))\in O(F(s)),\quad \delta B\ne F(\gamma B),
\] hence \((\delta B,F(\gamma B))\in O(F(s))\). Since \[
(F(\gamma B),F(\delta B))\in O(F(w)),\qquad l(F(sw))=l(F(s))+l(F(w)),
\] it follows that \((\delta B,F(\delta B))\in O(F(sw))\). Thus we have \(\delta:X_2\to X(F(sw))\). Let \[
X_2'=\{(B,\widetilde B)\in X_2\times X(sw)\mid F(\widetilde B)=\delta B\}
\] and let \(\delta':X_2'\to X(sw)\) be defined by \(\delta'(B,\widetilde B)=\widetilde B\); we define also \(\varphi:X_2'\to X_2\) by \(\varphi(B,\widetilde B)=B\). We have a cartesian\par
}

\EditionPage{8}{109}{%
diagram \[
\tag{1.6.2}
\begin{array}{ccc}
X_2' & \xrightarrow{\delta'} & X(sw)\\
{\scriptstyle\varphi}\downarrow & & \downarrow{\scriptstyle F}\\
X_2 & \xrightarrow{\delta} & X(F(sw)).
\end{array}
\]\par

For any \(\widetilde B\in X(sw)\) we denote by \(s\widetilde B\) the unique Borel subgroup such that \((\widetilde B,s\widetilde B)\in O(s)\), \((s\widetilde B,F\widetilde B)\in O(w)\). It is easy to see that \(\delta'^{-1}(\widetilde B)\) can be identified with the set of Borel subgroups \(B\) such that \((\widetilde B,B)\in O(s)\), \(B\ne s\widetilde B\); it follows that \(\delta'^{-1}(\widetilde B)\) is an affine line with a point removed. Thus \(X_2'\) is (via \(\delta'\)) a line bundle over \(X(sw)\) with the zero-section removed. Since the vertical arrows in (1.6.2) induce equivalences of étale sites, we have a canonical exact sequence (see [5]): \[
\tag{1.6.3}
\cdots\xrightarrow{\partial}H_c^{i-1}(X(sw))\longrightarrow H_c^i(X_2)\longrightarrow H_c^{i-2}(X(sw))(-1)
\xrightarrow{\partial}H_c^i(X(sw))\longrightarrow\cdots.
\]\par

It can be proved that the maps \(\partial\) in (1.6.3) are zero; this fact will not be used here.\par

Note that (1.6.1) and (1.6.3) are \(G^F\)-equivariant and that \(G^F\) acts trivially on \(\overline{\mathbb{Q}}_l(-1)\). It follows that for any \(g\in G^F\): \[
\operatorname{tr}(g^*,H_c^*(X_2))=0,
\]  \[
\operatorname{tr}(g^*,H_c^*(X_1))=\operatorname{tr}(g^*,H_c^*(X(w))).
\]\par

If one uses the exact sequence \[
\cdots\longrightarrow H_c^{i-1}(X_1)\longrightarrow H_c^i(X_2)\longrightarrow H_c^i(X(w'))\longrightarrow H_c^i(X_1)\longrightarrow\cdots,
\] it follows that \[
\begin{aligned}
\operatorname{tr}(g^*,H_c^*(X(w')))
&=\operatorname{tr}(g^*,H_c^*(X_1))+\operatorname{tr}(g^*,H_c^*(X_2))\\
&=\operatorname{tr}(g^*,H_c^*(X(w)))
\end{aligned}
\] and 1.6 is proved in this case.\par

The general case. It suffices to treat the case where \(w'=swF(s)\) for a fundamental reflection \(s\). By permuting, if necessary, \(w\) and \(w'\) (\(w=sw'F(s)\)), we may even assume that \(l(w')\ge l(w)\). If \(l(w')>l(w)\) we are in case 2. If \(l(w')=l(w)\), the following lemma shows that either we are in case 1 (with \(w\) and \(w'\) possibly interchanged) or that \(w=w'\).\par

\noindent\textbf{LEMMA 1.6.4.} Let \(s,t\) be two fundamental reflections in \(W\) and let \(w\in W\) be such that \(l(w)=l(swt)\). Then either \(w=swt\), or \[
l(sw)=l(w)-1,\qquad\text{or}\qquad l(wt)=l(w)-1.
\]\par
}

\EditionPage{9}{110}{%
Let \(w=s_1s_2\cdots s_k\) be a reduced expression for \(w\) (0.4). Assume that \(l(wt)=l(w)+1\); then \(wt=s_1s_2\cdots s_kt\) is also a reduced expression. We have \(l(swt)=l(wt)-1\). It follows (cf. Bourbaki, [1, Ch. IV, §1, Lemme 3]) that either there exists \(j\), \(1\le j\le k\), with \[
ss_1\cdots s_{j-1}=s_1\cdots s_{j-1}s_j,
\] or we have \[
ss_1\cdots s_k=s_1\cdots s_kt.
\] In the first case, we have \(w=ss_1\cdots s_{j-1}s_{j+1}\cdots s_k\) and \(l(sw)=l(w)-1\); in the second case we have \(w=swt\) and the lemma is proved.\par

1.7. Let us choose a maximal torus \(T^*\) in \(G\) and a Borel subgroup \(B^*\subset G\) containing \(T^*\), with unipotent radical \(U^*\). The quotient \(E=G/U^*\) is a \(T^*\)-torsor (= right principal homogeneous space of \(T^*\)) over \(X=G/B^*\). For \(x\in X\), the fibre \(E(x)\) of the projection \(E\to X\) is \[
E(x)=\{g\in G\mid ge^*=x\}/U^*,
\] where \(e^*\) is the point of \(X\) corresponding to \(B^*\).\par

Let \(\dot w\in N(T^*)\) define the element \(w\) in the Weyl group \(\mathbf W\) via the isomorphism \(\sigma(T^*,B^*):\mathbf W\xrightarrow{\sim}N(T^*)/T^*\). If \(x,y\in X\) are in relative position \(w\), the \(g\)'s in \(G\) such that \(ge^*=x\) and \(g\dot we^*=y\) form a torsor \(A(x,y)\) under \[
B^*\cap \dot wB^*\dot w^{-1}=T^*\cdot(U^*\cap \dot wU^*\dot w^{-1}).
\] For \(g\in A(x,y)\), the class of \(g\dot w\) in \(E(y)\) depends only on the class of \(g\) in \(E(x)\). This defines a map \(E(x)\to E(y)\), which we denote as right multiplication by \(\dot w\). We have the formulas \[
\tag{1.7.1}
(ut)\dot w=(u\dot w)\operatorname{ad}\dot w^{-1}(t),
\]  \[
\tag{1.7.2}
u(\dot w t)=(u\dot w)t.
\]\par

We will express (1.7.1) by saying that \(\cdot\dot w\) is a \(w\)-map of \(T^*\)-torsors. It is induced by a \(w\)-map of \(T^*\)-torsors over \(O(w)\)  \[
\cdot\dot w:\operatorname{pr}_1^*E\longrightarrow\operatorname{pr}_2^*E.
\] Assume that \(w=w_1w_2\), \(\dot w=\dot w_1\dot w_2\) and that \(l(w)=l(w_1)+l(w_2)\); then for \(x,y,z\in X\) with \((x,y)\in O(w_1)\) and \((y,z)\in O(w_2)\), we have \((x,z)\in O(w)\) and \[
\tag{1.7.3}
u\dot w=(u\dot w_1)\dot w_2.
\]\par

1.8. We now assume that \(T^*\) and \(B^*\) are \(F\)-stable. The identification \(\sigma(T^*,B^*)\) of \(T^*\) and \(N(T^*)/T^*\) with the torus \(\mathbf T\) and the Weyl group \(\mathbf W\) is then compatible with \(F\). For \(w\) in the Weyl group, we denote by \(\mathbf T(w)\) the torus \(\mathbf T\), provided with the rational structure for which the Frobenius is \(\operatorname{ad}(\dot w)F\). We have \[
\mathbf T(w)^F\simeq\{t\in T^*\mid \operatorname{ad}(\dot w)F(t)=t\}.
\] For \(x\in X\), the Frobenius map induces a map \(F:E(x)\to E(F(x))\), with \(F(ut)=F(u)F(t)\). For \(x\in X(w)\), we put\par
}

\EditionPage{10}{111}{%
\[
E(x,\dot w)=\{u\in E(x)\mid F(u)=u\dot w\}.
\]

This is a \(\mathbf T(w)^F\)-torsor. The \(E(x,\dot w)\) are the fibres of a map \[
\pi:\widetilde X(\dot w)\longrightarrow X(w),
\] with \(\widetilde X(\dot w)\subset E\mid X(w)\) a \(\mathbf T(w)^F\)-torsor over \(X(w)\). The action of \(G\) on \(E\) restricts to an action of \(G^F\) on \(\widetilde X(\dot w)\).\par

Up to isomorphism, the \(G^F\)-equivariant \(\mathbf T(w)^F\)-torsor \(\widetilde X(\dot w)\) over \(X(w)\) is independent of the lifting \(\dot w\) of \(w\) in \(N(T^*)\): for \(\dot w'=\dot wt\), there exists \(t_1\) with \(\operatorname{ad}\dot w^{-1}(t_1)F(t_1)^{-1}=t\) and the map \(u\mapsto ut_1\) induces an isomorphism \[
\widetilde X(\dot w)\xrightarrow{\sim}\widetilde X(\dot w').
\]\par

The groups \(G^F\) and \(\mathbf T(w)^F\) act on \(H_c^*(\widetilde X(\dot w),\overline{\mathbb{Q}}_l)\) by transport of structure. For any \(\theta\in\operatorname{Hom}(\mathbf T(w)^F,\overline{\mathbb{Q}}_l^*)\) we denote by \(H_c^*(\widetilde X(\dot w),\overline{\mathbb{Q}}_l)_\theta\) the subspace of \(H_c^*(\widetilde X(\dot w),\overline{\mathbb{Q}}_l)\) on which \(\mathbf T(w)^F\) acts by \(\theta\).\par

\noindent\textbf{DEFINITION 1.9.} \(R^\theta(w)\) is the virtual representation \[
\sum(-1)^iH_c^i(\widetilde X(\dot w),\overline{\mathbb{Q}}_l)_\theta
\] of \(G^F\) (an element of the Grothendieck group of representations of \(G^F\) over \(\overline{\mathbb{Q}}_l\)).\par

The character \(\theta\) can be used to transform the \(\mathbf T(w)^F\)-torsor \(\widetilde X(\dot w)\) into a local system of \(\overline{\mathbb{Q}}_l\)-vector spaces of rank one \(\mathcal F_\theta\) over \(X(w)\), provided with \[
\theta:\widetilde X(\dot w)\longrightarrow\mathcal F_\theta,\qquad \theta(xt)=\theta(x)\theta(t).
\] The morphism \(\pi:\widetilde X(\dot w)\to X(w)\) is finite and \[
\pi_*\overline{\mathbb{Q}}_l=\bigoplus_\theta\mathcal F_\theta.
\] The sheaf \(\mathcal F_\theta\) is the subsheaf of \(\pi_*\overline{\mathbb{Q}}_l\) on which \(\mathbf T^F\) acts by \(\theta\), hence \[
H_c^*(\widetilde X(\dot w),\overline{\mathbb{Q}}_l)_\theta=H_c^*(X(w),\mathcal F_\theta).
\] In particular, for \(\theta=1\), \[
R^1(w)=\sum(-1)^iH_c^i(X(w),\overline{\mathbb{Q}}_l)
\] so that Definition 1.9 is compatible with (1.5).\par

Example 1.10. For \(w=\dot w=e\), \(\pi:\widetilde X(\dot w)\to X(w)\) becomes the projection \[
\pi:G^F/U^{*F}\longrightarrow G^F/B^{*F},
\] and \(R^\theta(w)\) is the representation of \(G^F\) on the space of functions on \(G^F\) satisfying \[
f(gtu)=\theta(t)^{-1}f(g),
\]  \(G^F\) acting by \((g*f)(x)=f(g^{-1}x)\) (induced representation).\par

1.11. The Borel subgroup \(\operatorname{ad}gB^*\) is in \(X(w)\) if and only if \(\operatorname{ad}gB^*\) and \(\operatorname{ad}FgB^*\) are in relative position \(w\), i.e., if and only if \(g^{-1}Fg\in B^*\dot wB^*\) (where \(\dot w\in N(T^*)\) represents \(w\)): \[
\tag{1.11.1}
X(w)=\{g\in G\mid g^{-1}Fg\in B^*\dot wB^*\}/B^*.
\] If a Borel subgroup \(B\) is in \(X(w)\), one can find \(g\in G\) such that \(\operatorname{ad}gB^*\) is\par
}

\EditionPage{11}{112}{%
\(B\) and \(\operatorname{ad}g\operatorname{ad}\dot wB^*=FB\): \[
\tag{1.11.2}
X(w)=\{g\in G\mid g^{-1}Fg\in\dot wB^*\}/\bigl(B^*\cap\operatorname{ad}\dot wB^*\bigr),
\] where \[
B^*\cap\operatorname{ad}\dot wB^*=T^*\cdot(U^*\cap\operatorname{ad}\dot wU^*).
\] Changing \(g\) to \(gt\), we can normalize \(g\) so that we have also \(g^{-1}Fg\in\dot wU^*\): \[
\tag{1.11.3}
X(w)=\{g\in G\mid g^{-1}Fg\in\dot wU^*\}/\bigl(\mathbf T(w)^F\cdot(U^*\cap\operatorname{ad}\dot wU^*)\bigr).
\] A point in \(\widetilde X(\dot w)\) is defined by a Borel subgroup \(B\), plus \(g\in G\) such that \(\operatorname{ad}gB^*=B\), \(\operatorname{ad}g\operatorname{ad}\dot wB^*=FB\) and \(g\dot w=Fg\bmod U^*\): \[
\tag{1.11.4}
\widetilde X(\dot w)=\{g\in G\mid g^{-1}Fg\in\dot wU^*\}/\bigl(U^*\cap\operatorname{ad}\dot wU^*\bigr).
\]\par

\noindent\textbf{COROLLARY 1.12.} The following assertions are equivalent:\par

(i) \(X(w)\) is affine;\par

(ii) \(\widetilde X(\dot w)\) is affine;\par

(iii) Let \(\rho\) be the action of \(U^*\cap\operatorname{ad}\dot wU^*\) on \(U^*\) defined by \[
\rho(u)v=\operatorname{ad}\dot w^{-1}(u)vF(u^{-1});
\] then \(U^*/\rho(U^*\cap\operatorname{ad}\dot wU^*)\) is affine.\par

Put \(S=\{g\in G\mid g^{-1}Fg\in\dot wU^*\}\). The map \(f:S\to U^*:g\mapsto\dot w^{-1}g^{-1}Fg\) induces an isomorphism \(G^F\backslash S\to U^*\) and is such that, for \(u\in U^*\cap\operatorname{ad}\dot wU^*\), \[
f(gu)=\rho(u)^{-1}f(g).
\] Hence, \[
G^F\backslash\widetilde X(\dot w)\simeq U^*/(U^*\cap\operatorname{ad}\dot wU^*).
\] As \(\widetilde X(\dot w)/\mathbf T(w)^F=X(w)\), it only remains to use the fact that a space and a quotient of it by a finite group are simultaneously affine or not.\par

We will have to use another description of \(\pi:\widetilde X(\dot w)\to X(w)\). First, an easy lemma:\par

\noindent\textbf{LEMMA 1.13.} Let \(J\) be the set of pairs \((T,B)\), \(T\) an \(F\)-stable maximal torus and \(B\) a Borel subgroup containing \(T\). The map \(h\) which to \((T,B)\) associates the relative position of \(B\) and \(FB\) induces a bijection \[
G^F\backslash J\xrightarrow{\sim}\mathbf W.
\] The proof will be given in (1.15).\par

If we use \((T,B)\in J\) to identify \(\mathbf W\) with \(N(T)/T\), we have \[
\tag{1.13.1}
h(T,\operatorname{ad}\dot wB)=w^{-1}h(T,B)F(w),
\] where \(\dot w\in N(T)\) represents \(w\), hence\par

\noindent\textbf{COROLLARY 1.14.} The map \(h\) induces a bijection \[
\{G^F\text{-conjugacy classes of }F\text{-stable maximal tori}\}\xrightarrow{\sim}\mathbf W_F^{\sim}.
\]\par

Here is another description of \(h\): for \((T,B)\in J\), if \(\sigma:\mathbf T\to T\) and \(\sigma:\mathbf W\to W\)\par
}

\EditionPage{12}{113}{%
(\(W=N(T)/T\)) are defined by \((T,B)\), then \[
\sigma F\sigma^{-1}=\operatorname{ad}h(T,B)\circ F:\mathbf T\longrightarrow\mathbf T.
\] To give \(h(T,B)\) is the same as to give \(\operatorname{ad}h(T,B)\circ F\in\mathbf W\circ F\subset\operatorname{End}(\mathbf T)\), and to give the \(F\)-conjugacy class of \(h(T,B)\) is the same as to give \(\operatorname{ad}h(T,B)\circ F\) up to \(\mathbf W\)-conjugacy.\par

1.15. The space of maximal tori of \(G\) can be identified with the homogeneous space \(G/N(T^*)\), and the space of maximal tori marked by a containing Borel subgroup can be identified with \(G/T^*\). The group \(T^*\) being the connected component of \(N(T^*)\), (1.13) is a special case of the general result described below.\par

Let \(G_0\) be a connected algebraic group over \(\mathbb F_q\) and let \(x:\widetilde E_0\to E_0\) be a morphism of \(G_0\)-homogeneous spaces. We denote by \(G,\widetilde E,E\) the corresponding objects over \(k\), and we assume that the stabilizer \(S(\widetilde e)\) of \(\widetilde e\in\widetilde E\) is the connected component of the stabilizer \(S(\pi(\widetilde e))\) of \(\pi(\widetilde e)\in E\). Since any \(G_0\)-homogeneous space has a rational point, the existence of \(\widetilde E_0\) imposes no condition on \(E_0\).\par

The groups \(S(e)/S(e)^0\) form a local system on \(E\), which becomes constant on \(\widetilde E\); we denote by \(W\) its constant value on \(\widetilde E\). For \(\widetilde e\in\widetilde E\), we have an isomorphism \[
\alpha(\widetilde e):W\xrightarrow{\sim}S(\pi(\widetilde e))/S(\pi(e)),
\] and, for \(\widetilde f=g\widetilde e\), we have \(\alpha(\widetilde f)=\operatorname{ad}g\,\alpha(\widetilde e)\). We let the group \(W\) act on \(\widetilde E\) on the right, by \[
\widetilde e\,w=\alpha(\widetilde e)(w)\widetilde e;
\] in this way, \(\widetilde E\) becomes a \(W\)-torsor (= principal homogeneous space) over \(E\).\par

The group \(W\) is acted on by \(F\), with \(F(\widetilde e\,w)=F(\widetilde e)F(w)\). The set \(W_F^{\sim}\) of \(F\)-conjugacy classes in \(W\) is the set of orbits of the action of \(W\) on itself by \[
w\longmapsto w_1wF(w_1)^{-1}.
\]\par

\noindent\textbf{PROPOSITION 1.16.} For \(e\in E^F\) and \(\widetilde e\in E\) above it, define \(h(e,\widetilde e)\in W\) by \[
F(\widetilde e)=\widetilde e\cdot h(e,\widetilde e).
\]\par

(i) The map \(h\) induces a bijection from the set of \(G^F\)-orbits in \[
\{(e,\widetilde e)\mid e\in E^F,\ \pi(\widetilde e)=e\}
\] to \(W\).\par

(ii) We have \[
h(e,\widetilde e w)=w^{-1}h(e,\widetilde e)F(w).
\] Hence the map \(h\) induces a bijection from \(G^F\backslash E^F\) to the set of \(F\)-conjugacy classes in \(W\).\par

We will only prove (i). Let \(Y\) be the set of \(\widetilde e\in E\) such that \(\pi(\widetilde e)\in E^F\). If \(\widetilde e_0\in E^F\), the map \(g\mapsto g\widetilde e_0\) identifies \(X\) with \[
\{g\in G\mid g^{-1}Fg\in S(\pi(\widetilde e_0))\}/S(\widetilde e_0).
\] The Lang isogeny \(g\mapsto g^{-1}Fg\) is an isomorphism \(G^F\backslash G\to G\), hence the map \(ge_0\mapsto g^{-1}Fg\) induces a bijection \[
G^F\backslash X\xrightarrow{\sim}S(\pi(\widetilde e_0))/\bigl(S(\widetilde e_0)\text{ acting by }x\mapsto s^{-1}xF(s)\bigr).
\] The orbits of this action of \(S(\widetilde e_0)\) are just the usual cosets, and the resulting\par
}

\EditionPage{13}{114}{%
bijection \(G^F\backslash X\xrightarrow{\sim}W\) is the \(h\) above.\par

\noindent\textbf{DEFINITION 1.17.} Let \(T\) be an \(F\)-stable maximal torus and let \(B\) be a Borel subgroup containing \(T\), with unipotent radical \(U\); let \(w\) be the relative position of \(B\) and \(FB\).\par

(i) \(X_{T\subset B}\) is \(X(w)\). The map \(g\mapsto\operatorname{ad}gB\) induces isomorphisms \[
\begin{aligned}
X_{T\subset B}
&=\{g\in G\mid g^{-1}Fg\in B\cdot F(B)\}/B\\
&=\{g\in G\mid g^{-1}Fg\in FB\}/(B\cap FB)\\
&=\{g\in G\mid g^{-1}Fg\in FU\}/\bigl(T^F\cdot(U\cap FU)\bigr).
\end{aligned}
\]\par

(ii) \(\widetilde X_{T\subset B}\) is \[
\{g\in G\mid g^{-1}Fg\in FU\}/(U\cap FU).
\] We have a projection map \(\pi:\widetilde X_{T\subset B}\to X_{T\subset B}\), for which \(\widetilde X_{T\subset B}\) is a \(G^F\)-equivariant \(T^F\)-torsor over \(X_{T\subset B}\); \(G^F\) acts by left multiplication and \(T^F\) by right multiplication (it normalizes \(U\cap FU\)).\par

1.18. Let \(\dot w\in N(T^*)\) be a representative of \(w\). If \(x'\in G\) is such that \(\operatorname{ad}x'(T^*,B^*)=(T,B)\) then \(B\) and \(FB=\operatorname{ad}Fx'B^*\) are in relative position \(w\), and \(FB\) contains \(T\), hence \(FB=\operatorname{ad}x'\operatorname{ad}\dot wB^*\) and \[
x'^{-1}F(x')\in\dot wB^*\cap N(T^*)=\dot wT^*.
\] Replacing \(x'\) by \(x=x't\) (\(t\in T^*\)) one can achieve \(x^{-1}F(x)=\dot w\). The \(x\) such that \(\operatorname{ad}x(T^*,B^*)=(T,B)\) and \(x^{-1}F(x)=\dot w\) form a \(T^{*F}\)-torsor. For such an \(x\), \(\operatorname{ad}x\) induces an isomorphism \(\mathbf T(w)\to T\) (hence \(\mathbf T(w)^F\to T^F\)); this isomorphism is independent of \(x\).\par

\noindent\textbf{PROPOSITION 1.19.} Let \(T,B,U,w\) be as in (1.17) and \(x,\dot w\) as in (1.18). The map \(g\mapsto gx^{-1}\) induces an isomorphism from the \(G^F\)-equivariant \(\mathbf T(w)^F\)-torsor \(\widetilde X(\dot w)\) over \(X(w)\) (or rather its model (1.11)) to the \(G^F\)-equivariant \(T^F\)-torsor \(\widetilde X_{T\subset B}\) over \(X_{T\subset B}\) (or rather its model (1.17)).\par

This is a straightforward computation.\par

1.20. The cohomology of \(\widetilde X_{T\subset B}\) is acted on by \(G^F\) and \(T^F\). For any character \(\theta:T^F\to\overline{\mathbb Q}_l^*\) we put as in (1.8): \[
R_{T\subset B}^{\theta}=\sum_i(-1)^iH_c^i(\widetilde X_{T\subset B},\overline{\mathbb Q}_l)_\theta
\] (an element in the Grothendieck group of representations of \(G^F\) over \(\overline{\mathbb Q}_l\)).\par

By (1.19), for \(x\) as in (1.18), we have \[
R_{T\subset B}^{\theta}=R^{\theta\circ\operatorname{ad}x}(w).
\] We shall see in Chapter 4 that \(R_{T\subset B}^{\theta}\) is independent of \(B\). For \(g\in G^F\) such that \(\operatorname{ad}g\) carries \(T,B,\theta\) to \(T',B',\theta'\), we have clearly \[
R_{T\subset B}^{\theta}=R_{T'\subset B'}^{\theta'},
\] hence \(R_{T\subset B}^{\theta}\) will eventually depend only on the \(G^F\)-conjugacy class of \(T\) and on the orbit of \(\theta\) under \((N(T)/T)^F\).\par

The end of this chapter will be used in the proof of 7.10 only.\par
}

\EditionPage{14}{115}{%
1.21. Isogenies. Let \(T\) be an \(F\)-stable maximal torus of \(G\), and let \(Z\) be the centre of \(G\). We denote \(G\times^Z T\) the quotient of \(G\times T\) by the subgroup \[
\{(z,z^{-1})\mid z\in Z\}.
\]\par

Let \(B\) be a Borel subgroup containing \(T\). The action \(x\mapsto gxt\) of \(G^F\times T^F\) on \(\widetilde X_{T\subset B}\) is induced by an action of \((G\times^Z T)^F\), given by the same formula.\par

\noindent\textbf{COROLLARY 1.22.} On \(H^*(\widetilde X_{T\subset B},\overline{\mathbb Q}_l)_\theta\), \(Z^F\) acts by the character \(\theta\mid Z^F\). In particular, on any irreducible representation occurring in \(R_{T\subset B}^{\theta}\), \(Z^F\) acts by \(\theta\mid Z^F\).\par

Let \(\pi:\widetilde G\to G\) be the simply connected covering of the derived group of \(G\), \(\widetilde T=\pi^{-1}(T)\), \(\widetilde B=\pi^{-1}(B)\) and let \(\widetilde Z\) be the centre of \(\widetilde G\).\par

\noindent\textbf{PROPOSITION 1.23.} One has \[
T^F/\pi(\widetilde T^F)\xrightarrow{\sim}G^F/\pi(\widetilde G^F).
\]\par

Injectivity is clear; we have to check the surjectivity of the map \[
(T\times\widetilde G)^F\longrightarrow G^F
\] induced by \[
\varphi:T\times\widetilde G\longrightarrow G:\quad t,\widetilde g\longmapsto t\pi(\widetilde g).
\] Via \(\varphi\), \(T\times\widetilde G\) is a \(\widetilde T\)-torsor over \(G\) (with \[
(t,\widetilde g)*\widetilde t=(t\widetilde t,\widetilde t^{-1}\widetilde g)),
\] and one applies Lang's theorem to the connected group \(\widetilde T\).\par

1.24. Let \(h:A\to B\) be a homomorphism of finite groups and let \(X\) be a space on which \(A\) acts. The induced space \(\operatorname{Ind}_A^B(X)\) (unique up to unique isomorphism) is any \(B\)-space \(I\), provided with an \(A\)-equivariant map \(\varphi:X\to I\), such that for any \(B\)-space \(Y\), \[
\operatorname{Hom}_B(I,Y)\xrightarrow[\varphi]{\sim}\operatorname{Hom}_A(X,Y).
\] One has \[
\operatorname{Ind}_A^B(X)=\coprod_{b\in B/A}b\varphi(X),
\qquad
\varphi(X)\simeq\ker(h)\backslash X.
\]\par

\noindent\textbf{PROPOSITION 1.25.} The \((G\times^Z T)^F\)-space \(\widetilde X_{T\subset B}\) is induced by the \((\widetilde G\times^{\widetilde Z}\widetilde T)^F\)-space \(\widetilde X_{\widetilde T\subset\widetilde B}\).\par

\noindent\textbf{LEMMA 1.26.} In the diagram \[
\begin{array}{ccccccccccc}
0&\longrightarrow&\ker&\longrightarrow&\widetilde T^F&\longrightarrow&T^F&\longrightarrow&\operatorname{coker}&\longrightarrow&0\\
&&{\scriptstyle(1)}\downarrow&&{\scriptstyle(t,1)}\downarrow&&{\scriptstyle(t,1)}\downarrow&&{\scriptstyle(2)}\downarrow&&\\
0&\longrightarrow&\ker&\longrightarrow&(\widetilde T\times^{\widetilde Z}\widetilde G)^F&\longrightarrow&(T\times^Z G)^F&\longrightarrow&\operatorname{coker}&\longrightarrow&0\\
&&\downarrow&&\downarrow&&&&&&\\
&&(\widetilde G/\widetilde Z)^F&=&(G/Z)^F&&&&&&
\end{array}
\] the maps (1) and (2) are isomorphisms.\par

Indeed, \(\widetilde T\times^{\widetilde Z}\widetilde G\) (resp. \(T\times^Z G\)) is a \(\widetilde T\) (resp. \(T\))-torsor over \(\widetilde G/\widetilde Z=G/Z\), hence, since \(T\) and \(\widetilde T\) are connected, \((\widetilde T\times^{\widetilde Z}\widetilde G)^F\) is a \(\widetilde T^F\)-torsor over \((\widetilde G/\widetilde Z)^F\), and \((T\times^Z G)^F\) is the induced \(T^F\)-torsor.\par

Proof of 1.25. By 1.26, we are reduced to prove that \(\widetilde X_{T\subset B}\), as a \(T^F\)-space, is induced by the \(\widetilde T^F\)-space \(\widetilde X_{\widetilde T\subset\widetilde B}\). The spaces \(\widetilde X_{T\subset B}\) and \(\widetilde X_{\widetilde T\subset\widetilde B}\) are\par
}

\EditionPage{15}{116}{%
indeed respectively \(T^F\) and \(\widetilde T^F\)-torsor over \(X_{T\subset B}=X_{\widetilde T\subset\widetilde B}\), and \(X_{\widetilde T\subset\widetilde B}\) is hence the \(T^F\)-torsor induced by the \(\widetilde T^F\)-torsor \(X_{T\subset B}\).\par

\noindent\textbf{COROLLARY 1.27.} Let \(\theta\) be a character of \(G^F/\pi(\widetilde G^F)\). We denote again by \(\theta\) its restriction to \(T^F\). One has \[
R_{T\subset B}^{\theta\eta}=\theta\otimes R_{T\subset B}^{\eta}.
\]\par

It follows from 1.25 that \[
H^*(\widetilde X_{T\subset B},\overline{\mathbb Q}_l)
=
\operatorname{Ind}_{(\widetilde G\times^{\widetilde Z}\widetilde T)^F}^{(G\times^Z T)^F}
(\widetilde X_{\widetilde T\subset\widetilde B},\overline{\mathbb Q}_l).
\] The character \(\theta(gt)\) of \((G\times^{\widetilde Z} T)^F\) is trivial on the image of \(\widetilde G\times^{\widetilde Z}\widetilde T\). The induced representation we consider is hence isomorphic to its tensor product with \(\theta(gt)\), and 1.27 is a formal consequence of this.\par

\EditionChapter{2. Examples in the classical groups}

2.1. Let \(V\) be an \(n\)-dimenssional vector space over \(k\) and put \(G=\operatorname{GL}(V)\). If \(b=(b_1,\ldots,b_n)\) is a basis of \(V\), we may take for \(T^*\) the group of diagonal matrices and for \(B^*\) the group of upper triangular matrices. The Weyl group lifts into the subgroup of \(N(T^*)\) consisting of the \(\dot w\)'s inducing a permutation of basis vectors. In this case, \(\mathbf T,\mathbf W\) (1.1), \(X\) (1.2), \(E\), \(\cdot\dot w\) (1.7), have the following alternative description.\par

(a) \(\mathbf T=\mathbf G_m^n\), \(\mathbf W=\mathfrak S_n\), the fundamental reflections are the transpositions \((i,i+1)\) and the action of \(\mathbf W\) on \(\mathbf T\) is by permutation.\par

(b) \(X\) is the space of complete flags \(D_1\subset\cdots\subset D_{n-1}\) in \(V\): a flag \(D\) is an increasing filtration of \(V\) with \(\dim D_i=i\) for \(1\le i\le n-1\).\par

(c) \(E\) is the space of complete flags marked by non-zero vectors \[
e_i\in D_i/D_{i-1}=\operatorname{Gr}_i^D(V)\qquad(1\le i\le n),
\] where we use the convention \(D_0=0\), \(D_n=V\); \(\mathbf T\) acts on \(E\) by \[
(D,(e_i))(\lambda_i)=(D,(\lambda_i e_i)).
\] This is a \(G\)-equivariant \(\mathbf T\)-torsor over \(X\).\par

(d) If \(D'\) and \(D''\) are two flags, their relative position is labelled by the permutation \(w\) such that \[
\operatorname{Gr}_{w(i)}^{D'}\operatorname{Gr}_i^{D''}(V)\ne0.
\] The isomorphisms \[
\operatorname{Gr}_{w(i)}^{D'}(V)
\simeq\operatorname{Gr}_i^{D''}\operatorname{Gr}_{w(i)}^{D'}(V)
\simeq\operatorname{Gr}_{w(i)}^{D'}\operatorname{Gr}_i^{D''}(V)
\simeq\operatorname{Gr}_i^{D''}(V)
\] induce a \(w\)-isomorphism between the \(\mathbf T\)-torsor \(E(D')\) of markings of \(D'\) and \(E(D'')\): \[
e\longmapsto e\cdot\dot w.
\] When \(w\) is the \(n\)-cycle \((1,\ldots,n)\), \(D'\) and \(D''\) are in relative position \(w\) if and only if \[
D_i''+D_i'=D_{i+1}'\quad(1\le i<n-1),
\qquad
D_{n-1}''+D_1'=V.
\]\par

2.2. We now take \(k\) and \(\mathbb F_q\) as in (0.1.1) and assume that \(V\) is provided with an \(\mathbb F_q\)-structure. Frobenius maps are then defined. For \(w=(1,\ldots,n)\), the condition for a flag \(D\) to be in relative position \(w\) with its image \(FD\) by Frobenius is that \(D\) be the flag\par
}

\EditionPage{16}{117}{%
\[
D_1\subset D_1+FD_1\subset D_1+FD_1+F^2D_1\subset\cdots
\qquad\text{and that}\qquad
V=\bigoplus_0^{n-1}F^iD_1.
\] If we denote by \(\mathbf P(V)\) the set of homogeneous lines in \(V\), the map \(D\mapsto D_1\) is an isomorphism from \(X(w)\) to the set of all \(x\in\mathbf P(V)\) which do not lie on any \(\mathbb F_q\)-rational hyperplane. A marking \(e\) of \(F\) is such that \(F(e)=e\cdot\dot w\) if and only if \[
\begin{aligned}
e_2&\equiv F(e_1)\pmod{e_1},\\
e_3&\equiv F^2(e_1)\pmod{e_1,F(e_1)},\quad\ldots,\\
e_n&\equiv F^{n-1}(e_1)\pmod{e_1,F(e_1),\ldots,F^{n-2}(e_1)}
\end{aligned}
\] and \[
e_1\equiv F^n(e_1)\pmod{F(e_1),\ldots,F^{n-1}(e_1)};
\]  \(e\) is defined by \(e_1\in D_1\) subject to the condition that \[
e_1\wedge F(e_1)\wedge\cdots\wedge F^{n-1}(e_1)
=
F^n(e_1)\wedge F(e_1)\wedge\cdots\wedge F^{n-1}(e_1),
\] i.e.: \[
\tag{2.2.1}
F(e_1\wedge\cdots\wedge F^{n-1}(e_1))
=(-1)^{n-1}(e_1\wedge\cdots\wedge F^{n-1}(e_1)).
\] If \((x_i)\) are the coordinates of \(e_1\) with respect to some rational basis, the condition (2.2.1) can be rewritten \[
\tag{2.2.2}
(-1)^{n-1}\left(\det(x_i^{q^{j-1}})_{1\le i,j\le n}\right)^{q-1}=1.
\] The form on the left is invariant under \(\operatorname{GL}(n,\mathbb F_q)\). Up to a scalar factor, it is the product of all non-zero \(\mathbb F_q\)-rational linear forms. The map \((D,e)\mapsto e_1\) induces an isomorphism of \(\widetilde X(\dot w)\) with the affine hypersurface (2.2.2). This hypersurface is stable under \(x\mapsto\lambda x\) for \(\lambda\in\mathbb F_{q^n}^*\), and this is the action of \(\mathbf T(w)^F\).\par

Our work has been inspired by results of Drinfeld, who proved that the discrete series representations of \(\operatorname{SL}(2,\mathbb F_q)\) occur in the cohomology of the affine curve \[
xy^q-x^qy=1
\] (the form \[
xy^q-x^qy=\begin{vmatrix}x&x^q\\y&y^q\end{vmatrix}
\] is \(\operatorname{SL}_2(\mathbb F_q)\)-invariant).\par

Let \(\varphi(q,r,n)\) be the number of \(\mathbb F_{q^r}\)-rational points of \(X(w)\), \(r\ge1\). We have the following\par

\noindent\textbf{PROPOSITION 2.3.} \[
\varphi(q,r,n)=\prod_{1\le i\le n-1}(q^r-q^i).
\]\par

Proof. We define a partition \[
\mathbf P(V)=X_0\cup X_1\cup\cdots\cup X_{n-1}
\] as follows: \(X_i\) is the set of points \(x\in\mathbf P(V)\) such that \(x,F(x),F^2(x),\ldots\) span a linear subspace of dimension \(i\) of \(\mathbf P(V)\). This partition is invariant under Frobenius, and clearly \(X_i\) has precisely \[
\varphi(q,r,i+1)\frac{(1-q^{n-i})\cdots(1-q^n)}{(1-q)\cdots(1-q^{i+1})}
\]\par
}

\EditionPage{17}{118}{%
\(\mathbb F_q\)-rational points. It follows that \[
\frac{q^{nr}-1}{q-1}
=
\sum_{0\le i\le n-1}
\varphi(q,r,i+1)
\frac{(1-q^{n-i})\cdots(1-q^n)}{(1-q)\cdots(1-q^{i+1})}.
\] This shows that, for fixed \(n\), \(\varphi(q,r,n)\) is a polynomial of degree \((n-1)\) in \(Q=q^r\), with coefficient polynomials in \(q\), and leading term \(Q^{n-1}\). Since \(\varphi(q,r,n)=0\) for \(1\le r\le n-1\), this polynomial must be divisible by \((Q-q)\), \((Q-q^2)\), \(\ldots\), \((Q-q^{n-1})\). It follows that \[
\varphi(q,r,n)=(Q-q)(Q-q^2)\cdots(Q-q^{n-1})
\] and the proposition is proved.\par

2.4. Let \(V\) be a vector space as in (2.2), with a fixed non-degenerate symplectic form \(\langle\ ,\ \rangle\) defined over \(\mathbb F_q\). We must have \(n=2m\). The symplectic group \(\operatorname{Sp}(V)\) is then a group as in (0.1.1). Let \(Y\) be the set of all complete isotropic flags \(D_1\subset D_2\subset\cdots\subset D_m\) in \(V\) (\(\dim D_i=i\)) such that \[
D_1\ne FD_1\subset D_2,\quad
D_2\ne FD_2\subset D_3,\quad\ldots,\quad
D_{m-1}\ne FD_{m-1}\subset D_m,\quad
D_m\ne FD_m.
\] Then \(Y\) can be identified with \(X(w)\), where \(w\) is a Coxeter element in the Weyl group \(\mathbf W\) of \(\operatorname{Sp}(V)\); if we identify \(\mathbf W\) with the group of all permutations \(\sigma\) of \(-m,\ldots,-2,-1,1,2,\ldots,m\) such that \(\sigma(-i)=-\sigma(i)\) for all \(i\), then \(w\) is the permutation \[
\begin{aligned}
-i&\longmapsto -i+1 &&(2\le i\le m),&\qquad -1&\longmapsto m,\\
i&\longmapsto i-1 &&(2\le i\le m),& 1&\longmapsto -m.
\end{aligned}
\] On the other hand, \(Y\) (hence also \(X(w)\)) can be identified with the set of \(x\in\mathbf P(V)\) such that \[
\langle x,F(x)\rangle=\langle x,F^2(x)\rangle=\cdots=\langle x,F^{m-1}(x)\rangle=0,
\qquad
\langle x,F^m(x)\rangle\ne0.
\] For example, if \(\dim V=4\), this is just the set of all \(x\in\mathbf P(V)\) such that \(\langle x,F(x)\rangle=0\), \(\langle x,F^2(x)\rangle\ne0\). Note that the equation \(\langle x,F(x)\rangle=0\) defines a non-singular surface \(S\) in \(\mathbf P(V)\) and that \(\langle x,F^2(x)\rangle\ne0\) means that we remove from \(S\) a union of rational curves, one for each isotropic plane in \(V\), defined over \(\mathbb F_q\). One can prove that in this case (\(n=4\)), the number of \(\mathbb F_{q^r}\)-rational points of \(X(w)\) is given by: \[
q^{2r}-\frac q2(1+q)^2q^r+\frac q2(1-q)^2(-q)^r+q^4
=
\begin{cases}
(q^r-q^2)^2,&r\text{ even},\\
(q^r-q)(q^r-q^3),&r\text{ odd}.
\end{cases}
\]\par

\EditionChapter{3. A fixed point formula}

3.1. Let \(X\) be a scheme, separated and of finite type over \(k\) and let \(\sigma:X\to X\) be an automorphism of finite order of \(X\). We decompose \(\sigma\) as \(\sigma=s\cdot u\) where \(s\) and \(u\) are powers of \(\sigma\) respectively of order prime to \(p\) and\par
}

\EditionPage{18}{119}{%
a power of \(p\). The main result of this chapter is the following:\par

\noindent\textbf{THEOREM 3.2.} With the above notations (see also 0.5), \[
\operatorname{Tr}(\sigma^*,H_c^*(X,\overline{\mathbb Q}_l))
=
\operatorname{Tr}(u^*,H_c^*(X^s,\overline{\mathbb Q}_l)).
\]\par

The first step is to prove that the left hand side is an integer independent of \(l\). One might conjecture that for any endomorphism \(f\) of \(X\) and each \(i\), \(\operatorname{Tr}(f^*,H_c^i(X,\overline{\mathbb Q}_l))\) is integral and independent of \(l\), but such a more general and precise result is known only for \(X\) proper and smooth (Katz and Messing, Inv. Math. 23 (1974), 73-77).\par

\noindent\textbf{PROPOSITION 3.3.} With the notations of 3.1, \(\operatorname{Tr}(\sigma^*,H_c^*(X,\overline{\mathbb Q}_l))\) is an integer independent of \(l\) (\(l\ne p\)).\par

A standard specialization argument allows us to assume that \(p>1\) and that \(k\) is an algebraic closure of the prime field \(\mathbb F_p\). This is anyway the only case we will need in the rest of the paper. The scheme \(X\) and \(\sigma\) can then be defined over some finite extension \(\mathbb F_q\subset k\) of \(\mathbb F_p\); we denote by \(F:X\to X\) the corresponding Frobenius endomorphism.\par

Let us first assume that \(X\) is quasi-projective. Then, for \(n\ge1\), the composite \(F^n\circ\sigma\) is the Frobenius map relative to some new way of lowering the field of definition of \(X\) from \(k\) to \(\mathbb F_{q^n}\) and the Lefschetz fixed point formula for Frobenius ([5] and [11]) shows that \(\operatorname{Tr}((F^n\sigma)^*,H_c^*(X,\overline{\mathbb Q}_l))\) is the number of fixed points of \(F^n\sigma\) (\(n\ge1\)). The automorphisms \(F^*\) and \(\sigma^*\) of the cohomology commute; as a function of \(n\), \(\operatorname{Tr}((F^n\sigma)^*,H_c^*(X,\overline{\mathbb Q}_l))\) is hence of the form \(\sum_\lambda\alpha_\lambda\lambda^n\) where \(\lambda\) runs through the multiplicative group \(\overline{\mathbb Q}_l^*\) of \(\overline{\mathbb Q}_l\) and where \(\alpha_\lambda\) is zero for all \(\lambda\) except for a finite number of them.\par

The functions \(\mathbb N^+\to\overline{\mathbb Q}_l^*:n\mapsto\lambda^n\) (for \(\lambda\in\overline{\mathbb Q}_l^*\)) are linearly independent. This can be checked either by using Vandermonde determinants or by appealing to Dedekind's theorem on the linear independence of characters (Bourbaki, Algèbre V, § 7, 5). For \(n\ge1\) the numbers \(\sum_\lambda\alpha_\lambda\lambda^n=|X^{F^n\sigma}|\) are rational; hence for any automorphism \(\tau\) of \(\overline{\mathbb Q}_l\) one has \[
\sum_\lambda\tau(\alpha_\lambda)\tau(\lambda)^n
=
\sum_\lambda\tau(\alpha_{\tau^{-1}(\lambda)})\lambda^n
=
\sum_\lambda\alpha_\lambda\lambda^n
\qquad(n\ge1),
\] and, by the linear independence of the functions \(\lambda^n\), \[
\alpha_{\tau(\lambda)}=\tau(\alpha_\lambda).
\] In particular, \(\operatorname{Tr}(\sigma^*,H_c^*(X,\overline{\mathbb Q}_l))=\sum_\lambda\alpha_\lambda\) is invariant by any automorphism of \(\overline{\mathbb Q}_l\), hence rational. Similarly, as \(\operatorname{Tr}((F^n\sigma)^*,H_c^*(X,\overline{\mathbb Q}_l))=|X^{F^n\sigma}|\) is independent of \(l\), for any isomorphism \(\tau:\overline{\mathbb Q}_l\to\overline{\mathbb Q}_{l'}\) one has \[
\tau\operatorname{Tr}(\sigma^*,H_c^*(X,\overline{\mathbb Q}_l))
=
\operatorname{Tr}(\sigma^*,H_c^*(X,\overline{\mathbb Q}_{l'})),
\]\par
}

\EditionPage{19}{120}{%
hence the rational number \(\operatorname{Tr}(\sigma^*,H_c^*(X,\overline{\mathbb Q}_l))\) is independent of \(l\).\par

Since the automorphism \(\sigma\) has finite order, \(\operatorname{Tr}(\sigma^*,H_c^*(X,\overline{\mathbb Q}_l))\) is a sum of roots of unity, hence an algebraic integer. Being rational it is an ordinary integer.\par

If we do not assume \(X\) to be quasi-projective, we can either repeat the previous argument by working with algebraic spaces or reduce to the quasi-projective case: if \((X_i)\) is a finite partition of \(X\) into locally closed quasi-projective subschemes stable under \(\sigma\), then (cf. [5]) \[
\tag{3.3.1}
\operatorname{Tr}(\sigma^*,H_c^*(X,\overline{\mathbb Q}_l))
=
\sum_i\operatorname{Tr}(\sigma^*,H_c^*(X_i,\overline{\mathbb Q}_l)).
\]\par

3.4. Proof of 3.2. Let \((X_i)\) be a partition of \(X\) into locally closed subschemes stable under \(\sigma\) and such that, on each \(X_i\), \(\sigma\) defines a free action of a cyclic group. By applying (3.3.1) to the decompositions \(X=\bigcup X_i\) and \(X^s=\bigcup X_i^s\), one reduces to the case where \(\sigma\) generates a free action of a cyclic group \(H\). In this case, either\par

(a) \(\sigma\) is of order a power of \(p\), hence \(s=\operatorname{Id}\), \(u=\sigma\) and (3.2) is trivial, or\par

(b) the order of \(\sigma\) is divisible by a prime number \(l'\ne p\) and we must prove that \(\operatorname{Tr}(\sigma^*,H_c^*(X,\overline{\mathbb Q}_l))=0\). As \[
\operatorname{Tr}(\sigma^*,H_c^*(X,\overline{\mathbb Q}_l))
=
\operatorname{Tr}(\sigma^*,H_c^*(X,\overline{\mathbb Q}_{l'}))
\tag{3.3}
\] we may as well assume that \(l=l'\). Let us now grant the\par

\noindent\textbf{PROPOSITION 3.5.} Let \(H\) be a finite group acting freely on \(X\). Then \[
\chi(h)=\operatorname{Tr}(h^*,H_c^*(X,\overline{\mathbb Q}_l))
\] is the character of a virtual projective \(\mathbb Z_l[H]\)-module.\par

If \(H\) is abelian, and if \(H'\) is its largest subgroup of order prime to \(l\), a representation of \(H\) over \(\mathbb Z_l\) gives rise to a projective \(\mathbb Z_l[H]\)-module if and only if it is induced from a representation of \(H'\). Its character vanishes then on \(H-H'\) and we get the vanishing (b) by applying (3.5) to the cyclic group generated by \(\sigma\).\par

3.6. Let \(A\) be a torsion ring with unit and let \(\mathcal F\) be a sheaf of left \(A\)-modules on a scheme \(Y\), with \(Y\) separated and of finite type over \(k\). The \(A\)-modules \(H_c^i(Y,\mathcal F)\) are then the cohomology modules of a finer object \(R\Gamma_c(Y,\mathcal F)\) in the derived category \(D^b(A)\) of the category of \(A\)-modules. The following is the key to the proof of 3.5.\par

\noindent\textbf{PROPOSITION 3.7.} Assume \(A\) to be right and left noetherian. If \(\mathcal F\) is a constructible sheaf of projective \(A\)-modules, then \(R\Gamma_c(Y,\mathcal F)\) can be represented by a finite complex of projective \(A\)-modules of finite type.\par

We repeat the proof in [10, XVII (5.2.10)].\par
}

\EditionPage{20}{121}{%
(a) The \(H_c^i(Y,\mathcal F)\) are \(A\)-modules of finite type, and vanish for \(i>2\dim(Y)\) ([10, XVII (5.2.8.1) and (5.3.6)]).\par

(b) For any right \(A\)-module of finite type \(N\), one has \[
R\Gamma_c(Y,N\otimes_A\mathcal F)
\simeq
N\otimes_A^L R\Gamma_c(Y,\mathcal F)
\] ([10, XVII (5.2.9)]); the proof rests on replacing \(N\) by a free resolution of \(N\), to reduce to the trivial case where \(N\) is free of finite type; we are allowed to use such left infinite resolutions because \(R\Gamma_c\) is of finite cohomological dimension. The assumption that \(N\) is of finite type is in fact unnecessary.\par

(c) By (a) we can represent \(R\Gamma_c(Y,\mathcal F)\) by a complex of \(A\)-modules \(K^\bullet\), with \(K^i=0\) for \(i\notin[0,2\dim Y]\) and \(K^i\) free of finite type for \(i>0\). For any right \(A\)-module of finite type \(N\) and any \(i>0\), \[
\operatorname{Tor}_i^A(N,K^0)
=
H^{-i}(N\otimes^L R\Gamma_c(Y,\mathcal F))
\mathrel{=}_{(b)}
H^{-i}R\Gamma_c(Y,N\otimes\mathcal F)
=0.
\] The \(A\)-module \(K^0\) is hence flat; it is of finite type because \(H^0(K^\bullet)\) is, hence it is projective.\par

3.8. Proof of 3.5. We will assume that \(X\) is quasi-projective (the general case can be handled as in 3.3). Put \(Y=X/H\) and denote by \(\pi\) the projection \(\pi:X\to Y\). The group \(H\) acts on \(\pi_*\mathbb Z/l^n\) and this action turns \(\pi_*\mathbb Z/l^n\) into a locally constant sheaf of free \(\mathbb Z/l^n[H]\)-modules of rank one. By 3.7, \(R\Gamma_c(Y,\pi_*\mathbb Z/l^n)\) can be represented by a complex \(K_n^\bullet\) of projective \(\mathbb Z/l^n[H]\)-modules of finite type. One has \[
R\Gamma_c(Y,\pi_*\mathbb Z/l^n)
\simeq
R\Gamma_c(Y,\pi_*\mathbb Z/l^{n+1})
\otimes^L_{\mathbb Z/l^{n+1}}\mathbb Z/l^n
\] (3.7 (b)), and one checks easily ([11, XV, 3.3, Lemme 1]) that once \(K_n^\bullet\) is chosen, one can choose \(K_{n+1}^\bullet\) such that \(K_n^\bullet\) is the reduction mod \(l^n\) of \(K_{n+1}^\bullet\). Taking a projective limit, we get a complex \(K_\infty^\bullet\) of projective \(\mathbb Z_l[H]\)-modules and an \(H\)-equivariant isomorphism \[
H_c^*(X,\mathbb Z_l)
=
\varprojlim H_c^*(X,\mathbb Z/l^n)
=
\varprojlim H_c^*(Y,\pi_*\mathbb Z/l^n)
=
H^*(K_\infty^\bullet)
\] (the middle equality because \(\pi\) is finite). We then have \[
\chi(h)=\sum_i(-1)^i\operatorname{Tr}(h,K_\infty^i).
\]\par

3.9. Let \(X/k\) be as in (3.1), \(T\) a finite abelian group of order prime to \(p\), \(G\) a finite group and \(\rho\) an action of \(T\times G\) on \(X\). We assume that the action of \(T\) on \(X\) is free. We let \(T\times G\) act on \(H_c^*(X)\) by transport of structure. For any character \(\theta\) of \(T\) with values in \(\overline{\mathbb Q}_l^*\), we denote by \(H_\theta^*\) the subspace of \(H_c^*(X)\) on which \(T\) acts by \(\theta\); on \(H_\theta^*\), \(t^*\) is multiplication by \(\theta(t)^{-1}\).\par

Let us assume for simplicity that \(X\) is quasi-projective. We denote by \(\pi\) the projection \(\pi:X\to Y=X/T\). The group \(G\) acts on \(Y\). For each \(g\in G\), we denote by \(I(g)\) the set of connected components of the fixed point set \(Y^g\)\par
}

\EditionPage{21}{122}{%
and by \(Y_i^g\) the component corresponding to \(i\in I(g)\).\par

For \(y\in Y\), \(\pi^{-1}(y)\) is a principal homogeneous set for the action of the abelian group \(T\). If \(y\in Y^g\), \(g\) acts on \(\pi^{-1}(y)\) and commutes with \(T\), it hence acts like some \(t(g,y)\in T\): \[
gx=t(g,x)x,\qquad x\in\pi^{-1}(y);
\]  \(t(g,y)\) is constant for \(y\) in each connected component \(Y_i^g\) of \(Y^g\); we denote it by \(t(g,i)\).\par

The Theorem 3.2 will be used in the following form.\par

\noindent\textbf{COROLLARY 3.10.} With the above notations, let \(g=su\) be the decomposition of \(g\in G\) as the product of commuting elements respectively of order prime to \(p\) and a power of \(p\). Then, \[
\operatorname{Tr}(g^*,H_\theta^*)
=
\sum_{i\in I(s)}\theta(t(s,i)^{-1})\operatorname{Tr}(u^*,H_c^*(Y_i^s)).
\]\par

For \(t\in T\), the decomposition 3.1 of \(\sigma=t\cdot g\) is \(tg=(st)\cdot u\). We have \[
X^{st}
=
\coprod_{i\in I(s)}\pi^{-1}(Y_i^s)^{st}
=
\coprod_{\substack{i\in I(s)\\t=t(s,i)^{-1}}}\pi^{-1}(Y_i^s),
\] hence \[
\begin{aligned}
\operatorname{Tr}(g^*,H_\theta^*)
&=\frac1{|T|}\sum_t\theta(t)^{-1}\operatorname{Tr}(g^*t^{*-1},H_c^*(X))\\
&=\frac1{|T|}\sum_t\theta(t)\operatorname{Tr}(g^*t^*,H_c^*(X))\\
&=\frac1{|T|}\sum_t\theta(t)\sum_{\substack{i\in I(s)\\t=t(s,i)^{-1}}}
\operatorname{Tr}(u^*,H_c^*(\pi^{-1}(Y_i^s)))\\
&=\frac1{|T|}\sum_{i\in I(s)}\theta(t(s,i)^{-1})
\operatorname{Tr}(u^*,H_c^*(\pi^{-1}(Y_i^s))).
\end{aligned}
\] The decomposition (3.1) of the automorphism \(tu\) of \(\pi^{-1}(Y_i^s)\) is \(t\cdot u\). For \(t\ne e\), \(t\) has no fixed point, hence \[
\operatorname{Tr}((tu)^*,H_c^*(\pi^{-1}(Y_i^s)))=0
\qquad(t\ne e).
\] The endomorphism \[
\frac1{|T|}\sum_t t^*\quad\text{of}\quad H_c^*(\pi^{-1}(Y_i^s))
\] is a retraction onto the subspace \(\pi^*H_c^*(Y_i^s)\), hence \[
\frac1{|T|}\operatorname{Tr}(u^*,H_c^*(\pi^{-1}(Y_i^s)))
=
\operatorname{Tr}(u^*,H_c^*(Y_i^s))
\] and, substituting into the expression for \(\operatorname{Tr}(g^*,H_\theta^*)\), we get (3.10).\par

3.11. The end of this chapter will not be used in this paper. Let \(G\) be\par
}

\EditionPage{22}{123}{%
a finite group acting on \(X/k\) as in (3.1). For simplicity we assume \(X\) to be quasi-projective. Assume further that \(G\) acts freely on \(X\), and put \(Y=X/G\). The covering \(X\) of \(Y\) is said to be tame if \(Y\) can be imbedded as a Zariski open dense subset of a proper scheme \(\overline Y\) in such a way that the \(p\)-Sylow-subgroups of \(G\) act freely on the normalisation \(\overline X\) of \(\overline Y\) in \(X\): \[
\begin{array}{ccc}
X&\overset{j}{\hookrightarrow}&\overline X\\
\downarrow&&\downarrow\\
Y&\hookrightarrow&\overline Y.
\end{array}
\]\par

\noindent\textbf{PROPOSITION 3.12.} With the above notations, if \(X/Y\) is tame, then the virtual representation \(\sum_i(-1)^iH_c^i(X)\) of \(G\) is a multiple of the regular representation.\par

It suffices to prove that \(\operatorname{Tr}(g^*,H_c^*(X))=0\) for \(g\ne e\). If \(g\) is not of order a power of \(p\), this follows from (3.2). If \(g\) is of order a power of \(p\), \(g\) acts without fixed points on \(\overline X\) and \[
\operatorname{Tr}(g^*,H_c^*(X))
=
\operatorname{Tr}(g^*,H^*(\overline X,j_!\overline{\mathbb Q}_l))
=0
\] by the Lefschetz fixed point theorem applied to \(\overline X\) and the sheaf \(j_!\overline{\mathbb Q}_l\).\par

3.13. Historical remark. The method we have followed, using (3.5), was first used by Zarelua (On finite groups of transformations, Proc. Int. Symp. on Topology and its Applications, 1968, pp. 334-339) and independently by J. L. Verdier to prove that if a finite group \(G\) acts freely on a topological space \(X\) of finite cohomological dimension, and if the \(H^i(X,\mathbb Z)\) are of finite type, then \(\sum_i(-1)^iH^i(X,\mathbb Z)\otimes\mathbb Q\) is a multiple of the regular representation of \(G\).\par

\EditionChapter{4. The character formula}

The assumptions (0.1.1) are in force in this chapter and the next four.\par

\noindent\textbf{DEFINITION 4.1.} Let \(T\) be an \(F\)-stable maximal torus of \(G\). The Green function \(Q_{T,G}(u)\) is the restriction to the unipotent elements of the character of the virtual representation \(R_{T\subset B}^1\), where \(B\) is any Borel subgroup containing \(T\).\par

This Green function does not depend on \(B\) (by (1.6) and (1.17 (i))); it only depends on the \(G^F\)-conjugacy classes of \(T\) and \(u\). The natural map \(x\mapsto\overline x\) from \(G\) to its adjoint group induces a bijection on the unipotent sets, and \[
\tag{4.1.1}
Q_{T,G}(u)=Q_{\overline T,G^{\mathrm{ad}}}(\overline u).
\] The Green function is integer-valued (3.3) and is a restriction of a character\par
}

\EditionPage{23}{124}{%
of \(G^F\), hence \(Q_{T,G}(u)=Q_{T,G}(u^n)\) if \((n,p)=1\).\par

Let \(\alpha\) be the smallest integer \(\ge1\) such that \(F^\alpha\) is the identity on \(\mathbf W\). From the proof of (3.3) we see that \(\sum_{n\ge1}X_{T\subset B}^{F^{\alpha n}u}t^n\) is a rational function of \(t\) and that \[
\tag{4.1.2}
Q_{T,G}(u)
=-\left\{\sum_{n\ge1}X_{T\subset B}^{F^{\alpha n}u}t^n\right\}_{t=\infty}.
\]\par

In this chapter, we will express the character of \(R_{T\subset B}^\theta\) in terms of \(\theta\) and of Green functions.\par

\noindent\textbf{THEOREM 4.2.} Let \(x=su\) be the Jordan decomposition of \(x\in G^F\). Then \[
\operatorname{Tr}(x,R_{T\subset B}^\theta)
=
\frac1{|Z^0(s)^F|}
\sum_{\substack{g\in G^F\\\operatorname{ad}gT\subset Z^0(s)}}
Q_{\operatorname{ad}gT,Z^0(s)}(u)\operatorname{ad}g(\theta)(s).
\]\par

Let \(\varepsilon_x\) be the function on \(G^F\) whose value is 1 at \(x\in G^F\) and 0 elsewhere. One can rewrite 4.2 as giving the following formula for the character of \(R_{T\subset B}^\theta\): \[
\tag{4.2.1}
\operatorname{Tr}(\ ,R_{T\subset B}^\theta)
=
\sum_{t\in T^F}\frac1{|Z^0(t)^F|}\theta(t)
\sum_{\substack{u\in Z^0(t)^F\\g\in G^F}}
Q_{T,Z^0(t)}(u)\varepsilon_{\operatorname{ad}g(tu)}
\] where we put \(Q_{T,Z^0(t)}(u)=0\) whenever \(u\) is not unipotent.\par

\noindent\textbf{COROLLARY 4.3.} \(R_{T\subset B}^\theta\) is independent of the choice of \(B\) (\(B\supset T\)).\par

From now on we shall write \(R_T^\theta\) for \(R_{T\subset B}^\theta\). We will deduce 4.2 from 3.7 and the following geometrical facts.\par

\noindent\textbf{PROPOSITION 4.4.} (i) If a Borel subgroup \(B_1\) of \(G\) contains \(s\), then \(B_1\cap Z^0(s)\) is a Borel subgroup of \(Z^0(s)\).\par

(ii) Pick \(T\subset B\) in \(G\). Any Borel subgroup \(B_1\) such that \(s\in B_1\) is of the form \(\operatorname{ad}gB\) with \(\operatorname{ad}gT\subset Z^0(s)\) (i.e., \(g^{-1}sg\in T\)). The left coset \[
\tau_{T,B}(B_1)\mathrel{\mathop{=}^{\mathrm{def.}}}Z^0(s)\cdot g
\] depends only on \(T\), \(B\), and \(B_1\).\par

(iii) The map \(B_1'\mapsto B_1'\cap Z^0(s)\) is an isomorphism from the space of Borel subgroups containing \(s\), such that \(\tau_{T,B}(B_1')=\tau_{T,B}(B_1)\) and the space of Borel subgroups of \(Z^0(s)\).\par

Proof. (i) is well known. Put \(B_1=\operatorname{ad}g_1B\). We have \(g_1^{-1}sg_1\in B\) hence there exists \(u\) in the unipotent radical of \(B\) such that \(u^{-1}g_1^{-1}sg_1u\in T\). We take \(g=g_1u\). Put \(T'=\operatorname{ad}gT\). If \(g'=hg\) is also such that \(\operatorname{ad}g'B=B_1\) (resp. \(\operatorname{ad}g'T\subset Z^0(s)\)) then \(h\in B_1\) (resp. \(h\in Z^0(s)N(T')\), by the conjugacy of maximal tori in \(Z^0(s)\)). If \(W=N(T')/T'\) is the Weyl group of \(G\) and \(W^0=(N(T')\cap Z^0(s))/T'\) that of \(Z^0(s)\), the Bruhat decomposition for \(Z^0(s)\) reads \[
Z^0(s)=(B_1\cap Z^0(s))W^0(B_1\cap Z^0(s)),
\] and\par
}

\EditionPage{24}{125}{%
\[
(B_1Z^0(s))\cap N(T')
=
B_1W^0(B_1\cap Z^0(s))\cap N(T')
\subset
B_1W^0B_1\cap N(T')
=
W^0(T')
\subset Z^0(s).
\] Hence \[
B_1\cap(Z^0(s)\cdot N(T'))\subset Z^0(s)
\] and \[
B_1\cap(Z^0(s)\cdot N(T'))=B_1\cap Z^0(s).
\] The element \(g\) such that \(\operatorname{ad}gB=B_1\) and \(\operatorname{ad}gT\subset Z^0(s)\) is unique modulo \(B_1\cap Z^0(s)\), hence (ii) and (iii).\par

A \(Z^0(s)\)-left coset \(\tau=Z^0(s)g\), with \(\operatorname{ad}gT\subset Z^0(s)\), defines an isomorphism \(\overline\tau\) from the maximal torus of \(G\) to that of \(Z^0(s)\); \(\operatorname{ad}g\) maps \(T\) (contained in the Borel subgroup \(B\)) to \(\operatorname{ad}gT\subset Z^0(s)\) (contained in the Borel subgroup \(\operatorname{ad}gB\cap Z^0(s)\)). This construction will allow us to compare the relative position of Borel subgroups in \(G\) and \(Z^0(s)\).\par

\noindent\textbf{PROPOSITION 4.5.} Let \(B_1'\) and \(B_1''\) be Borel subgroups containing \(s\), \[
\tau'=\tau_{T,B}(B_1'),\qquad \tau''=\tau_{T,B}(B_1'');
\] let \(w\) be the relative position of \(B_1'\) and \(B_1''\) (in \(G\)) and \(w_1\) that of \(B_1'\cap Z^0(s)\) and \(B_1''\cap Z^0(s)\) (in \(Z^0(s)\)). Then \[
w_1=\overline\tau' w\overline\tau''^{-1}.
\]\par

Replacing \((T,B)\) by a conjugate, we may assume that \(B=B_1'\) and that \(T\subset B_1'\cap Z^0(s)\cap B_1''\). We will identify the maximal torus of \(G\) and that of \(Z^0(s)\) with \(T\), using \(T\subset B\) and \(T\subset(B\cap Z^0(s))\). We then have \(\overline\tau'=\operatorname{Id}\). Replacing \(B_1''\) by an \(N(T)\cap Z^0(s)\)-conjugate, we may further assume that \(w_1=1\), i.e., \(B_1'\cap Z^0(s)=B_1''\cap Z^0(s)\). We must prove that \(w=\overline\tau''\), which is clear.\par

\noindent\textbf{LEMMA 4.6.} If \(B\) and \(B'\), containing \(T\), are in the same relative position as \(B_1\) and \(B_1'\) containing \(s\), then \(\tau_{T,B}(B_1)=\tau_{T,B'}(B_1')\).\par

Replacing \((B,T,B')\) by a conjugate, we may again assume that \(B=B_1\) and that \(T\subset B_1\cap Z^0(s)\cap B_1'\). In this case \(B'=B_1'\).\par

To prove (4.2) we first compute the space \(X_{T\subset B}^s\) of Borel subgroups \(B_1\) containing \(s\) and such that \(B_1\) and \(FB_1\) are in the same relative position as \(B\) and \(FB\).\par

\noindent\textbf{PROPOSITION 4.7.} Let \(X_{T\subset B}^s(g)\) be the subspace of \(X_{T\subset B}^s\) consisting of those \(B_1\) with \(\tau_{T,B}(B_1)=Z^0(s)g\). Then \[
\tag{4.7.1}
X_{T\subset B}^s
=
\coprod_{\substack{g\in Z^0(s)^F\backslash G^F\\\operatorname{ad}gT\subset Z^0(s)}}
X_{T\subset B}^s(g)
\] and, for \(g\in G^F\), the map \(B_1\mapsto B_1\cap Z^0(s)\) induces an isomorphism\par
}

\EditionPage{25}{126}{%
\[
\tag{4.7.2}
X_{T\subset B}^s(g)\xrightarrow{\sim}
X_{\operatorname{ad}gT\subset \operatorname{ad}gB\cap Z^0(s)}.
\]

By 4.6, we have \(F\tau_{T,B}(B_1)=\tau_{T,FB}(FB_1)=\tau_{T,B}(B_1)\). By Lang's Theorem, the \(Z^0(s)\)-principal homogeneous space \(\tau_{T,B}(B_1)\), being \(F\)-stable, has a rational point, hence (4.7.1). By (4.4 (iii)) and (4.5), the map (4.7.2) is an isomorphism from \(X_{T\subset B}^s(g)\) to some \(X(w)\) (relative to \(Z^0(s)\)); to check which \(w\) appears, we observe that \(\operatorname{ad}gB\in X_{T\subset B}^s(g)\).\par

Proof of 4.2. The theorem is now an immediate application of (4.7) and (3.10) applied to the action (1.17) of \(G^F\times T^F\) on \(\widetilde X_{T\subset B}\).\par

\EditionChapter{5. Characters of tori}

We will assume chosen isomorphisms \[
\tag{5.0.1}
k^*\xrightarrow{\sim}(\mathbb Q/\mathbb Z)_{p'},
\] and \[
\tag{5.0.2}
(\text{roots of unity of order prime to }p\text{ in }\overline{\mathbb Q}_l^*)
\xrightarrow{\sim}(\mathbb Q/\mathbb Z)_{p'}.
\]\par

5.1. Let \(T\) be a torus over \(k\). Besides the character group \(X(T)=\operatorname{Hom}(T,\mathbb G_m)\), we will consider its dual \(Y(T)=\operatorname{Hom}(\mathbb G_m,T)\). The duality is given by \[
\langle\alpha,h\rangle=\alpha\circ h\in\operatorname{Hom}(\mathbb G_m,\mathbb G_m)=\mathbb Z
\] (\(\alpha\in X(T)\) and \(h\in Y(T)\)). The maps \((h,x)\mapsto h(x)\) and \(t\mapsto(\alpha\mapsto\alpha(t))\) induce isomorphisms \[
\tag{5.1.1}
Y(T)\otimes k^*=T(k)=\operatorname{Hom}(X(T),k^*).
\]\par

If \(T\) is a maximal torus in \(G\), the roots of \(T\) in \(G\) are elements of \(X(T)\), while the coroots belong to \(Y(T)\). If \(\alpha\) is a root, the corresponding coroot \(H_\alpha\) is characterised as follows: there is a homomorphism \(u:\operatorname{SL}(2)\to G\), which maps the subgroup \(\begin{pmatrix}1&*\\0&1\end{pmatrix}\) onto the root subgroup \(U_\alpha\), and whose restriction to the group of diagonal matrices (identified with \(\mathbb G_m\) by \(x\mapsto\begin{pmatrix}x&0\\0&x^{-1}\end{pmatrix}\)) is \(H_\alpha\).\par

5.2. Using the chosen isomorphism (5.0.1), we can rewrite (5.1.1) as \[
\tag{5.2.1}
T(k)=Y(T)\otimes(\mathbb Q/\mathbb Z)_{p'}\simeq(Y(T)\otimes\mathbb Q/Y(T))_{p'}.
\] If \(T\) is obtained by extension of scalars from a torus \(T_0/\mathbb F_q\), the Frobenius map \(F\) induces a map \(F:Y(T)\to Y(T)\) (\(Y\) is a covariant functor) and \(T^F\) is the subgroup \(\operatorname{Ker}(F-1)\) of \(T(k)=Y(T)\otimes(\mathbb Q/\mathbb Z)_{p'}\); since \(F\) is divisible by \(p\), it is also the kernel of \(F-1\) in \(Y(T)\otimes\mathbb Q/\mathbb Z\):\par
}

\EditionPage{26}{127}{%
\[
\tag{5.2.2}
0\longrightarrow T^F\longrightarrow Y(T)\otimes\mathbb Q/\mathbb Z
\overset{F-1}{\longrightarrow}Y(T)\otimes\mathbb Q/\mathbb Z\longrightarrow0.
\]

By applying the Snake Lemma to \[
\begin{array}{ccccccccc}
0&\longrightarrow&Y(T)&\longrightarrow&Y(T)\otimes\mathbb Q&\longrightarrow&Y(T)\otimes\mathbb Q/\mathbb Z&\longrightarrow&0\\
&&{\scriptstyle F-1}\downarrow&&{\scriptstyle F-1}\downarrow&&{\scriptstyle F-1}\downarrow&&\\
0&\longrightarrow&Y(T)&\longrightarrow&Y(T)\otimes\mathbb Q&\longrightarrow&Y(T)\otimes\mathbb Q/\mathbb Z&\longrightarrow&0,
\end{array}
\] we get another exact sequence \[
\tag{5.2.3}
0\longrightarrow Y(T)\overset{F-1}{\longrightarrow}Y(T)\longrightarrow T^F\longrightarrow0.
\]\par

The maps in (5.2.2), (5.2.3) depend on the choice (5.0.1). An intrinsic expression for these sequences would be as follows:\par

\(T^F\) is the image of the middle map in the exact sequence \[
0\longrightarrow Y(T)_{\widehat{p'}}(1)
\overset{F-1}{\longrightarrow}Y(T)_{\widehat{p'}}(1)
\overset{(F-1)^{-1}}{\longrightarrow}Y(T)_{\widehat{p'}}(1)\otimes\mathbb Q/\mathbb Z
\overset{F-1}{\longrightarrow}Y(T)_{\widehat{p'}}(1)\otimes\mathbb Q/\mathbb Z
\longrightarrow0.
\]\par

Let us map the sequences (5.2.2), (5.2.3) into \(\mathbb Q/\mathbb Z\). The isomorphism (5.0.2) provides an isomorphism \[
(T^F)^{\sim}\mathrel{\mathop{=}^{\mathrm{def}}}
\operatorname{Hom}(T^F,\overline{\mathbb Q}_l^*)
=
\operatorname{Hom}(T^F,\mathbb Q/\mathbb Z),
\] and we get exact sequences \[
\tag{5.2.2}^*
0\longrightarrow X(T)\overset{F-1}{\longrightarrow}X(T)\longrightarrow(T^F)^{\sim}\longrightarrow0,
\]  \[
\tag{5.2.3}^*
0\longrightarrow(T^F)^{\sim}\longrightarrow X(T)\otimes\mathbb Q/\mathbb Z
\overset{F-1}{\longrightarrow}X(T)\otimes\mathbb Q/\mathbb Z\longrightarrow0.
\]\par

The dual torus \(T^*\) of \(T\) is defined by the rule \(X(T^*)=Y(T)\), (hence \(Y(T^*)=X(T)\)); its \(\mathbb F_q\)-structure is defined by \((F\text{ on }Y(T^*))=({}^tF\text{ on }X(T))\). Via some isomorphism \[
\tag{5.2.4}
(T^F)^{\sim}=T^{*F},
\] the sequences (5.2.3), (5.2.2) for \(T^*\) are identical with (5.2.2)\(^*\), (5.2.3)\(^*\).\par

5.3. If we go from \(\mathbb F_q\) to \(\mathbb F_{q^n}\), \(F\) is replaced by \(F^n\). Composition with the norm map \[
N=\frac{F^n-1}{F-1}=\sum_{0}^{n-1}F^i:T^{F^n}\longrightarrow T^F
\] is an injection \({}^tN:(T^F)^{\sim}\to(T^{F^n})^{\sim}\). We have commutative diagrams\par
}

\EditionPage{27}{128}{%
\[
\begin{array}{ccccccccccccc}
0&\longrightarrow&Y&\overset{F^n-1}{\longrightarrow}&Y&\longrightarrow&T^{F^n}&\longrightarrow&Y\otimes\mathbb Q/\mathbb Z&\overset{F^n-1}{\longrightarrow}&Y\otimes\mathbb Q/\mathbb Z&\longrightarrow&0\\
&&{\scriptstyle N}\downarrow&&\Vert&&{\scriptstyle N}\downarrow&&{\scriptstyle N=\sum_{0}^{n-1}F^i}\downarrow&&\Vert&&\\
0&\longrightarrow&Y&\longrightarrow&Y&\longrightarrow&T^F&\longrightarrow&Y\otimes\mathbb Q/\mathbb Z&\longrightarrow&Y\otimes\mathbb Q/\mathbb Z&\longrightarrow&0,
\end{array}
\]
\[
\begin{array}{ccccccccccccc}
0&\longrightarrow&X&\longrightarrow&X&\longrightarrow&(T^F)^{\sim}&\longrightarrow&X\otimes\mathbb Q/\mathbb Z&\longrightarrow&X\otimes\mathbb Q/\mathbb Z&\longrightarrow&0\\
&&\Vert&&{\scriptstyle {}^tN=\sum_{0}^{n-1}{}^tF^i}\downarrow&&{\scriptstyle {}^tN}\downarrow&&\Vert&&{\scriptstyle {}^tN}\downarrow&&\\
0&\longrightarrow&X&\longrightarrow&X&\longrightarrow&(T^{F^n})^{\sim}&\longrightarrow&X\otimes\mathbb Q/\mathbb Z&\longrightarrow&X\otimes\mathbb Q/\mathbb Z&\longrightarrow&0.
\end{array}
\]
where \(X\) and \(Y\) stand for \(X(T)\) and \(Y(T)\). In particular, for \(\theta\) a character of \(T^F\), \(\theta\) and \(\theta\circ N\) have the same image in \(X(T)\otimes\mathbb Q/\mathbb Z\), and (5.2.4) gives rise to a commutative diagram
\[
\begin{array}{ccc}
(T^F)^{\sim}&\simeq&(T^*)^F\\
{\scriptstyle {}^tN}\downarrow&&{\scriptstyle N}\downarrow\\
(T^{F^n})^{\sim}&\simeq&(T^*)^{F^n}.
\end{array}
\]

\noindent\textbf{PROPOSITION 5.4.} Let \(T\) and \(T'\) be two \(F\)-stable maximal tori of \(G\), and let \(\theta,\theta'\) be characters of \(T^F,(T')^F\). We identify them with characters of \(Y(T)\) and \(Y(T')\), by (5.2.3). The following conditions are equivalent:\par

(i) For some \(x\in G\), with \(\operatorname{ad}x(T)=T'\), the map induced by \(\operatorname{ad}x:Y(T)\to Y(T')\) carries \(\theta\) to \(\theta'\);\par

(ii) For some \(n\), the pairs \((T,\theta\circ N)\), \((T',\theta'\circ N)\), where \(N\) is the norm from \(T^{F^n}\) to \(T^F\) (resp. \((T')^{F^n}\) to \((T')^F\)) are \(G^{F^n}\)-conjugate.\par

By 5.3, the condition (i) is invariant under the replacement of \(F\) by \(F^n\) and of \(\theta,\theta'\) by \(\theta\circ N,\theta'\circ N\). It hence suffices to check (5.4) in the trivial case where \(T\) and \(T'\) are split.\par

\noindent\textbf{DEFINITION 5.5.} The pairs \((T,\theta)\), \((T',\theta')\) are said to be geometrically conjugate when the equivalent conditions of (5.4) hold.\par

5.6. Let \((T,\theta)\) be as above. The choice of a Borel subgroup \(B\) containing \(T\) defines an isomorphism of \(T\) with the torus \(\mathbf T\); the corresponding isomorphism \(Y(T)\to Y(\mathbf T)\) carries \(\theta\) to a character of \(Y(\mathbf T)\) with values in the roots of unity of order prime to \(p\) in \(\overline{\mathbb Q}_l\). The isomorphism (5.0.2) identifies it with an element of \(X(\mathbf T)\otimes(\mathbb Q/\mathbb Z)_{p'}\). Its class \([\theta]\) in \((X(\mathbf T)\otimes(\mathbb Q/\mathbb Z)_{p'})/\mathbf W\) is independent of the choice of \(B\). It is invariant under \(F\). Let us put \[
\mathcal S\mathrel{\mathop{=}^{\mathrm{def}}}
[(X(\mathbf T)\otimes(\mathbb Q/\mathbb Z)_{p'})/\mathbf W]^F
\xrightarrow{\sim}
[(X(\mathbf T)\otimes\mathbb Q/\mathbb Z)/\mathbf W]^F.
\]\par

\noindent\textbf{PROPOSITION 5.7.} (i) The map \(\theta\mapsto[\theta]\) induces a bijection from the set of geometric conjugacy classes of pairs \((T,\theta)\) to \(\mathcal S\).\par
}

\EditionPage{28}{129}{%
(ii) (Compare Steinberg [15, p. 93]). The number of geometric conjugacy classes of pairs \((T,\theta)\) is \(|Z^{0F}|q^l\) where \(Z^0\) is the identity component of the centre of \(G\) and \(l\) is the semisimple rank of \(G\).\par

(i) The injectivity is clear on 5.4(i). Let us prove surjectivity. If the class mod \(\mathbf W\) of \(x\in X(\mathbf T)\otimes(\mathbb Q/\mathbb Z)_{p'}\) is invariant by \(F\), one indeed has \({}^t(wF)x=x\) for some \(w\in\mathbf W\), and \(x\) corresponds to a character of \(\mathbf T(w)^F\), hence of \(T^F\), for some \(F\)-stable maximal torus \(T\) in \(G\) (1.14).\par

(ii) As in (5.2), we may now identify geometric conjugacy classes of pairs \((T,\theta)\) with \(\mathbb F_q\)-rational points of the quotient \(\mathbf T^*/\mathbf W\). Our task is to compute \[
|(\mathbf T^*/\mathbf W)^F|
=
\sum_i(-1)^i\operatorname{Tr}(F^*,H_c^i(\mathbf T^*/\mathbf W))
=
\sum_i(-1)^i\operatorname{Tr}(F^*,H_c^i(\mathbf T^*)^{\mathbf W}).
\] Let \(G'\) be the derived group of \(G\); the torus \(\mathbf T^*\) is isogenous to (hence has the same cohomology as) the product of \(Z^{0*}\) with \(\mathbf T'^*\) (\(\mathbf T'\) being the torus of \(G'\)). Since any torus has as many rational points as its dual, we get by Künneth's Theorem \[
|(\mathbf T^*/\mathbf W)^F|
=
|Z^{0F}|\sum_i(-1)^i\operatorname{Tr}(F^*,H_c^i(\mathbf T'^*)^{\mathbf W}).
\] We have \[
H^*(\mathbf T'^*,\overline{\mathbb Q}_l)
=
\Lambda^*H^1(\mathbf T'^*,\overline{\mathbb Q}_l)
=
\Lambda^*(Y(\mathbf T')\otimes\overline{\mathbb Q}_l(-1));
\] hence by Poincaré duality, \[
H_c^{2l-i}(\mathbf T'^*,\overline{\mathbb Q}_l)
=
\operatorname{Hom}(\Lambda^iY(\mathbf T'),\overline{\mathbb Q}_l(i-l)),
\] and \[
\operatorname{Tr}(F^*,H_c^{2l-i}(\mathbf T'^*)^{\mathbf W})
=
q^{l-i}\operatorname{Tr}(F,\Lambda^iY(\mathbf T')^{\mathbf W}).
\] The next lemma shows that this trace is zero for \(i\ne0\), hence the \(q^l\) factor.\par

\noindent\textbf{LEMMA 5.8.} Let \(W\subset\operatorname{Aut}(V)\) be a finite group generated by reflections of a real vector space \(V\) of dimension \(d\). We assume that \(V^W=0\). Then \((\Lambda^iV)^W=0\) for \(i\ne0\).\par

For a proof, see Bourbaki [1, Ch. V, Ex. 3 of § 2].\par

5.9. Let \(T\) be an \(F\)-stable maximal torus in \(G\) and let \(\alpha\) be a root of \(T\). For \(\theta\) a character of \(T^F\), we will say that \(\theta\) is orthogonal to \(H_\alpha\): \(\langle H_\alpha,\theta\rangle=0\) if \(\theta\), viewed as a character of \(Y(T)\) (5.2.3), is identically 1 on \(H_\alpha\). The character \(\theta\) of \(Y(T)\) is then invariant by the corresponding reflection \(s_\alpha\).\par

5.10. Let \(\pi:\widetilde G\to G\) be the simply connected covering of the derived group of \(G\), and let \(\widetilde T\) be the inverse image of \(T\) in \(\widetilde G\). The group \(\widetilde T\) is\par
}

\EditionPage{29}{130}{%
connected: it is a maximal torus in \(\widetilde G\). It acts on \(T\times\widetilde G\) by \[
\widetilde t*(t,g)=(t\pi(\widetilde t)^{-1},\widetilde t g),
\] and the map \((t,g)\mapsto tg\) induces an isomorphism \(\widetilde T\backslash(T\times\widetilde G)\xrightarrow{\sim}G\), hence (Lang's Theorem) \(\widetilde T^F\backslash(T^F\times\widetilde G^F)\xrightarrow{\sim}G^F\). We have \[
\tag{5.10.1}
T^F/\pi(\widetilde T^F)\xrightarrow{\sim}G^F/\pi^F(\widetilde G).
\]\par

\noindent\textbf{PROPOSITION 5.11.} (i) A character of \(T^F\) is the restriction to \(T^F\) of a character of \(G^F/\pi(\widetilde G^F)\) if and only if it is orthogonal to all coroots.\par

(ii) Let \(\theta\) be a character of \(G^F/\pi(\widetilde G^F)\), \(T\) and \(T'\) two \(F\)-stable maximal tori and let \(x\in G\) be such that \(\operatorname{ad}x(T)=T'\). Then \(\operatorname{ad}x:Y(T)\to Y(T')\) carries the restriction of \(\theta\) to \(T^F\) (viewed as a character of \(Y(T)\)) onto the restriction of \(\theta\) to \((T')^F\): the restrictions of \(\theta\) to the \(F\)-stable maximal tori are all geometrically conjugate.\par

The assertion (i) follows from (5.10.1), and the fact that \(Y(\widetilde T)\subset Y(T)\) is spanned by the coroots \(H_\alpha\). To prove (ii), we will use the criterion 5.4 (ii). Let \(n\) be such that \(x\in G^{F^n}\). We will prove that if \(y\in T^{F^n}\), \(N(y)\) and \(N(\operatorname{ad}x(y))\) have the same image in \(G^F/\pi(\widetilde G^F)\), hence that \(\operatorname{ad}x\) transforms \((\theta\mid T^F)\circ N\) into \((\theta\mid(T')^F)\circ N\).\par

The map \[
t\longmapsto N(t)^{-1}N(\operatorname{ad}x(t)):T\longrightarrow G
\] lifts uniquely into a map \(\varphi:T\to\widetilde G\) mapping \(e\) to \(e\). Indeed,\par

(a) it factors through \(T/Z\);\par

(b) the map \(t\mapsto N(t)^{-1}N(\operatorname{ad}x(t)):\widetilde T\to\widetilde G\) factors through \(\widetilde T/\widetilde Z=T/Z\) and the desired lifting is the composite \(T\to T/Z=\widetilde T/\widetilde Z\to\widetilde G\).\par

Similarly, the map \[
(a,t)\longmapsto a^{-1}N(t)^{-1}N(\operatorname{ad}x(t))\operatorname{ad}x(a):T\times T\longrightarrow G
\] lifts into \(\psi:T\times T\to\widetilde G\), with \(\psi(e,t)=\varphi(t)\). The identity \[
\begin{aligned}
F(N(t)^{-1}N(\operatorname{ad}x(t)))
&=N(Ft)^{-1}N(F\operatorname{ad}x(t))\\
&=(t^{-1}F^nt)^{-1}N(t)^{-1}N(\operatorname{ad}x(t))\operatorname{ad}x(t^{-1})F^n\operatorname{ad}x(t)\\
&=(t^{-1}F^nt)^{-1}N(t)^{-1}N(\operatorname{ad}x(t))\operatorname{ad}x(t^{-1}F^nt)
\end{aligned}
\] lifts into \[
F\varphi(t)=\psi(t^{-1}F^nt,\varphi(t)).
\] Putting \(t=y\), we have \(y^{-1}F^ny=e\), hence \(F\varphi(y)=\varphi(y)\) and \(\varphi(y)\in\widetilde G^F\). We have \(N(y)^{-1}N(\operatorname{ad}x(y))\in\pi(\widetilde G^F)\), as required.\par

Remark 5.12. The argument above can be used to show the existence of a norm map \[
N:G^{F^n}/\pi(\widetilde G^{F^n})\longrightarrow G^F/\pi(\widetilde G^F)
\] which for each \(F\)-stable maximal torus \(T\) gives rise to a commutative diagram\par
}

\EditionPage{30}{131}{%
\[
\begin{array}{ccc}
T^{F^n}/\pi(\widetilde T^{F^n})&\xrightarrow{\sim}&G^{F^n}/\pi(\widetilde G^{F^n})\\
{\scriptstyle N}\downarrow&&{\scriptstyle N}\downarrow\\
T^F/\pi(\widetilde T^F)&\xrightarrow{\sim}&G^F/\pi(\widetilde G^F).
\end{array}
\]

\noindent\textbf{THEOREM 5.13.} Let \((T,\theta)\) be as above. We use (5.2.3) to identify \(\theta\) with a character of \(Y(T)\). If the centre of \(G\) is connected, then the stabilizer of \(\theta\) in the Weyl group \(W=N(T)/T\) is generated by the reflections \(s_\alpha\), for \(H_\alpha\) a coroot orthogonal to \(\theta\).\par

The key point in the proof is the same as in Steinberg's Theorem that if the derived group \(G'\) is simply connected, then the centralizer of any semisimple element is connected.\par

Let us identify \(\theta\) with the corresponding element in \(X(T)\otimes(\mathbb Q/\mathbb Z)_{p'}\) (invariant under \(F\)). Let \(\theta_1\in X(T)\otimes\mathbb Q\) be a representative for it. This \(\theta_1\) has no \(p\) in the denominator. We have the dictionary:\par

(a) \(\langle H_\alpha,\theta\rangle=0\Longleftrightarrow\langle H_\alpha,\theta_1\rangle\in\mathbb Z\);\par

(b) \(\theta\) is fixed by \(w\in W\Longleftrightarrow\) for some \(x\in X(T)\), \(w\theta_1+x=\theta_1\).\par

Let \(X^{\mathrm{ad}}\) be the subgroup of \(X\) generated by the roots. It is the character group of the image \(T^{\mathrm{ad}}\) of \(T\) in the adjoint group. The character group of the centre \(Z=\operatorname{Ker}(T\to T^{\mathrm{ad}})\) of \(G\) is \(X/X^{\mathrm{ad}}\); the character group of \(Z/Z_{\mathrm{red}}^0\) is the torsion subgroup of \(X/X^{\mathrm{ad}}\), hence\par

(c) \(Z\) is connected and smooth (resp. connected) if and only if \(X/X^{\mathrm{ad}}\) is torsion free (resp. has no torsion prime to \(p\)).\par

For each root \(\alpha\), \(s_\alpha(\theta_1)=\theta_1-\langle H_\alpha,\theta_1\rangle\alpha\). The number \(\langle H_\alpha,\theta_1\rangle\) has no \(p\) in the denominator; hence \(\langle H_\alpha,\theta_1\rangle\alpha\in X(T)\) if and only if \(\langle H_\alpha,\theta_1\rangle\in\mathbb Z\); \(s_\alpha(\theta)=\theta\) if and only if \(\langle H_\alpha,\theta\rangle=0\). It remains to check that the stabilizer of \(\theta\) is generated by the reflections it contains.\par

Fix \(N\) such that \((1/p^N)X^{\mathrm{ad}}\supset X\cap(X^{\mathrm{ad}}\otimes\mathbb Q)\) and such that \(p^N\equiv1\pmod{\text{the order of }\theta}\). For any \(w\in W\), \(w\theta_1-\theta_1\) is in \(X^{\mathrm{ad}}\otimes\mathbb Q\), hence if \(w\theta_1-\theta_1\) is in \(X\), \(w(p^N\theta_1)-(p^N\theta_1)\in X^{\mathrm{ad}}\). Replacing \(\theta_1\) by \(p^N\theta_1\), which is also a representative of \(\theta\), we may assume that \(\theta\) is fixed by \(w\in W\) if and only if \(w\theta_1+x=\theta_1\) for some \(x\in X^{\mathrm{ad}}\). The stabilizer of \(\theta\) in \(W\) is the image in \(W\) of the stabilizer of \(\theta_1\) in the affine Weyl group. It remains to apply Bourbaki [1, Ch. VI, Ex. 1 of § 2].\par

Remark 5.14. The proof shows that, if we use a suitable Borel subgroup \(B\supset T\) to identify \(W\) with \(\mathbf W\), the stabilizer of \(\theta\) becomes the subgroup of \(\mathbf W\) generated by some of the reflections corresponding either to a simple root or to the negative of the highest coroot.\par

\noindent\textbf{DEFINITION 5.15.} (i) The character \(\theta\) of \(T^F\) is nonsingular if it is not\par
}

\EditionPage{31}{132}{%
orthogonal to any coroot.\par

(ii) \(\theta\) is in general position if it is not kept fixed by any non-trivial element of \((N(T)/T)^F\).\par

\noindent\textbf{PROPOSITION 5.16.} If the centre of \(G\) is connected, a character is non-singular if and only if it is in general position.\par

The stabilizer of \(\theta\) in \(N(T)/T\) is a group generated by reflections, and is stable by Frobenius. We have to prove that the subgroup fixed by Frobenius is non-trivial. This follows from the next lemma.\par

\noindent\textbf{LEMMA 5.17.} Let \(V\ne\{0\}\) be a euclidian vector space, and let \(W\) be a finite group generated by orthogonal reflections. We assume that \(V^W=\{0\}\). If \(a\) belongs to the normalizer \(A\) of \(W\) in \(O(V)\), the centralizer \(W^a\) of \(a\) in \(W\) is non-trivial.\par

(a) Reduction to the irreducible case. Let \(V=\bigoplus V_i\) be the decomposition of \(V\) as a direct sum of irreducible root systems; we have \(W=\prod W_i\). The automorphism \(a\) permutes the \(V_i\); by taking a direct factor, we may assume that it permutes them cyclically: \(V=\bigoplus_{i\in\mathbb Z/n}V_i\) and \(aV_i=V_{i+1}\). An element \(w=(w_i)\in W=\prod W_i\) is in \(W^a\) if and only if \(w_i=\operatorname{ad}a^i(w_0)\) and \(w_0\) is fixed by \(a^n\); we are reduced to proving that \(W_0^{a^n}\) is non-trivial.\par

(b) In the irreducible case, one always has either \(A=W\cup -W\) or \(-1\in W\); both cases are clear.\par

\noindent\textbf{COROLLARY 5.18.} For any \(G\), if \(\theta\) is in general position, then \(\theta\) is non-singular.\par

Let us embed \(G\) in a group \(G_1\) with connected centre and the same derived group. For instance, one can take \(G_1=G\times T/\{(z,z^{-1})\mid z\in Z(G)\}\). The torus \(T\) is then contained in an \(F\)-stable maximal torus \(T_1\) of \(G_1\), and \(\theta\) is the restriction to \(T^F\) of some character \(\theta_1\) of \(T_1^F\). If \(\theta\) is in general position, \(\theta_1\) is so a fortiori, hence non-singular (5.16). A character \(\theta_1\) of \(T_1^F\) is non-singular if and only if its restriction to \(T^F\) is, hence (5.18).\par

When all roots of \(G\) have the same length, geometric conjugacy can be given the following convenient description.\par

\noindent\textbf{DEFINITION 5.19.} (Assuming all roots to be of the same length) Let \(T\) be an \(F\)-stable maximal torus and let \(\theta\) be a character of \(T^F\). The connected centralizer \(S(T,\theta)\) of \((T,\theta)\) in \(G\) is the \(F\)-stable reductive subgroup of \(G\) with maximal torus \(T\) and with root subgroups relative to \(T\) the \(U_\alpha\) for which \(\langle H_\alpha,\theta\rangle=0\).\par

If \(\langle H_\alpha,\theta\rangle=\langle H_\beta,\theta\rangle=0\) and \(\alpha+\beta\) is a root, then \(H_{\alpha+\beta}=H_\alpha+H_\beta\),\par
}

\EditionPage{32}{133}{%
hence \(\langle H_{\alpha+\beta},\theta\rangle=0\) and the definition makes sense. Let \(\pi:\widetilde{S(T,\theta)}\to S(T,\theta)\) be the simply connected covering of the derived subgroup of \(S(T,\theta)\). By 5.11 (i), \(\theta\) extends to a character \(\theta_S\) of \(S(T,\theta)^F/\pi(\widetilde{S(T,\theta)}^F)\).\par

\noindent\textbf{PROPOSITION 5.20.} If the centre of \(G\) is connected and all roots are of the same length, then a pair \((T',\theta')\) is geometrically conjugate to \((T,\theta)\) if and only if \((S(T,\theta),\theta_S)\) and \((S(T',\theta'),\theta'_{S'})\) are \(G^F\)-conjugate.\par

The “if” part follows from 5.11. Conversely, let \((T',\theta')\) be geometrically conjugate to \((T,\theta)\): for some \(x\in G\), \(\operatorname{ad}x(T)=T'\) and \(\operatorname{ad}x:Y(T)\to Y(T')\) maps \(\theta\) to \(\theta'\) (via (5.2.3)). If \(F':T'\to T'\) is the Frobenius map of \(T'\), \(\operatorname{ad}x^{-1}(F')=\operatorname{ad}wF\) for some \(w\) in the Weyl group of \(T\). We have \({}^tF\theta=\theta\), and \({}^t(\operatorname{ad}wF)\theta=\theta\), hence \(\theta\) is fixed by \(w\), which by part (i) belongs to the Weyl group of \(T\) in \(S(T,\theta)\). Let \(T''=\operatorname{ad}y(T)\), \(y\in S(T,\theta)\), be an \(F\)-stable torus in \(S(T,\theta)\), with Frobenius \(F''\), such that \((\operatorname{ad}y)^{-1}(F'')=(\operatorname{ad}x)^{-1}(F')\). Replacing \((T,\theta)\) by \((T'',\operatorname{ad}y(\theta))\) (by applying 5.11) and replacing \(x\) by \(xy^{-1}\), we are reduced to the case where \(\operatorname{ad}xF=F'\operatorname{ad}x\). In this case, \((T,\theta)\) and \((T',\theta')\) are \(G^F\)-conjugate; the identity \(\operatorname{ad}x(Ft)=F'\operatorname{ad}x(t)\) for \(t\in T\) amounts to \(x^{-1}Fx\in T\); and replacing \(x\) by \(xt^{-1}\) with \(t\in T\), \(t^{-1}Ft=x^{-1}Fx\), makes \(x\) rational.\par

When roots are not all of the same length, one has to work in the dual group \(G^*\) of \(G\).\par

\noindent\textbf{DEFINITION 5.21.} A group \(G^*\) dual to \(G\) is a reductive group \(G^*\) defined over \(\mathbb F_q\), whose maximal torus \(\mathbf T^*\) is provided with an isomorphism with the dual of the maximal torus \(\mathbf T\) of \(G\), this isomorphism carrying simple roots to simple coroots.\par

Here are some properties of this duality. Let \(G\) and \(G^*\) be dual.\par

(5.21.1) \(G\) and \(G^*\) have the same Weyl group \(\mathbf W\). The action of \(\mathbf W\) on \(Y(\mathbf T^*)\) is the contragredient of its action on \(Y(\mathbf T)\).\par

(5.21.2) However, Frobenius does not act in the same way on \(\mathbf W\) in \(G\) and \(G^*\): it acts in inverse ways (the origin of this is that \(F\) on \(Y(\mathbf T^*)\) is the transpose of \(F\) on \(Y(\mathbf T)\)).\par

(5.21.3) The conjugacy classes of pairs \((T,B)\) (\(T\) an \(F\)-stable maximal torus, \(B\) a Borel subgroup containing it) in \(G\) and \(G^*\) correspond. With the notations of (1.13), (1.14), to the class of \((T,B)\) we associate the class of \((T^*,B^*)\) with \[
{}^t(\operatorname{ad}h(T,B)^{-1}\circ F)
=
\operatorname{ad}h(T^*,B^*)^{-1}\circ{}^tF;
\] the tori \(T\) and \(T^*\) are in duality.\par

(5.21.4) This induces a bijection between the rational conjugacy classes of maximal tori in \(G\) and \(G^*\). If \(T\) and \(T'\) are in corresponding classes, we have a natural class (mod \((N(T)/T)^F\)) of isomorphisms between \(T^*\) and \(T'\).\par
}

\EditionPage{33}{134}{%
(5.21.5) Let \(T\) be an \(F\)-stable torus on \(G\), and let \(\theta\) be a character of \(T^F\). Fix \(T'\), a corresponding torus in \(G^*\). The character \(\theta\) defines an element of \(T^{*F}\), then a \((N(T')/T')^F\)-conjugacy class of elements \(\theta'\) of \(T'\). In that way, we get a bijection between \(G^F\)-conjugacy classes of pairs \((T,\theta)\) as above, and \(G^{*F}\)-conjugacy classes of pairs \((T',\theta')\), \(T'\) an \(F\)-stable maximal torus of \(G^*\) and \(\theta'\) an element of \(T'^F\). This correspondence is compatible with extension of scalars from \(\mathbb F_q\) to \(\mathbb F_{q^n}\) (cf. 5.3.).\par

(5.21.6) By forgetting \(T'\), we see that each \(G^F\)-conjugacy class of pairs \((T,\theta)\) defines an element \(\theta'\in G^{*F}\), well-defined up to \(G^{*F}\)-conjugacy.\par

\noindent\textbf{PROPOSITION 5.22.} Two pairs \((T_1,\theta_1)\) and \((T_2,\theta_2)\) are geometrically conjugate if and only if \(\theta_1'\) and \(\theta_2'\) are geometrically conjugate.\par

An extension of scalars (cf. 5.4, 5.5) reduces us to the case where \(T_1\) and \(T_2\) are split. Conjugating, we may assume further that \(T_1=T_2\). In this case, geometric conjugacy is \(W\)-conjugacy, and the proposition is clear.\par

\noindent\textbf{PROPOSITION 5.23.} Let \(G\) and \(G^*\) be dual. Then, the centre of \(G\) is connected and smooth (resp. connected) if and only if the derived group of \(G^*\) is simply connected (resp. if its simply connected covering is unseparable).\par

Put \(Y=Y(T^*)\), \(T^*\) the maximal torus of \(G^*\); let \(Y_1\) be the subgroup generated by the coroots; and put \(Y_2=Y\cap Y_1\otimes\mathbb Q\). Then, \(Y_2\) is the \(Y\)-group of the maximal torus of the derived group, and \(Y_1\) that of the maximal torus of its universal covering. The group of this covering is the dual of the Pontrjagin dual of \(Y_2/Y_1\), and (5.23) follows by comparison with point (c) in the proof of (5.13).\par

\noindent\textbf{COROLLARY 5.24.} If the centre of \(G\) is connected, two pairs \((T_1,\theta_1)\), \((T_2,\theta_2)\) as in (5.22) are geometrically conjugate if and only if \(\theta_1'\) and \(\theta_2'\) are \(G^{*F}\)-conjugate.\par

Indeed, by a theorem of Steinberg, geometric conjugacy amounts in this case (5.23) to conjugacy, for semi-simple elements of \(G^*\) (their centralizers are connected).\par

\noindent\textbf{DEFINITION 5.25.} Let \((T',\theta')\) correspond to \((T,\theta)\) as in (5.21.5). The pair \((T,\theta)\) is maximally split if the \(\mathbb F_q\)-rank of \(T\) is equal to that of the centralizer of \(\theta'\), i.e., if \(T'\) is maximally split in that centralizer.\par

Any geometric conjugacy class of pairs \((T,\theta)\) contains maximally split pairs.\par

\noindent\textbf{PROPOSITION 5.26.} If the centre of \(G\) is connected, two geometrically conjugate maximally split pairs \((T_1,\theta_1)\), \((T_2,\theta_2)\) are \(G^F\)-conjugate.\par
}

\EditionPage{34}{135}{%
Via (5.21.5) and (5.24), this amounts to the known fact that two maximally split maximal tori in the centralizer \(Z(\theta')\) of a semi-simple element \(\theta'\) of \(G^{*F}\) are conjugate by an element of \(Z(\theta)^F\).\par

\noindent\textbf{PROPOSITION 5.27.} If \((T,\theta)\) is maximally split, and the image of \(T\) in the adjoint group is anisotropic, then \(\theta\) is non-singular.\par

Fix \((T',\theta')\) as in (5.21.5). By assumption, \(Z^0(\theta')\) contains no non-central \(F\)-stable split torus, hence \(Z^0(\theta')=T'\). The roots of \(T'\) in \(Z^0(\theta')\) correspond to the coroots of \(T\) orthogonal to \(\theta\), hence the proposition.\par

\noindent\textbf{COROLLARY 5.28.} If \((T,\theta)\) is maximally split, \(T\) is containd in an \(F\)-stable Levi subgroup \(L\) of an \(F\)-stable parabolic subgroup \(P\), and, in \(L\), \((T,\theta)\) is non-singular.\par

One chooses \(P\) and \(L\) so that the image of \(T\) in the adjoint group of \(L\) is anisotropic. Being maximally split in \(G\), \((T,\theta)\) is a fortiori maximally split in \(L\), and one applies (5.27).\par

The following result will be used in the next chapter.\par

\noindent\textbf{PROPOSITION 5.29.} Let \(T'\) be a subtorus of a torus \(T\), with \(T\) defined over \(\mathbb F_q\) (no assumption on \(T'\)). Let \(\theta\) be a character of \(T^F\). It is trivial on \(T'\cap T^F\) if and only if, when viewed as a character of \(Y(T)\), it is trivial on \[
((F-1)Y(T')\otimes\mathbb Q\cap Y(T)).
\]\par

The condition is that \(\theta(x)=0\) for \[
(F-1)^{-1}(x)\in(Y(T')\otimes\mathbb Q+Y(T)),
\] i.e., for \[
x\in((F-1)Y(T')\otimes\mathbb Q\cap Y(T))+((F-1)Y(T)).
\] The character \(\theta\) being trivial on \((F-1)Y(T)\), this means \(\theta(x)\) is trivial for \(x\in((F-1)Y(T')\otimes\mathbb Q\cap Y(T))\).\par

\EditionChapter{6. Intertwining numbers}

6.1. Let \(T,T'\) be two \(F\)-stable maximal tori in \(G\) and let \(\theta,\theta'\) be characters of \(T^F\) and \(T'^F\). We put \(N_G(T,T')=\{g\in G\mid Tg=gT'\}\) and \[
W_G(T,T')=T\backslash N_G(T,T')=N_G(T,T')/T'.
\] We will drop the index \(G\) if there is no ambiguity. \(F\) acts on \(W(T,T')\), and \[
W(T,T')^F=T^F\backslash N(T,T')^F=N(T,T')^F/T'^F.
\]\par

\noindent\textbf{THEOREM 6.2.} If \(\theta^{-1}\) is not geometrically conjugate to \(\theta'\), then \[
\left[
H_c^*(\widetilde X_{T\subset B})_\theta\otimes
H_c^*(\widetilde X_{T'\subset B'})_{\theta'}
\right]^{G^F}=0,
\]\par
}

\EditionPage{35}{136}{%
i.e., if an irreducible representation of \(G^F\) occurs in \(H_c^*(\widetilde X_{T\subset B})_\theta\), its dual does not occur in \(H_c^*(\widetilde X_{T'\subset B'})_{\theta'}\).\par

The virtual representation \(R_T^{\theta^{-1}}\) is the dual of \(R_T^\theta\) (as can be seen on the character formula 4.2), hence we have as a corollary:\par

\noindent\textbf{COROLLARY 6.3.} If \(\theta\) and \(\theta'\) are not geometrically conjugate, no irreducible representation of \(G^F\) can occur in both virtual representations \(R_T^\theta\) and \(R_{T'}^{\theta'}\).\par

The proof of 6.2 will make use of the following homotopy argument:\par

\noindent\textbf{PROPOSITION 6.4.} Let \(H\) be a connected algebraic group acting on a scheme \(Y\), separated and of finite type over \(k\). For any \(h\in H\), the action of \(h\) on \(H_c^*(Y,\mathbb Z/n)\) is trivial.\par

Let \(\pi\) be the projection of \(H\times Y\) onto \(H\) and let \(f:H\times Y\to H\times Y\) be defined by \(f(h,y)=(h,hy)\): \[
\begin{array}{ccccc}
&H\times Y&\xrightarrow{f}&H\times Y&\\
&\searrow_{\pi}&&\swarrow_{\pi}&\\{}[-1ex]
&&H&&
\end{array}
\] By the change-of-basis theorem in cohomology with compact support applied to \[
\begin{array}{ccc}
H\times Y&\longrightarrow&Y\\
\downarrow&&\downarrow\\
H&\longrightarrow&\operatorname{Spec}k,
\end{array}
\]  \(R^i\pi_!\mathbb Z/n\) is the constant sheaf \(H_c^i(Y,\mathbb Z/n)\) on \(H\). The automorphism \(f\) acts on it and, at \(h\in H\), it acts the way \(h\) acts on \(H_c^i(Y,\mathbb Z/n)\). At the identity element of \(H\), the action of \(f\) is trivial. An endomorphism of a constant sheaf over a connected base is constant, hence \(f\) acts trivially everywhere, and the proposition is proved.\par

\noindent\textbf{COROLLARY 6.5.} The conclusion of (6.4) holds in \(l\)-adic cohomology.\par

The proof is a passage to limit.\par

6.6. Proof of 6.2. On \(\widetilde X_{T\subset B}\times\widetilde X_{T'\subset B'}\) we have commuting actions of \(G^F\), \(T^F\), and \(T'^F\) (\(G^F\) acts diagonally). By Künneth's Theorem, we have to prove that \[
\left[
H_c^*(\widetilde X_{T\subset B}\times\widetilde X_{T'\subset B'})_{\theta,\theta'}
\right]^{G^F}=0,
\] where the symbol \(M_{\theta,\theta'}\), applied to any \(T^F\times T'^F\)-module \(M\), means the subspace\par
}

\EditionPage{36}{137}{%
of \(M\), where \(T^F\) and \(T'^F\) act by \(\theta\) and \(\theta'\).\par

\[
\left\{
\begin{aligned}
&\text{Let }U\text{ (resp. }U'\text{) be the unipotent radical of }B\text{ (resp. }B'\text{).}\\
&\text{Let }S_{T\subset B}=\{g\in G\mid g^{-1}Fg\in FU\},\quad
S_{T'\subset B'}=\{g'\in G\mid g'^{-1}Fg'\in FU'\}.
\end{aligned}
\right.
\] The unipotent group \((U\cap FU)\times(U'\cap FU')\) acts freely on \(S_{T\subset B}\times S_{T'\subset B'}\) and the orbit space is \(\widetilde X_{T\subset B}\times\widetilde X_{T'\subset B'}\) (see (1.17)). It follows that \[
H_c^i(\widetilde X_{T\subset B}\times\widetilde X_{T'\subset B'})(-d)
\simeq
H_c^{i+2d}(S_{T\subset B}\times S_{T'\subset B'}),
\] where \(d\) is the dimension of \((U\cap FU)\times(U'\cap FU')\). This isomorphism is compatible with the action of \(G^F\times T^F\times T'^F\); moreover this group acts trivially on \(\overline{\mathbb Q}_l(-d)\). Hence it is sufficient to prove that \[
H_c^*(S_{T\subset B}\times S_{T'\subset B'})_{\theta,\theta'}^{G^F}
=
H_c^*((S_{T\subset B}\times S_{T'\subset B'})/G^F)_{\theta,\theta'}
=0.
\] The map \[
(g,g')\longmapsto(x,x',y),\qquad
x=g^{-1}Fg,\quad x'=g'^{-1}Fg',\quad y=g^{-1}g',
\] defines an isomorphism of \((S_{T\subset B}\times S_{T'\subset B'})/G^F\) with \[
\bar S=\{(x,x',y)\in FU\times FU'\times G\mid xF(y)=yx'\}.
\] Under this isomorphism, \(T^F\times T'^F\) acts on \(\bar S\) by the formula \[
\tag{6.6.1}
(x,x',y)\longmapsto(t^{-1}xt,t'^{-1}x't',t^{-1}yt'),
\qquad(t,t')\in T^F\times T'^F.
\] Let \(U'^{-}\) be opposed to \(U'\) with respect to \(T'\). Any \(g\in G\) can be written uniquely in the form \(g=u_gn_gu'_g\), with with \(u_g\in U\cap n_gU'^{-}n_g^{-1}\), \(n_g\in N(T,T')\) (see (6.1)), \(u'_g\in U'\); this follows easily from Bruhat's Lemma. For any \(w\in W(T,T')\), let \(G_w\) be the set of all \(g\in G\) such that \(n_g\) represents \(w\). Then \((G_w)_{w\in W(T,T')}\) is a finite partition of \(G\) into locally closed subschemes. It has the property (6.6.2) below, expressing the fact that the closure of a Bruhat cell is a union of Bruhat cells.\par

(6.6.2.) For a suitable ordering of \(W(T,T')\), the unions \[
\Sigma_w=\bigcup_{w'<w}G_{w'}
\] are closed.\par

Let \(\bar S_w=\{(x,x',y)\in\bar S\mid y\in G_w\}\). Then \((\bar S_w)_{w\in W(T,T')}\) is a finite partition of \(\bar S\) into locally closed subschemes, stable under \(T^F\times T'^F\). It inherits a property analogous to (6.6.2), that is, the unions \(\bigcup_{w'<w}\bar S_{w'}\) are closed for any \(w\).\par

The spectral sequence associated to the filtration of \(\bar S\) by these unions shows that, in order to prove that \(H_c^*(\bar S)_{\theta,\theta'}=0\), it is sufficient to prove that \(H_c^*(\bar S_w)_{\theta,\theta'}=0\) for any \(w\in W(T,T')\).\par

Let \[
H_w=\{(t,t')\in T\times T'\mid t'F(t')^{-1}=F(\dot w)^{-1}tF(t)^{-1}F(\dot w)\},
\] where \(\dot w\in N(T,T')\) represents \(w\). Then \(H_w\) is a closed subgroup of \(T\times T'\), containing \(T^F\times T'^F\). We define an action of \(H_w\) on \(\bar S_w\) as follows. Let \((t,t')\in H_w\); for \((x,x',y)\in\bar S_w\) define\par
}

\EditionPage{37}{138}{%
\[
f_{t,t'}(x,x',y)=(\widetilde x,\widetilde x',\widetilde y),
\] where \[
\begin{aligned}
\widetilde x&=t^{-1}xF(u_y)tF(t^{-1}u_y^{-1}t),\\
\widetilde x'&=t'^{-1}x'F(u'_y)^{-1}t'F(t'^{-1}u'_yt'),\\
\widetilde y&=t^{-1}yt'.
\end{aligned}
\] It is easy to check that \((\widetilde x,\widetilde x',\widetilde y)\in\bar S_w\), hence \(f_{t,t'}:\bar S_w\longrightarrow\bar S_w\). It is also easy to check that \[
f_{t_1,t'_1}\circ f_{t_2,t'_2}=f_{t_1t_2,t'_1t'_2}
\] for \((t_i,t'_i)\in H_w\), \(i=1,2\), hence \(f_{t,t'}\) defines an action of \(H_w\) on \(\bar S_w\). It is clear that this action extends the action (6.6.1) of \(T^F\times T'^F\) on \(\bar S_w\). The theorem now follows from 6.5 and the\par

\noindent\textbf{LEMMA 6.7.} If the character \(\theta\theta'\) of \(T^F\times T'^F\) is trivial on \(H_w^0\cap(T^F\times T'^F)\), then \(\theta'^{-1}=\theta\circ\operatorname{ad}(F(w))\) (as characters of \(Y(T')\)).\par

The subgroup \(H_w\) of \(T\times T'\) is the kernel of the composite map \[
T\times T'\xrightarrow{x^{-1}Fx}T\times T'\xrightarrow{t^{-1}\operatorname{ad}F(w)(t')}T.
\] Applying the functor \(Y\), we get that \(Y(H_w^0)\) is the kernel \(K\) of \[
Y(T)\times Y(T')\xrightarrow{Fx-x}Y(T)\times Y(T')\xrightarrow{\operatorname{ad}F(w)(x')-x}Y(T).
\] By (5.29), the triviality of \(\theta\theta'\) on \(H_w^0\cap(T^F\times T'^F)\) means that, when we view \(\theta\theta'\) as a character of \(Y(T)\times Y(T')\), it is trivial on \[
(F-1)(K\otimes\mathbb Q)\cap(Y(T)\times Y(T')).
\] Since the map \(F-1\) is injective, this intersection is \[
\operatorname{Ker}(\operatorname{ad}F(w)(x')-x:Y(T)\times Y(T')\longrightarrow Y(T))
\] and the assumption becomes the triviality of \(\theta\theta'\) on \(\operatorname{Ker}(\operatorname{ad}F(w)(x')-x)\), i.e., the identity \[
\theta\circ\operatorname{ad}F(w)(\theta')=1.
\]\par

We will now conjugate (6.2) with the character formula (4.2) to get quantitative results.\par

\noindent\textbf{THEOREM 6.8.} Let \(\theta\in(T^F)^\vee\), \(\theta'\in(T'^F)^\vee\). Then \[
\langle R_T^\theta,R_{T'}^{\theta'}\rangle_{G^F}
=
\#\{w\in W(T,T')^F\mid \operatorname{ad}w(\theta')=\theta\}.
\]\par

\noindent\textbf{THEOREM 6.9.} \[
\tag{6.9.1}
\frac1{|G^F|}\sum_{\substack{u\in G^F\\\text{unipotent}}}Q_{T,G}(u)Q_{T',G}(u)
=
\frac{|N(T,T')^F|}{|T^F|\,|T'^F|}.
\]\par

We will prove (6.8) and (6.9) simultaneously, by induction on the dimension of \(G\).\par

\noindent\textbf{LEMMA 6.10.} If (6.9) holds, for \(G\) replaced by \(Z^0(s)\), where \(s\in G^F\) is any semisimple element not contained in the centre \(Z\) of \(G\), then \[
\tag{6.10.1}
\langle R_T^\theta,R_{T'}^{\theta'}\rangle_{G^F}
=
\#\{w\in W(T,T')^F\mid \operatorname{ad}w(\theta')=\theta\}+\alpha
\]\par
}

\EditionPage{38}{139}{%
where \[
\alpha=
\sum_{s\in T^F\cap Z}\theta(s)\theta'(s^{-1})
\left(
\frac1{|G^F|}\sum_{\substack{u\in G^F\\\text{unipotent}}}Q_{T,G}(u)Q_{T',G}(u)
-
\frac{|N(T,T')^F|}{|T^F|\,|T'^F|}
\right).
\]\par

In particular, (6.9) implies (6.8).\par

According to (4.2) we have: \[
\begin{aligned}
\langle R_T^\theta,R_{T'}^{\theta'}\rangle
&=\frac1{|G^F|}\sum_{g_1\in G^F}\operatorname{Tr}(g_1,R_T^\theta)\operatorname{Tr}(g_1^{-1},R_{T'}^{\theta'})\\
&=\frac1{|G^F|}\sum_{\substack{s\in G^F\\\text{semis.}}}\frac1{|Z^0(s)^F|^2}
\sum_{\substack{g,g'\in G^F\\g^{-1}sg\in T^F\\g'^{-1}sg'\in T'^F}}
\theta(g^{-1}sg)\theta'(g'^{-1}sg')^{-1}\\
&\qquad\times
\sum_{\substack{u\in Z^0(s)^F\\\text{unipotent}}}
Q_{gTg^{-1},Z^0(s)}(u)Q_{g'T'g'^{-1},Z^0(s)}(u).
\end{aligned}
\] With our assumption, this can be written as \[
\begin{aligned}
\alpha+\frac1{|G^F|}\sum_{\substack{s\in G^F\\\text{semis.}}}\frac1{|Z^0(s)^F|}
&\sum_{\substack{g,g'\in G^F\\g^{-1}sg\in T^F\\g'^{-1}sg'\in T'^F}}
\theta(g^{-1}sg)\theta'(g'^{-1}sg')^{-1}\\
&\times
\frac{|N_{Z^0(s)}(gTg^{-1},g'T'g'^{-1})^F|}{|T^F|\,|T'^F|}.
\end{aligned}
\] The formula \((g,g',n_1)\longmapsto(g,n,n_1)\), \(n=g'^{-1}n_1g\), establishes a one-to-one correspondence between the sets \[
\begin{aligned}
\{(g,g',n_1)\in G^F\times G^F\times G^F\mid
&\ g^{-1}sg\in T^F,\ g'^{-1}sg'\in T'^F,\\
&\ n_1\in N_{Z^0(s)}(gTg^{-1},g'T'g'^{-1})^F\},
\end{aligned}
\]  \[
\{(g,n,n_1)\in G^F\times N_G(T,T')^F\times Z^0(s)^F\mid g^{-1}sg\in T^F\}.
\] It follows that \[
\begin{aligned}
\langle R_T^\theta,R_{T'}^{\theta'}\rangle
&=\alpha+\frac1{|G^F|}\sum_{\substack{s\in G^F\\\text{semis.}}}\frac1{|Z^0(s)^F|}
\sum_{\substack{g\in G^F\\n\in N_G(T,T')^F\\g^{-1}sg\in T^F}}
\theta(g^{-1}sg)\theta'(ng^{-1}sgn^{-1})^{-1}
\frac{|Z^0(s)^F|}{|T^F|\,|T'^F|}\\
&=\alpha+\frac1{|G^F|}\sum_{\substack{t\in T^F\\n\in N_G(T,T')^F}}
\theta(t)\theta'(ntn^{-1})^{-1}
\frac{|G^F|}{|T^F|\,|T'^F|}\\
&=\alpha+\frac1{|T'^F|}\sum_{n\in N_G(T,T')^F}\frac1{|T^F|}
\sum_{t\in T^F}\theta(t)\theta'(ntn^{-1})^{-1}\\
&=\alpha+\#\{w\in W(T,T')^F\mid \operatorname{ad}w(\theta')=\theta\},
\end{aligned}
\] and (6.10.1) is proved.\par

Proof of 6.9. Let \(\bar G=G/Z\); the centre of \(\bar G\) consists of the identity element. Both sides of (6.9.1) remain unchanged when \(G\) is replaced by \(\bar G\); for\par
}

\EditionPage{39}{140}{%
the left hand side this follows from (4.1.1), while for the right hand side this is easy to check. Hence in order to prove (6.9), we may assume that the centre of \(G\) has a single element; moreover we may assume by induction that the conclusion of (6.9) is true for \(G\) replaced by \(Z^0(s)\), where \(s\in G^F\) is an arbitrary non-central semisimple element, and for any two \(F\)-stable maximal tori in \(Z^0(s)\). (To start the induction we may take \(G\) to be the group with only one element in which case (6.9) is clear.) It follows from (6.3) that for any non-trivial character \(\theta\) of \(T^F\) one has \[
\langle R_T^\theta,R_{T'}^1\rangle=0.
\] Substituting this information into (6.10), we see that if \(T^F\) has any non-trivial character, then \(\alpha=0\) which proves (6.9). A similar argument applies when \(|T'^F|\ne1\). It remains to check the special case where \(|T^F|=|T'^F|=1\). In this case, we must have \(q=2\), and \(T,T'\) must be \(\mathbb F_q\)-split tori. In particular, \(T\) and \(T'\) are conjugate under \(G^F\), hence \[
\#\{w\in W(T,T')^F\mid \operatorname{ad}w(\theta')=\theta\}
=|W(T)^F|=|W(T)|.
\] Moreover \(R_T^1\) and \(R_{T'}^1\) are both equal to the representation of \(G^F\) induced by the unit representation of \(B^F\), where \(B\) is an \(F\)-stable Borel subgroup of \(G\). It follows that \[
\tag{1.10}
\langle R_T^1,R_{T'}^1\rangle=|B^F\backslash G^F/B^F|=|W(T)|.
\] Substituting in (6.10) we get (6.9) in this case. This completes the proof of (6.9) (and of (6.8) by (6.10)).\par

\EditionChapter{7. Computations on semisimple elements}

The following result describes the Euler characteristic of the scheme \(X_{T\subset B}\). Let \(\sigma(G)\) (resp. \(\sigma(T)\)) be the \(\mathbb F_q\)-rank of \(G\) (resp. \(T\)).\par

\noindent\textbf{THEOREM 7.1.} \[
\chi(X_{T\subset B})=Q_{T,G}(e)
=(-1)^{\sigma(G)-\sigma(T)}\frac{|G^F|}{\operatorname{St}_G(e)|T^F|}.
\] (Here \(\operatorname{St}_G\) is the Steinberg representation of \(G^F\), identified with its character.)\par

We may assume that the centre of \(G\) is reduced to a single element and that the statement is true when \(G\) is replaced by \(Z^0(s)\), where \(s\in G^F\) is any non-central semisimple element.\par

We first assume that \(T^F\) has a non-trivial character \(\theta\). The Steinberg representation occurs in \[
R_{T^*}^1=\operatorname{Ind}_{B^{*F}}^{G^F}(1)
\] (\(T^*\subset B^*\) as in 1.8), hence, by (6.3), it does not occur in \(R_T^\theta\): \(\langle R_T^\theta,\operatorname{St}_G\rangle=0\). It is known that \[
\operatorname{St}_G(g)=
\begin{cases}
(-1)^{\sigma(G)-\sigma(Z^0(g))}\operatorname{St}_{Z^0(g)}(e),&g\in G^F\text{ semisimple},\\
0,&g\in G^F\text{ non-semisimple}.
\end{cases}
\] If we use 4.2, it follows:\par
}

\EditionPage{40}{141}{%
\[
\sum_{s\in G^F}(-1)^{\sigma(G)-\sigma(Z^0(s))}
\frac{\operatorname{St}_{Z^0(s)}(e)}{|Z^0(s)^F|}
\sum_{\substack{g\in G^F\\g^{-1}sg\in T}}
Q_{gTg^{-1},Z^0(s)}(e)\theta(g^{-1}sg)=0
\] (sum over the semisimple elements \(s\) of \(G^F\)).\par

By our assumption, we may substitute \[
Q_{gTg^{-1},Z^0(s)}(e)
=(-1)^{\sigma(Z^0(s))-\sigma(T)}
\frac{|Z^0(s)^F|}{\operatorname{St}_{Z^0(s)}(e)|T^F|},
\] for all \(s\ne e\): \[
\frac{(-1)^{\sigma(G)-\sigma(T)}}{|T^F|}
\sum_{\substack{s\in G^F\\s\ne e}}
\sum_{\substack{g\in G^F\\g^{-1}sg\in T^F}}
\theta(g^{-1}sg)
+\operatorname{St}_G(e)Q_{T,G}(e)=0,
\]  \[
(-1)^{\sigma(G)-\sigma(T)}\frac{|G^F|}{|T^F|}
\sum_{\substack{t\in T^F\\t\ne e}}\theta(t)
+\operatorname{St}_G(e)Q_{T,G}(e)=0.
\] Since the character \(\theta\) in non-trivial, \(\sum_{t\in T^F}\theta(t)=0\) and the desired formula for \(Q_{T,G}(e)\) follows.\par

In the case where \(|T^F|=1\), \(T\) must be an \(\mathbb F_q\)-split torus and \(q=2\). In this case, \(Q_{T,G}(e)=|G^F|/|B^{*F}|\), where \(B^*\) is any \(F\)-stable Borel subgroup (1.8). This agrees with the statement of the theorem and ends the proof.\par

\noindent\textbf{COROLLARY 7.2.} For any semisimple element \(s\in G^F\), \[
\operatorname{Tr}(s,R_T^\theta)
=(-1)^{\sigma(Z^0(s))-\sigma(T)}
\frac1{\operatorname{St}_{Z^0(s)}(e)|T^F|}
\sum_{\substack{g\in G^F\\g^{-1}sg\in T^F}}\theta(g^{-1}sg).
\] This follows from (7.1) and 4.2. An equivalent statement is:\par

\noindent\textbf{PROPOSITION 7.3.} \[
(-1)^{\sigma(G)-\sigma(T)}R_T^\theta\otimes\operatorname{St}_G
=\operatorname{Ind}_{T^F}^{G^F}(\theta).
\]\par

\noindent\textbf{PROPOSITION 7.4.} If \(\theta\) is a non-singular character of \(T^F\) (5.19) then \((-1)^{\sigma(G)-\sigma(T)}R_T^\theta\) can be represented by an actual \(G^F\)-module. If \(\theta\) is in general position (5.15), \((-1)^{\sigma(G)-\sigma(T)}R_T^\theta\) is irreducible.\par

By embedding \(G\) in a group \(G_1\) with connected centre and the same derived group, as in (5.18), we reduce to the case where \(\theta\) is in general position. It remains to use (6.8) and (7.1).\par

\noindent\textbf{PROPOSITION 7.5.} Let \(s\in G^F\) be a semisimple element. The character of \[
\tag{7.5.1}
\frac1{\operatorname{St}_G(s)}
\sum_{\substack{T\\s\in T}}
\sum_{\theta\in(T^F)^\vee}
\theta(s)^{-1}(-1)^{\sigma(G)-\sigma(T)}R_T^\theta
\] equals \(|Z(s)^F|\) on elements conjugate to \(s\) in \(G^F\) and zero on all other elements of \(G^F\).\par

Let \(\mu\) be the character of (7.5.1) and let \(\mu'\) be the class function on \(G^F\) defined by\par
}

\EditionPage{41}{142}{%
\[
\mu'(g)=
\begin{cases}
|Z(s)^F|,&g\text{ conjugate to }s\text{ in }G^F,\\
0,&\text{otherwise}.
\end{cases}
\] In order to prove that \(\mu=\mu'\) it is sufficient to prove that \[
\langle\mu,\mu\rangle=\langle\mu,\mu'\rangle=\langle\mu',\mu'\rangle
\qquad\text{(see (0.3)).}
\] We have \(\langle\mu',\mu'\rangle=|Z(s)^F|\); from (6.8) we see that \[
\begin{aligned}
\langle\mu,\mu\rangle
&=\frac1{\operatorname{St}_G(s)^2}
\sum_{\substack{T\ni s\\T'\ni s}}
\sum_{\substack{\theta\in(T^F)^\vee\\\theta'\in(T'^F)^\vee}}
\theta(s^{-1})\theta'(s)
\#\{n\in N(T,T')^F\mid \operatorname{ad}n(\theta')=\theta\}|T^F|^{-1}\\
&=\frac1{\operatorname{St}_G(s)^2}
\sum_{\substack{T\ni s\\\theta\in(T^F)^\vee\\n\in G^F\\nsn^{-1}\in T}}
\theta(s^{-1})\theta(nsn^{-1})|T^F|^{-1}\\
&=\frac1{\operatorname{St}_G(s)^2}
\#\{T\ni s,\ n\in G^F\mid sn=ns\}\\
&=\frac{|Z(s)^F|}{\operatorname{St}_G(s)^2}
\#\{F\text{-stable maximal tori in }Z^0(s)\}
=|Z(s)^F|
\end{aligned}
\] (by [15, Cor. 14.16]).\par

By (7.2), we have \[
\begin{aligned}
\langle\mu,\mu'\rangle=\mu(s)
&=\frac1{\operatorname{St}_G(s)}
\sum_{\substack{T\\s\in T}}
\sum_{\theta\in(T^F)^\vee}
\theta(s^{-1})
\frac{(-1)^{\sigma(G)-\sigma(Z^0(s))}}{\operatorname{St}_{Z^0(s)}(e)|T^F|}
\sum_{\substack{g\in G^F\\g^{-1}sg\in T^F}}\theta(g^{-1}sg)\\
&=\frac1{\operatorname{St}_G(s)^2}
\sum_{\substack{T\\s\in T}}
\sum_{\substack{g\in G^F\\g^{-1}sg\in T^F}}
|T^F|^{-1}\theta(s^{-1})\theta(g^{-1}sg)\\
&=\frac1{\operatorname{St}_G(s)^2}
\#\{T\ni s,\ g\in G^F\mid sg=gs\}\\
&=\frac{|Z(s)^F|}{\operatorname{St}_G(s)^2}
\#\{F\text{-stable maximal tori in }Z^0(s)\}
=|Z(s)^F|,
\end{aligned}
\] and the proposition is proved.\par

\noindent\textbf{COROLLARY 7.6.} For any \(\rho\in\mathcal R(G^F)\) and any semisimple element \(s\in G^F\), \[
\tag{7.6.1}
\operatorname{Tr}(s,\rho)
=\frac1{\operatorname{St}_G(s)}
\sum_{\substack{T\\s\in T}}
\sum_{\theta\in(T^F)^\vee}
\theta(s)(-1)^{\sigma(G)-\sigma(T)}\langle\rho,R_T^\theta\rangle.
\] In particular, if \(s\in G^F\) and is regular semisimple, contained in a unique torus \(T\), \[
\tag{7.6.2}
\operatorname{Tr}(s,\rho)=
\sum_{\theta\in(T^F)^\vee}\theta(s)\langle\rho,R_T^\theta\rangle,
\]  \[
\tag{7.6.3}
\dim\rho=
\frac1{\operatorname{St}_G(e)}
\sum_T\sum_{\theta\in(T^F)^\vee}
(-1)^{\sigma(G)-\sigma(T)}\langle\rho,R_T^\theta\rangle.
\]\par

\noindent\textbf{COROLLARY 7.7.} For any irreducible representation \(\rho\) of \(G^F\) there exists an \(F\)-stable maximal torus \(T\) and a character \(\theta\) of \(T^F\) such that \(\langle\rho,R_T^\theta\rangle\ne0\).\par

This follows from (7.6.3).\par
}

\EditionPage{42}{143}{%
\noindent\textbf{DEFINITION 7.8.} An irreducible representation \(\rho\) of \(G^F\) is said to be unipotent if \(\langle\rho,R_T^1\rangle\ne0\) for some \(F\)-stable maximal torus \(T\).\par

By (6.3) and (7.7), \(\rho\) is unipotent if and only if \(\langle\rho,R_T^\theta\rangle=0\) for any \(F\)-stable maximal torus \(T\) and any \(\theta\in(T^F)^\vee\), \(\theta\ne1\). For example, any irreducible component of \(\operatorname{Ind}_{B^{*F}}^{G^F}(1)\) (\(B^{*F}\) an \(F\)-stable Borel subgroup) is a unipotent representation. For unipotent representations, (7.6) becomes:\par

\noindent\textbf{PROPOSITION 7.9.} Let \(\rho\) be a unipotent representation of \(G^F\) and let \(s\in G^F\) be a semisimple element. Then \[
\operatorname{Tr}(s,\rho)=
\frac1{\operatorname{St}_G(s)}
\sum_{\substack{T\\s\in T}}
(-1)^{\sigma(G)-\sigma(T)}\langle\rho,R_T^1\rangle.
\] In particular, if \(s\in G^F\) is regular semisimple contained in a unique maximal torus \(T\), \(\operatorname{Tr}(s,\rho)=\langle\rho,R_T^1\rangle\) and \[
\dim\rho=
\frac1{\operatorname{St}_G(e)}
\sum_T(-1)^{\sigma(G)-\sigma(T)}\langle\rho,R_T^1\rangle.
\]\par

It follows that if \(B\) is an \(F\)-stable Borel subgroup of \(G\) and \(T\) is an \(F\)-stable maximal torus in \(B\), we have \(\operatorname{Tr}(s,\rho)=\langle\rho,\operatorname{Ind}_{B^F}^{G^F}(1)\rangle\); in the case where \(\langle\rho,\operatorname{Ind}_{B^F}^{G^F}(1)\rangle\ne0\), this is a result of Curtis, Kilmoyer and Seitz (see C. W. Curtis, On the Values of Certain Irreducible Characters of Finite Chevalley Groups, Ist. Naz. di Alta Mat. Symp. Mat. XIII, 1974, 343--355).\par

\noindent\textbf{PROPOSITION 7.10.} Let \(G^{\mathrm{ad}}\) be the adjoint group of \(G\). Then, the restriction to \(G^F\) of a unipotent representation of \((G^{\mathrm{ad}})^F\) is irreducible; non-isomorphic unipotent representations have non-isomorphic restrictions; and every unipotent representation of \(G\) is such a restriction.\par

It suffices to check that, if \(\sigma,\tau\) are unipotent representations of \((G^{\mathrm{ad}})^F\), then \(\langle\sigma,\tau\rangle_{G^{\mathrm{ad}}F}=\langle\sigma,\tau\rangle_{G^F}\). Let \(G_1\) be the image of \(G^F\) in \((G^{\mathrm{ad}})^F\). One has \[
\begin{aligned}
\langle\sigma,\tau\rangle_{G^F}
&=\langle\sigma,\tau\rangle_{G_1}
=\left\langle\operatorname{Ind}_{G_1}^{G^{\mathrm{ad}}F}(\operatorname{Res}\sigma),\tau\right\rangle_{G^{\mathrm{ad}}F}\\
&=\sum_{\theta\in(G^{\mathrm{ad}}F/G_1^F)^\vee}
\langle\sigma\otimes\theta,\tau\rangle_{G^{\mathrm{ad}}F}.
\end{aligned}
\] By the exclusion theorem 6.3 and (1.23), (1.27), \(\langle\sigma\otimes\theta,\tau\rangle_{G^{\mathrm{ad}}F}=0\) for \(\theta\ne1\), hence the proposition.\par

\noindent\textbf{PROPOSITION 7.11.} Let \(\rho\) be a virtual representation of \(G^F\) such that \(\operatorname{Tr}(su,\rho)=\operatorname{Tr}(s,\rho)\) for any \(s\in G^F\) semisimple, \(u\in G^F\) unipotent, \(su=us\). Let \(T\) be an \(F\)-stable maximal torus and let \(\theta\) be a character of \(T\). Then \[
\langle\rho,R_T^\theta\rangle_{G^F}
=(-1)^{\sigma(G)-\sigma(T)}\langle\rho\otimes\operatorname{St}_G,R_T^\theta\rangle_{G^F}
=\langle\rho,\theta\rangle_{T^F}.
\]\par

The Steinberg character is integral valued, hence by (7.3),\par
}

\EditionPage{43}{144}{%
\[
(-1)^{\sigma(G)-\sigma(T)}
\langle\rho\otimes\operatorname{St}_G,R_T^\theta\rangle_{G^F}
=
\langle\rho,(-1)^{\sigma(G)-\sigma(T)}\operatorname{St}_G\otimes R_T^\theta\rangle_{G^F}
=
\langle\rho,\operatorname{Ind}_{T^F}^{G^F}(\theta)\rangle_{G^F}
=
\langle\rho,\theta\rangle_{T^F}.
\]

We now prove the identity \[
\tag{7.11.1}
\langle\rho,R_T^\theta\rangle_{G^F}=\langle\rho,\theta\rangle_{T^F}.
\] The special case of this identity \[
\tag{7.11.2}
\langle 1,R_T^1\rangle_{G^F}=1
\] follows from the fact that the Euler characteristic of \(X_{T\subset B}/G^F\) (which is the same as the Euler characteristic of \[
\{g\in G\mid g^{-1}Fg\in FU\}/G^F\times T^F=FU/T^F
\] by (1.17)) equals 1. If we use (4.2), (7.11.2) can be written in the form \[
\tag{7.11.3}
\frac1{|G^F|}
\sum_{\substack{s\in G^F\\\mathrm{semis.}}}
\frac1{|Z^0(s)^F|}
\sum_{\substack{g\in G^F\\g^{-1}sg\in T^F}}
\sum_{\substack{u\in Z^0(s)^F\\\mathrm{unipotent}}}
Q_{gTg^{-1},Z^0(s)}(u)=1.
\] This implies, by induction on the dimension of \(G\), that \[
\tag{7.11.4}
\frac1{|G^F|}
\sum_{\substack{u\in G^F\\\mathrm{unipotent}}}Q_{T,G}(u)
=
\frac1{|T^F|}.
\] Indeed, by substituting (7.11.4) in (7.11.3) we find a true identity: \[
\frac1{|G^F|}
\sum_{\substack{s\in G^F\\\mathrm{semis.}}}
\frac1{|T^F|}
\#\{g\in G^F\mid g^{-1}sg\in T^F\}=1.
\] We apply (4.2) again and use (7.11.4) to compute: \[
\begin{aligned}
\langle\rho,R_T^\theta\rangle_{G^F}
&=\frac1{|G^F|}
\sum_{\substack{s\in G^F\\\mathrm{semis.}}}
\frac1{|Z^0(s)^F|}
\sum_{\substack{g\in G^F\\g^{-1}sg\in T^F}}
\sum_{\substack{u\in Z^0(s)^F\\\mathrm{unipotent}}}
Q_{gTg^{-1},Z^0(s)}(u)\theta(g^{-1}sg)^{-1}\operatorname{Tr}(s,\rho)\\
&=\frac1{|G^F|}
\sum_{\substack{s\in G^F\\\mathrm{semis.}}}
\sum_{\substack{g\in G^F\\g^{-1}sg\in T^F}}
\frac1{|T^F|}\theta(g^{-1}sg)^{-1}\operatorname{Tr}(s,\rho)\\
&=\frac1{|G^F|}
\sum_{t\in T^F}\theta(t)^{-1}\operatorname{Tr}(t,\rho)
=\langle\rho,\theta\rangle_{T^F}.
\end{aligned}
\]\par

\noindent\textbf{PROPOSITION 7.12.} Let \(\rho\) be as in (7.11). Then \[
\tag{7.12.1}
\rho=
\sum_{(T)}\frac1{|W(T)^F|}
\sum_{\theta\in(T^F)^\vee}
\langle\rho,\theta\rangle_{T^F}R_T^\theta,
\]  \[
\tag{7.12.2}
\rho\otimes\operatorname{St}_G=
\sum_{(T)}\frac{(-1)^{\sigma(G)-\sigma(T)}}{|W(T)^F|}
\sum_{\theta\in(T^F)^\vee}
\langle\rho,\theta\rangle_{T^F}R_T^\theta,
\] where \(\sum_{(T)}\) means sum over all \(G^F\)-conjugacy classes of \(F\)-stable maximal tori \(T\).\par

The character of \(\rho\otimes\operatorname{St}_G\) is a linear combination of the form \(\sum_{(T,\theta)}c_{T,\theta}R_T^\theta\) (sum over all \(G^F\)-conjugacy classes of pairs \((T,\theta)\)). The coefficients \(c_{T,\theta}\) are\par
}

\EditionPage{44}{145}{%
determined by (7.11) and (7.12.2) follows.\par

Now let \(\rho'\) (resp. \(\rho''\)) be the right hand side of (7.12.1) (resp. (7.12.2)). We have clearly \[
\tag{7.12.3}
\langle\rho',\rho'\rangle=\langle\rho'',\rho''\rangle
\] and by (7.11), \[
\tag{7.12.4}
\langle\rho,\rho'\rangle=\langle\rho\otimes\operatorname{St}_G,\rho''\rangle.
\] We note also that \[
\tag{7.12.5}
\langle\rho,\rho\rangle
=
\langle\rho\otimes\operatorname{St}_G,\rho\otimes\operatorname{St}_G\rangle.
\] This follows from the fact that the number of unipotents in \(Z^0(s)^F\) equals \(\operatorname{St}_G(s)^2\) for any semisimple element \(s\in G^F\) (see [15, Thm. 15.1]). By (7.12.2) we have \[
\langle\rho\otimes\operatorname{St}_G-\rho'',\rho\otimes\operatorname{St}_G-\rho''\rangle=0.
\] This, together with (7.12.3), (7.12.4) and (7.12.5) implies: \[
\langle\rho-\rho',\rho-\rho'\rangle=0,
\] hence \[
\rho=\rho'.
\]\par

Remark 7.13. It is known that \[
\langle\rho,\rho\rangle_{G^F}
=
\sum_{(T)}\frac1{|W(T)^F|}\langle\rho,\rho\rangle_{T^F}
\] (N. Kawanaka, A theorem on finite Chevalley groups, Osaka J. Math. 10 (1973), 1--13); this could be also deduced from (7.12.1) and (6.8).\par

\noindent\textbf{COROLLARY 7.14.} \[
\tag{7.14.1}
1=\sum_{(T)}\frac1{|W(T)^F|}R_T^1.
\]  \[
\tag{7.14.2}
\operatorname{St}_G
=\sum_{(T)}\frac{(-1)^{\sigma(G)-\sigma(T)}}{|W(T)^F|}R_T^1.
\] In particular, \[
\langle\operatorname{St}_G,R_T^1\rangle=(-1)^{\sigma(G)-\sigma(T)}.
\]\par

Here is an alternative proof of (7.14.1). In the language of (1.4), (7.14.1) asserts that \[
\sum_{w\in W}R^1(w)=|W|\,1.
\] Since \((X(w))_{w\in W}\) is a partition of the flag manifold into locally closed subschemes, we have\par
}

\EditionPage{45}{146}{%
\[
\sum_{w\in W}R^1(w)=\sum_i(-1)^iH_c^i(X_G)
\] as virtual representations of \(G^F\). The action of \(G^F\) on \(X_G\) is the restriction of the action of the connected group \(G\); by (6.5), \(G^F\) acts trivially on \(H_c^*(X_G)\) and it remains to use the fact that the Euler characteristic of \(X_G\) equals \(|W|\).\par

\noindent\textbf{COROLLARY 7.15.} \[
\tag{7.15.1}
\operatorname{St}_G
=
\sum_{(T)}\frac{(-1)^{\sigma(G)-\sigma(T)}}{|W(T)^F|}
\operatorname{Ind}_{T^F}^{G^F}(1),
\]  \[
\tag{7.15.2}
\operatorname{St}_G\otimes\operatorname{St}_G
=
\sum_{(T)}\frac1{|W(T)^F|}
\operatorname{Ind}_{T^F}^{G^F}(1).
\] This follows from (7.1) by tensoring with \(\operatorname{St}_G\), and using (7.3). The formula (7.15.1) is due to B. Srinivasan (On the Steinberg character of a finite simple group of Lie type, J. Australian Math. Soc. 12 (1971), 1--14).\par

\EditionChapter{8. Induced and cuspidal representations}

8.1. Let \(P\) be an \(F\)-stable parabolic subgroup of \(G\) and let \(T\subset P\) be an \(F\)-stable maximal torus. We denote by \(U_P\) the unipotent radical of \(P\). The quotient group \(P/U_P\) is a connected reductive algebraic group acted on by \(F\). Let \(\pi:P\to P/U_P\) be the canonical projection. \(\pi\) induces an isomorphism \(T\xrightarrow{\sim}\pi(T)\) hence also an isomorphism \(T^F\xrightarrow{\sim}\pi(T)^F\). Let \(\theta\) be a character of \(T^F\) and let \(\bar\theta\) be the corresponding character of \(\pi(T)^F\). We denote by \(R_{T,P}^\theta\) the image of the virtual representation \(R_{\pi(T)}^{\bar\theta}\) of \((P/U_P)^F\) under the canonical embedding \(\mathcal R((P/U_P)^F)\subset\mathcal R(P^F)\). With these notations, we have the following generalisation of 1.10:\par

\noindent\textbf{PROPOSITION 8.2.} \[
R_T^\theta=\operatorname{Ind}_{P^F}^{G^F}(R_{T,P}^\theta).
\]\par

We choose a Borel subgroup \(B\) in \(G\) such that \(T\subset B\subset P\). Let \(\mathscr P\) be the set of all parabolic subgroups \(P'\) in \(G\) such that \(P'\) and \(P\) are conjugate under \(G^F\); this is a finite set. We have \[
\widetilde X_{T\subset B}=\coprod_{P'\in\mathscr P}\widetilde X(P')
\] where \[
\widetilde X(P')=
\{g\in G\mid g^{-1}Fg\in FU,\ gPg^{-1}=P'\}/U\cap FU,
\] with \(U\) the unipotent radical of \(B\).\par

(Note that \(g^{-1}F(g)\in FU\) implies \(F(gPg^{-1})=gPg^{-1}\).) Let \(P'\in\mathscr P\) and let \(g_1\in G^F\) be such that \(g_1Pg_1^{-1}=P'\). If \(g\in G\), \(g^{-1}F(g)\in FU\), \(gPg^{-1}=P'\), we have \[
g_1^{-1}g\in P,\qquad (g_1^{-1}g)^{-1}F(g_1^{-1}g)\in FU;
\]\par
}

\EditionPage{46}{147}{%
by sending \(g\) to \(\pi(g_1^{-1}g)\) we get an isomorphism \[
\widetilde X(P')\xrightarrow{\sim}\widetilde X_{\pi(T)\subset\pi(B)}.
\] It follows easily that the \(G^F\)-module \(H_c^i(\widetilde X_{T\subset B})\) is canonically isomorphic to the \(G^F\)-module induced by \(H_c^i(\widetilde X_{\pi(T)\subset\pi(B)})\) regarded as a \(P^F\)-module; moreover this isomorphism is compatible with the action of \(T^F\). The proposition follows.\par

\noindent\textbf{THEOREM 8.3.} Let \(T\) be an \(F\)-stable maximal torus in \(G\) such that \(T\) is not contained in any \(F\)-stable, proper parabolic subgroup of \(G\). Let \(\theta\) be a non-singular character of \(T^F\) (cf. 5.19). Then \((-1)^{\sigma(G)-\sigma(T)}R_T^\theta\) can be represented by a cuspidal \(G^F\)-module.\par

Applying (7.5) for \(s\) equal to the identity element of \(G^F\), we see that the regular representation of \(G^F\) is equal to: \[
\operatorname{Ind}_e^{G^F}(1)
=
\frac1{\operatorname{St}_G(e)}
\sum_T\sum_{\theta\in(T^F)^\vee}
(-1)^{\sigma(G)-\sigma(T)}R_T^\theta.
\] We apply this formula with \(G\) replaced by \(P/U_P\) (\(P\) as in (8.1)); we then regard the regular representation of \((P/U_P)^F\) as a representation of \(P^F\) on which \(U_P^F\) acts trivially (i.e., as \(\operatorname{Ind}_{U_P^F}^{P^F}(1)\)) and we induce it to \(G^F\): \[
\begin{aligned}
\operatorname{Ind}_{U_P^F}^{G^F}(1)
&=\operatorname{Ind}_{P^F}^{G^F}\bigl(\operatorname{Ind}_{U_P^F}^{P^F}(1)\bigr)\\
&=\frac1{\operatorname{St}_{P/U_P}(e)}
\sum_{T'\subset P/U_P}
\sum_{\theta'\in(T'^F)^\vee}
(-1)^{\sigma(P/U_P)-\sigma(T')}
\operatorname{Ind}_{P^F}^{G^F}(R_{T'}^{\theta'}),
\end{aligned}
\] where \(R_{T'}^{\theta'}\) is regarded as an element in \(\mathcal R(P^F)\) in the natural way. We now use (8.2) and the fact that any \(F\)-stable maximal torus \(T'\) in \(P/U_P\) can be lifted to an \(F\)-stable maximal torus \(T\) in \(P\) in precisely \(|U_P^F|\) different ways: \[
\tag{8.3.1}
\operatorname{Ind}_{U_P^F}^{G^F}(1)
=
\frac1{\operatorname{St}_G(e)}
\sum_{T\subset P}
\sum_{\theta\in(T^F)^\vee}
(-1)^{\sigma(G)-\sigma(T)}R_T^\theta.
\] If \(P\ne G\) and \(T\) is as in the statement of the theorem, we see from (8.3.1) and (6.8) that \[
\tag{8.3.2}
\left\langle(-1)^{\sigma(G)-\sigma(T)}R_T^\theta,
\operatorname{Ind}_{U_P^F}^{G^F}(1)\right\rangle=0
\qquad\text{for any }\theta\in(T^F)^\vee.
\] If \(\theta\) is non-singular, \((-1)^{\sigma(G)-\sigma(T)}R_T^\theta\) can be represented by an actual \(G^F\)-module (7.4). This \(G^F\)-module is cuspidal by (8.3.2).\par

\EditionChapter{9. A Vanishing Theorem}

9.1. Let \(G\) be a reductive algebraic group over \(k\). Let \(w=s_1\cdots s_n\) be a minimal expression for an element of the Weyl group \(W\) of \(G\) (0.4). The following desingularization of the closure \(\overline{O(w)}\) of \(O(w)\subset X\times X\) (1.2) has been considered by H. C. Hansen [6] and M. Demazure [3].\par

\noindent\textbf{DEFINITION 9.2.} \(\overline O(s_1,\ldots,s_n)\) is the space of sequences \((B_0,\ldots,B_n)\) of\par
}

\EditionPage{47}{148}{%
Borel subgroups, with \(B_{i-1}\) and \(B_i\) in relative position \(e\) or \(s_i\).\par

The maps \[
\overline O(s_1,\ldots,s_n)
\longrightarrow
\overline O(s_1,\ldots,s_{n-1})
\longrightarrow\cdots\longrightarrow
\overline O(\,)=X
\] express \(\overline O=\overline O(s_1,\ldots,s_n)\) as an iterated fibre space over \(X\), with fibre \(\mathbb P^1\) (1.2). The subspace \(D_i\subset\overline O\) where \(B_{i-1}=B_i\) is the inverse image of a section of \(\overline O(s_1,\ldots,s_i)\to\overline O(s_1,\ldots,s_{i-1})\); it is a smooth divisor, and \(D=\bigcup D_i\) is a divisor with normal crossings. The map \((B_0,\ldots,B_n)\mapsto(B_0,B_n)\) induces an isomorphism \(\overline O-D\xrightarrow{\sim}O(w)\), and hence maps \(\overline O\) to \(\overline{O(w)}\): \(\overline O\) is a resolution of singularities for \(\overline{O(w)}\).\par

9.3. Let \(T^*\), \(B^*\) and \(U^*\) be as in (1.7). For any character \(\lambda:T^*\to\mathbb G_m\), the \(T^*\)-torsor \(E\) over \(X\) defined in (1.7) gives rise to a line bundle \(E_\lambda\), provided with \(\lambda:E\to E_\lambda\) such that \(\lambda(et)=\lambda(e)\lambda(t)\). For \(\dot w\in N(T^*)\) with image \(w\) in \(W=N(T^*)/T^*\), the \(w\)-map of \(T^*\)-torsors over \(O(w)\subset X\times X\) constructed in (1.7) induces an isomorphism \[
\Psi(\dot w):\operatorname{pr}_1^*E_\lambda
\xrightarrow{\sim}
\operatorname{pr}_2^*E_{\lambda\circ\operatorname{ad}\dot w}
=
\operatorname{pr}_2^*E_{w^{-1}(\lambda)}
\] which makes the following diagram commute: \[
\begin{array}{ccc}
\operatorname{pr}_1^*E&\xrightarrow{\cdot\dot w}&\operatorname{pr}_2^*E\\
\lambda\downarrow&&\downarrow w^{-1}(\lambda)\\
\operatorname{pr}_1^*E_\lambda&\xrightarrow{\Psi(\dot w)}&\operatorname{pr}_2^*E_{w^{-1}(\lambda)}.
\end{array}
\] We will investigate its behaviour at infinity.\par

9.4. Let \(X(T^*)\) be the character group of \(T^*\) and let \(C\subset X(T^*)\otimes\mathbb R\) be the fundamental chamber (corresponding to \(B^*\)). If \(C_1\) and \(C_2\) are two chambers, we write \(D(C_1,C_2)\) for the intersection of the (closed) radicial half spaces containing both \(C_1\) and \(C_2\), and \(D^0(C_1,C_2)\) for its interior.\par

\noindent\textbf{PROPOSITION 9.5.} Let \(\overline{O(w)}'\) be the normalization of the closure \(\overline{O(w)}\) of \(O(w)\) in \(X\times X\). The map \(\Psi(\dot w):\operatorname{pr}_1^*E_\lambda\to\operatorname{pr}_2^*E_{w^{-1}(\lambda)}\) extends over \(\overline{O(w)}'\) if and only if \(\lambda\in D(C,-wC)\). It vanishes outside of \(O(w)\) if and only if \(\lambda\in D^0(C,-wC)\).\par

Let us consider \(\Psi(\dot w)\) as a rational section of \(\operatorname{pr}_1^*E_{\lambda^{-1}}\otimes\operatorname{pr}_2^*E_{w^{-1}(\lambda)}\); it suffices to prove that, with the notations of (9.2), for each \(i\), the order \(\nu_i\) of \(\Psi(\dot w)\) along the divisor \(D_i\subset\overline O(s_1,\ldots,s_n)\) is \(\ge0\) for \(\lambda\in D(C,-wC)\), \(>0\) for \(\lambda\in D^0(C,-wC)\). The convex region \(D(C,-wC)\) is the intersection of the radicial half spaces containing \(C\) but not \(wC\). Put \(w_i=s_1\cdots s_i\). The gallery \((C,w_1C,w_2C,\ldots,wC)\) connects \(C\) to \(wC\); its walls are all the radicial hyperplanes that separate \(C\) and \(wC\). The wall between \(w_{i-1}C\) and \(w_iC\) is defined\par
}

\EditionPage{48}{149}{%
by the root \(w_{i-1}(\alpha_i)\), where \(\alpha_i\) is the simple root corresponding to \(s_i\); \(C\) and \(w_{i-1}C\) are on the same side of it. The convex region \(D(C,-wC)\) is hence defined by the inequalities \[
\langle\mu,H_{w_{i-1}(\alpha_i)}\rangle\ge0.
\] (Recall that \(H_\alpha\) denotes the coroot corresponding to a root \(\alpha\), cf. (5.1).) We will prove that \[
\tag{9.5.1}
\nu_i=\langle\lambda,H_{w_{i-1}(\alpha_i)}\rangle.
\] Let us lift the decomposition \(w=(s_1\cdots s_{i-1})s_i(s_{i+1}\cdots s_n)\) into a decomposition \(\dot w=\dot w_{i-1}\dot s_i\dot w_i'\) in \(N(T^*)\). We have \[
\Psi(\dot w)=\Psi(\dot w_i')\Psi(\dot s_i)\Psi(\dot w_{i-1}).
\] When we approach a general point of \(D_i\), the isomorphisms \(\Psi(\dot w_i')\) and \(\Psi(\dot w_{i-1})\) extend, and we are left to prove that the order along \(D_i\) of \[
\Psi(\dot s_i):E_{w_{i-1}^{-1}(\lambda)}\longrightarrow E_{w_i^{-1}(\lambda)}
\] is \[
\langle\lambda,H_{w_{i-1}(\alpha_i)}\rangle
=
\langle w_{i-1}^{-1}(\lambda),H_{\alpha_i}\rangle;
\] we are reduced to the case where \(w\) is a fundamental reflection: all computations occur within a minimal parabolic subgroup \(P\), or \(P/U_P\), or the derived group of \(P/U_P\), or its universal covering \(\operatorname{SL}(2)\). Let us make a direct check for \(G=\operatorname{SL}(2)\), \(w\ne e\) and \(\lambda\) the fundamental weight.\par

(a) We take \(\operatorname{SL}(2)\) in its obvious representation \(k^2\), \(T^*=\) diagonal matrices, \(B^*=\) upper triangular matrices. The weight \(\lambda\) is \[
\begin{pmatrix}a&*\\0&a^{-1}\end{pmatrix}\longmapsto a,
\qquad\text{and}\qquad
\langle\lambda,H_\alpha\rangle=1.
\] (b) \(X=\mathbb P^1\), the space of homogeneous lines of \(k^2\) and the fibre of \(E_\lambda\) at \(x\) is the corresponding line.\par

(c) Up to a scalar, \(\Psi(\dot w)\) associates to \(u\in(E_\lambda)_x\) the linear form \(u\wedge v\) on \((E_\lambda)_y\), for \(x\ne y\). It vanishes simply for \(y\to x\).\par

9.6. The assumptions (0.1.1) are in force from now on. Under the assumptions of (1.8), we have \[
F^*E_\lambda=E_{F^*\lambda}\qquad(\text{where }F^*\lambda=\lambda\circ F).
\] On the inverse image \(\overline{X(w)}'\) of the graph of the Frobenius map \(F:X\to X\) in \(\overline{O(w)}'\), the line bundle \(\operatorname{pr}_1^*E_{\lambda^{-1}}\otimes\operatorname{pr}_2^*E_{w^{-1}(\lambda)}\) is hence isomorphic to \[
\operatorname{pr}_1^*(E_{F^*w^{-1}\lambda-\lambda}).
\] It will be ample if \(E_{F^*w^{-1}\lambda-\lambda}\) is ample on \(X\), i.e., if \(F^*w^{-1}\lambda-\lambda\in-C^0\). On the\par
}

\EditionPage{49}{150}{%
other hand, if \(\lambda\in D^0(C,-wC)\), the section \(\Psi(\dot w)\) of it has as zero set the complement of \(X(w)\) in the projective variety \(\overline{X(w)}'\). If both conditions can be simultaneously fulfilled, then \(X(w)\) is affine. In terms of \(\mu=-w^{-1}(\lambda)\), the conditions read \[
\tag{9.6.1}
\mu\in D^0(C,-w^{-1}C),
\]  \[
\tag{9.6.2}
F^*\mu-w\mu\in C^0.
\]\par

\noindent\textbf{THEOREM 9.7.} If there is \(\mu\in X(T^*)\otimes\mathbb R\) satisfying (9.6.1) and (9.6.2) then \(X(w)\) is affine. As a consequence, \(X(w)\) is affine as soon as \(q\) is larger than the Coxeter number \(h\) of \(G\).\par

It remains to prove that if \(q\ge h\), then some \(\mu\) satisfies (9.6.1) and (9.6.2). The first condition will be fulfilled if \(\mu\in C^0\), that is, if \(\langle\mu,H_\alpha\rangle>0\) for each simple root \(\alpha\). We take \(\mu\) so that \(\langle\mu,H_\alpha\rangle=1\) for each simple root \(\alpha\). We then have \(\langle F\mu,H_\alpha\rangle=q\), and \(\langle w\mu,H_\alpha\rangle=\langle\mu,w^{-1}H_\alpha\rangle\). If \(H=\sum n_\alpha H_\alpha\) is the highest coroot, \(\sum n_\alpha=h-1\) and \[
\langle F\mu-w\mu,H_\alpha\rangle
\ge q-\langle\mu,H\rangle
=q-h+1>0,
\] hence (9.6.2).\par

\noindent\textbf{THEOREM 9.8.} If the character \(\theta:T(w)^F\to\overline{\mathbb Q}_l^*\) is non-singular, then \[
H_c^*(\widetilde X(\dot w),\overline{\mathbb Q}_l)_\theta
\xrightarrow{\sim}
H^*(\widetilde X(\dot w),\overline{\mathbb Q}_l)_\theta.
\]\par

\noindent\textbf{COROLLARY 9.9.} If \(X(w)\) is affine and \(\theta\) is non-singular, then \[
H_c^i(\widetilde X(\dot w),\overline{\mathbb Q}_l)_\theta=0
\qquad\text{for }i\ne l(w).
\]\par

Let us deduce (9.9) from (9.8). With the notations of (1.9), (9.8) means that \[
\tag{9.9.1}
H_c^*(X(w),\mathcal F_\theta)
\xrightarrow{\sim}
H^*(X(w),\mathcal F_\theta).
\] If \(X(w)\) is affine, then \(H^i(X(w),\mathcal F_\theta)=0\) for \(i>\dim X(w)=l(w)\) ([10, XIV, 3.2]). The sheaf \(\mathcal F_\theta\) is locally constant, its dual is \(\mathcal F_{\theta^{-1}}\) and \(X(w)\) is smooth and purely of dimension \(l(w)\). By Poincaré duality, \(H_c^i(X(w),\mathcal F_\theta)\) is dual (up to a twist) to \(H^{2l(w)-i}(X(w),\mathcal F_{\theta^{-1}})\), hence vanishes for \(2l(w)-i>l(w)\), i.e., for \(i<l(w)\). This proves the corollary.\par

9.10. We first construct a nice compactification of \(X(w)\). Let \(w=s_1\cdots s_n\) be a minimal expression for \(w\). We define \(\overline X(s_1,\ldots,s_n)\) to be the space of sequences \((B_0,\ldots,B_n)\) of Borel subgroups of \(G\), with \(B_n=FB_0\) and with \(B_{i-1}\) and \(B_i\) in relative position \(s_i\) or \(e\). It is the inverse image of the graph of Frobenius by the map \[
\alpha:(B_0,\ldots,B_n)\longmapsto(B_0,B_n):
\overline O(s_1,\ldots,s_n)\longrightarrow X\times X.
\]\par
}

\EditionPage{50}{151}{%
In other words, if \(\Gamma\xrightarrow{i}X\times X\) is the inclusion in \(X\times X\) of the graph of Frobenius, it is the fibre product \(\overline O\times_{X\times X}\Gamma\): \[
\begin{array}{ccc}
\overline X(s_1,\ldots,s_n)&\longrightarrow&\Gamma\\
\downarrow&&\downarrow i\\
\overline O(s_1,\ldots,s_n)&\longrightarrow&X\times X.
\end{array}
\]\par

\noindent\textbf{LEMMA 9.11.} \(\Gamma\) is transverse to \(\overline O(s,\ldots,s)\), as well as to any intersection of the divisors \(D_i\). The fibre product \(\overline X(s,\ldots,s_n)\) is hence a smooth compactification of \(X(w)\), with a divisor with normal crossings (sum of the traces \(\overline D_i\) of the \(D_i\)) at infinity.\par

We will show that \(\Gamma\) is transverse to any smooth \(G\)-equivariant scheme \(\pi:Y\to X\times X\) over \(X\times X\) (with \(G\) acting diagonally on \(X\times X\)); that is, if \(\pi(y)\in\Gamma\), the sum of the tangent space to \(\Gamma\) at \(\pi(y)\) and of the image of \(d\pi\) is the whole tangent space of \(X\times X\) at \(\pi(y)\). Indeed\par

(a) the tangent space of \(\Gamma\) at \(\pi(y)\) is \(T_x\times\{0\}\);\par

(b) the image of \(d\pi\) contains the image of \(\operatorname{Lie}(G)\) by the derivative at \(e\in G\) of \(g\mapsto g\pi y\) (by the equivariance of \(Y\)). Since the space \(X\) is homogeneous with reduced stabilizers this image projects onto \(T_x\) by the second projection.\par

9.12. We now investigate how the covering \(\widetilde X(\dot w)\) (with structural group \(T(w)^F\)) ramifies along the divisor \(\overline D_i\). The structural group being of order prime to \(p\), the ramification is tame. On each connected component of \(\overline D_i\), it gives rise to a homomorphism \(\tau_i:\widehat{\mathbb Z}(1)\to T(w)^F\).\par

Let \(x(t)\in\overline X(s_1,\ldots,s_n)\) be a one parameter family with \(x(0)\) a point of \(\overline D_i\) (not on any \(\overline D_j\), \(j\ne i\)), and with the tangent vector \(\dot x(0)\) transverse to \(\overline D_i\). In technical terms: for \(S=\operatorname{spectrum}\) of the Henselization of \(k[t]\) at \((t)\), \(x\) is a morphism \(x:S\to\overline X(s_1,\ldots,s_n)\) such that the closed point \(t=0\) is mapped to a point of \(\overline D_i\) not on any \(\overline D_j\), \(j\ne i\), and that the inverse image of \(\overline D_i\) is the reduced scheme \(t=0\). Put \(x(t)=(B_0(t),\ldots,B_n(t))\) and let \(E_i\) be the pull back by \(t\mapsto B_i(t)\) of the \(T\)-torsor \(E\) on \(X\) (1.7). If we factor \(\dot w\) as in 9.5, the isomorphism \(\Psi(\dot w):E_0\to E_n\) (over the generic point of \(S\)) factors as \[
\Psi(\dot w)=\Psi(\dot w_i')\Psi(\dot s_i)\Psi(\dot w_{i-1})
\] and \(\Psi(\dot w_{i-1})\) and \(\Psi(\dot w_i')\) extend over \(S\); \(\Psi(\dot s_i)\) doesn't; rather, the composite \(x\mapsto\Psi(\dot s_i)(xH_{-\alpha}(t))\) does (9.5), hence \(\Psi(\dot w)\) is of the form \[
\Psi(\dot w)(x)=\Psi_0\bigl(xH_{w_{i-1}(\alpha_i)}(t)\bigr)
\] where \(\Psi_0:E_0\to E_n\) extends over \(S\). The pull back under \(x\) of the \(T(w)^F\)-torsor \(\widetilde X(\dot w)\), to the generic point of \(S\), is given by the equation \[
\tag{9.12.1}
F(u)=\Psi_0\bigl(uH_{w_{i-1}(\alpha_i)}(t)\bigr).
\] Let \(u_0\) be a solution, over \(S\), of the equation\par
}

\EditionPage{51}{152}{%
\[
F(u_0)=\Psi_0(u_0)
\] (this exists, because \(S\) is strictly Henselian and \(F\) and \(\Psi_0\) extend over \(S\)); if we put \(u=u_0y\), (9.12.1) becomes \[
\tag{9.12.2}
y^{-1}\operatorname{ad}\dot w F(y)=H_{w_{i-1}(\alpha_i)}(t).
\] In other words:\par

\noindent\textbf{LEMMA 9.13.} The \(T(w)^F\)-torsor \(\widetilde X(\dot w)\) ramifies along \(\overline D_i\) in the same way as the pull back under \(H_{w_{i-1}(\alpha_i)}:\mathbb G_m\to T\) of the Lang covering of \(T(w)\) ramifies at 0.\par

The tame fundamental group of \(\mathbb G_m\) is \(\widehat{\mathbb Z}_{p'}(1)\), that of \(T\) is \(Y_{p'}^{\wedge}(1)\) (for \(Y\) the dual of the character group of \(T\)) and \(H_{w_{i-1}(\alpha_i)}\) induces \(x\mapsto xH_{w_{i-1}(\alpha_i)}\): \[
\widehat{\mathbb Z}_{p'}(1)\longrightarrow Y_{p'}^{\wedge}(1).
\] The Lang covering of \(T(w)^F\) gives rise to the map \(Y_{p'}^{\wedge}(1)\to T(w)^F\) described in (5.2). Hence the pull back (9.13) corresponds to the composite \[
\widehat{\mathbb Z}_{p'}(1)
\xrightarrow{H_{w_{i-1}(\alpha_i)}}
Y_{p'}^{\wedge}(1)
\longrightarrow T(w)^F.
\]\par

For \(\theta\) a character of \(T(w)^F\), \(\mathcal F_\theta\) will ramify along \(\overline D_i\) if and only if \(\theta\) is not trivial on \(H_{w_{i-1}(\alpha_i)}\widehat{\mathbb Z}_{p'}(1)\). In that case, all higher direct images \(R^ij_*\mathcal F_\theta\) of \(\mathcal F_\theta\) (\(i\ge0\)) under \(j:X(w)\to\overline X(s_1,\ldots,s_n)\) will vanish on \(\overline D_i\). In particular, we have\par

\noindent\textbf{LEMMA 9.14.} If \(\theta\) is non-singular, \(j_*\mathcal F_\theta=j_!\mathcal F_\theta\) (extension by zero) and \(R^ij_*\mathcal F_\theta=0\) for \(i>0\).\par

Proof of 9.8 (in the guise 9.9.1). By (9.14), the Leray spectral sequence for \(j\) reads \[
H^*(\overline X(s_1,\ldots,s_n),j_!\mathcal F_\theta)
\xrightarrow{\sim}
H^*(X(w),\mathcal F_\theta).
\] The space \(\overline X(s_1,\ldots,s_n)\) being a compactification of \(X(w)\), the left hand side is by definition \(H_c^*(X(w),\mathcal F_\theta)\).\par

Remark 9.15.1. By using arguments parallel to those of (1.6), one can show that the conclusion of (9.9): “for \(\theta\) non-singular, \(H_c^i(\widetilde X(\dot w),\overline{\mathbb Q}_l)_\theta=0\) for \(i\ne l(w)\)” holds true as soon as there is \(w'\) in the \(F\)-conjugacy class of \(w\) such that \(X(w')\) is affine. The criterion (9.7) can be used to check that this is always the case for the classical groups and for \(G_2\) except the case \((G_2,q=2,w=\text{Coxeter})\) for which a different method applies.\par

Remark 9.15.2. Let \((T',\theta')\) be maximally split (5.17); under an isomorphism \(T(w)\xrightarrow{\sim}T'\) (inducing \(T(w)^F\xrightarrow{\sim}T'^{F}\)), as in (1.18), \(\theta'\) becomes a character \(\theta\) of \(T(w)^F\). Proposition 5.24 and the proof of the Induction Theorem\par
}

\EditionPage{52}{153}{%
8.2 show that the conclusion of (9.8) still holds for \(\theta\). The same applies to (9.9).\par

\noindent\textbf{THEOREM 9.16.} If \(u\in G^F\) is a regular unipotent element, then, for any \(F\)-stable maximal torus \(T\), we have \(Q_{T,G}(u)=1\).\par

We must prove that for any \(w\in\mathbf W\), \(\operatorname{Tr}(u,H_c^*(X(w)))=1\). We know that \(\operatorname{Tr}(u,H_c^*(X(w)))\) is an integer, and that \[
\sum_{w\in\mathbf W}\operatorname{Tr}(u,H_c^*(X(w)))=|\mathbf W|
\] (see (7.14) and its proof). Hence it suffices to prove the\par

\noindent\textbf{LEMMA 9.17.} Under the assumptions of (9.16), \(\operatorname{Tr}(u,H_c^*(X(w)))>0\), for any \(w\in\mathbf W\).\par

Take \(w=s_1\cdots s_n\) as in (9.10) and let \(\overline X(s_1,\ldots,s_n)\) be the corresponding compactification of \(X(w)\). For \(P\subset[1,n]\) we write \(\overline D_P\) for the intersection, in \(\overline X(s_1,\ldots,s_n)\), of the divisors \(\overline D_i\) (\(i\in P\)) (9.11). For \(P=\varnothing\), \(\overline D_P\) is \(\overline X(s_1,\ldots,s_n)\). The additivity property of cohomology with compact support shows that \[
\tag{9.17.1}
\operatorname{Tr}(u,H_c^*(X(w)))
=
\sum_{P\subset[1,n]}(-1)^{|P|}\operatorname{Tr}(u,H^*(\overline D_P)).
\]\par

Let \(b\) be the unique fixed point of \(u\) in \(X\). Then \(a=(b,\ldots,b)\) is the unique fixed point of \(u\) in \(\overline X(s_1,\ldots,s_n)\); it is contained in each \(\overline D_P\). The variety \(\overline D_P\) being smooth and compact, \(\operatorname{Tr}(u,H^*(\overline D_P))\) is just the multiplicity of the unique fixed point \(a\) of \(u\) acting on \(\overline D_P\); to compute it, we may replace \(\overline D_P\) by its completion \(\widehat D_P\) at \(a\). Let \((x_i)\) be a formal coordinate system for \(\overline X(s_1,\ldots,s_n)\) at \(a\), such that \(x_i=0\) is an equation for \(\widehat D_i\). In terms of these coordinates, \(u\) maps the point with coordinates \(x_1,\ldots,x_n\) into the point with coordinates \(P_1(x_1,\ldots,x_n),\ldots,P_n(x_1,\ldots,x_n)\), for some formal power series \(P_i\). Since the divisor \(\widehat D_i\) is preserved by \(u\), \(P_i\) is divisible by \(x_i\). The jacobian matrix of \(u\) at \((0,\ldots,0)\) is hence diagonal; \(u\) being of order a power of \(p\), it is the identity, and we have \[
P_i=x_i(1+Q_i).
\] with \(Q_i(0,\ldots,0)=0\). The fixed point scheme of \(u\) is given by the equations \(x_iQ_i=0\).\par

Let \(E_i\) be the divisor \(Q_i=0\). It is non-empty. The formula (9.17.1) can be rewritten \[
\begin{aligned}
\operatorname{Tr}(u,H_c^*(X(w)))
&=\sum_{P\subset[1,n]}(-1)^{|P|}
  \prod_{i\in P}\widehat D_i
  \prod_{j\notin P}(\widehat D_j+E_j)\\
&=\sum_{P\subset Q\subset[1,n]}(-1)^{|P|}
  \prod_{i\in Q}\widehat D_i
  \prod_{j\notin Q}E_j\\
&=\sum_{Q\subset[1,n]}
  \left(\sum_{P\subset Q}(-1)^{|P|}\right)
  \prod_{i\in Q}\widehat D_i
  \prod_{j\notin Q}E_j\\
&=\prod_jE_j>0,
\end{aligned}
\]\par
}

\EditionPage{53}{154}{%
(the product of \(n\) divisors stands for the multiplicity of their intersection).\par

\noindent\textbf{PROPOSITION 9.18.} Let \(s\in G^F\) be a semisimple element. Let \(\nu\) be the class function on \(G^F\) defined by: \[
\nu(g)=
\begin{cases}
|Z^0Z^0(s)^F|\,q^{l(Z^0(s))}\dfrac{|Z(s)^F|}{|Z^0(s)^F|},
&\text{if }g\in G^F\text{ is regular with semisimple}\\
&\text{part conjugate in }G^F\text{ to }s,\\[2mm]
0,&\text{otherwise}.
\end{cases}
\] Then \(\nu\) is the character of \[
\frac1{|Z^0(s)^F|}
\sum_{\substack{T\\s\in T}}|T^F|
\sum_{\theta\in(T^F)^{\sim}}\theta(s)^{-1}R_T^\theta.
\] (Here \(Z^0Z^0(s)\) is the connected centre of the connected centralizer of \(s\), and \(l(Z^0(s))\) is the semisimple rank of \(Z^0(s)\).) The proof is similar to that of (7.5); it uses (9.16) and (4.2) instead of (7.2).\par

For any \(\rho\in\mathcal R(G^F)\) we have \[
\tag{9.18.1}
\langle\rho,\nu\rangle
=
\frac{|Z^0Z^0(s)^F|}{|Z^0(s)^F|}
q^{l(Z^0(s))}
\sum_{u\in A(s)}\rho(su)
\] where \(A(s)\) is the set of regular unipotent elements in \(Z^0(s)^F\).\par

\noindent\textbf{COROLLARY 9.19.} For any \(\rho\in\mathcal R(G^F)\), the average value of the character of \(\rho\) on the regular elements in \(G^F\) with semisimple part \(s\), equals \[
\frac1{|A(s)|}\sum_{u\in A(s)}\operatorname{Tr}(su,\rho)
=
\frac1{|Z^0(s)^F|}
\sum_{\substack{T\\s\in T}}|T^F|
\sum_{\theta\in(T^F)^{\sim}}\theta(s)\langle\rho,R_T^\theta\rangle.
\] This follows from (9.18), (9.18.1) and the following lemma.\par

\noindent\textbf{LEMMA 9.20.} The number of regular unipotent elements in \(G^F\) equals \[
|G^F|/|Z^{0F}|\,q^l.
\] (Here \(Z^0\) is the connected centre of \(G\) and \(l\) is the semisimple rank of \(G\).)\par

To count regular unipotents in \(G^F\), we observe that any such element is contained in a unique Borel subgroup and then count the regular unipotents in \(G^F\) contained in a fixed \(F\)-stable Borel subgroup.\par

\EditionChapter{10. A decomposition of the Gelfand-Graev representation}

10.1. Let \((G^F)^{\sim}\) be the set of isomorphism classes of irreducible representations of \(G^F\) (over \(\overline{\mathbb Q}_l\)). For any \(\rho\in(G^F)^{\sim}\), there exists an \(F\)-stable maximal torus \(T\) and \(\theta\in(T^F)^{\sim}\) such that \(\langle\rho,R_T^\theta\rangle\ne0\) (7.7); moreover the geometric conjugacy class \([\theta]\) of \((T,\theta)\) (see 5.5), is uniquely determined by \(\rho\) (6.3). We thus get a well-defined surjective map \[
\tag{10.1.1}
(G^F)^{\sim}\longrightarrow\mathfrak S
\]\par
}

\EditionPage{54}{155}{%
where \(\mathfrak S\) is the set of all geometric conjugacy classes of pairs \((T,\theta)\).\par

10.2. In the rest of this chapter we shall assume that the centre \(Z\) of \(G\) is connected. Let \(T^*,B^*\) be as in (1.8). Let \(U^*\) be the unipotent radical of \(B^*\) and let \(U_\bullet^*\) be the subgroup of \(U^*\) generated by the root subgroups corresponding to non-simple roots. The quotient \(U^*/U_\bullet^*\) is commutative and is a direct product over the simple roots \(\alpha\): \(\prod_\alpha U_\alpha^*\), with \(U_\alpha^*\) one-dimensional. Let \(I\) be the set of orbits of \(F\) on the simple roots. For any \(i\in I\), let \(U_i^*=\prod_{\alpha\in i}U_\alpha^*\); then \(U^*/U_\bullet^*=\prod_{i\in I}U_i^*\). This decomposition is \(F\)-stable, hence we have also \(U^{*F}/U_\bullet^{*F}=\prod_{i\in I}U_i^{*F}\).\par

We consider the Gelfand-Graev representation \(\Gamma_G=\operatorname{Ind}_{U^{*F}}^{G^F}(\chi)\) where \(\chi\) is any character of \(U^{*F}\) which is trivial on \(U_\bullet^{*F}\) and defines a non-trivial character of \(U_i^{*F}\) for all \(i\in I\). All such \(\chi\) are conjugate under \(T^{*F}\), since \(Z\) is connected, hence the \(G^F\)-module \(\Gamma_G\) is well-defined up to isomorphism.\par

Let \(\Delta_G\) be the class function on \(G^F\) which equals \(|Z^F|q^l\) on any regular unipotent element in \(G^F\) and vanishes on all other elements. (\(l\) is the semisimple rank of \(G\).) The following result shows that this is the character of a virtual representation of \(G^F\) (which will be also denoted by \(\Delta_G\)):\par

\noindent\textbf{PROPOSITION 10.3.} For any subset \(J\subset I\), let \(P(J)\supset B^*\) be the parabolic subgroup generated by \(B^*\) and by the root subgroups corresponding to minus the simple roots in \(F\)-orbits in \(J\). Let \(L(J)\) be \(P(J)\) modulo its unipotent radical. Then \[
\tag{10.3.1}
\Delta_G
=
\sum_{J\subset I}(-1)^{|J|}
\operatorname{Ind}_{P(J)^F}^{G^F}(\Gamma_{L(J)}).
\] (Note that the centre of \(L(J)\) is connected, since that of \(G\) is, hence \(\Gamma_{L(J)}\) is well-defined.)\par

Remark 10.4. The following identity is a formal consequence of (10.3.1): \[
\Gamma_G
=
\sum_{J\subset I}(-1)^{|J|}
\operatorname{Ind}_{P(J)^F}^{G^F}(\Delta_{L(J)}).
\]\par

10.5. Proof of 10.3. For any subset \(J\subset I\), let \(\mathcal C(J)\) be the set of characters \(\chi\) of \(U^{*F}\) with \(\chi|U_\bullet^{*F}=1\) and such that \(\chi\) induces a non-trivial character of \(U_i^*\) (\(i\in I\)) if and only if \(i\in J\). It is easy to see that, for any \(\chi\in\mathcal C(J)\), we have \[
\operatorname{Ind}_{P(J)^F}^{G^F}(\Gamma_{L(J)})
=
\operatorname{Ind}_{U^{*F}}^{G^F}(\chi).
\]\par

Let \(u_0\in U^{*F}\) be a fixed regular unipotent element; in other words, the image of \(u_0\) in \(U_i^{*F}\) is non-zero for any \(i\in I\). It follows that \[
\sum_{\chi\in\mathcal C(J)}\chi(u_0)=(-1)^{|J|}.
\] Hence (10.3.1) is equivalent to:\par
}

\EditionPage{55}{156}{%
\[
\tag{10.5.1}
\Delta_G
=
\sum_{\chi\in\mathcal C}\chi(u_0)\operatorname{Ind}_{U^{*F}}^{G^F}(\chi)
\] where \(\mathcal C=\bigcup_{J\subset I}\mathcal C(J)\). The character of the right hand side of (10.5.1) vanishes at non-unipotent elements. Its value at \(u\in U^{*F}\) is: \[
\tag{10.5.2}
\begin{aligned}
&|U^{*F}|^{-1}
\sum_{\substack{g\in G^F\\g^{-1}ug\in U^{*F}}}
\sum_{\chi\in\mathcal C}\chi(g^{-1}ugu_0)\\
&\qquad=q^l|U^{*F}|^{-1}
\#\{g\in G^F\mid g^{-1}ug\in U^{*F},\ g^{-1}ugu_0\in U_\bullet^{*F}\}.
\end{aligned}
\] This is zero unless \(u\in U^{*F}\) is regular unipotent; we now assume that \(u\in U^{*F}\) is regular unipotent. For any \(g\in G^F\) with \(g^{-1}ug\in U^{*F}\), we have \(u\in B^*\cap gB^*g^{-1}\); but \(B^*\) is the only Borel subgroup containing \(u\), hence \(g\in B^{*F}\). If \(g=tu'\), \(t\in T^{*F}\), \(u'\in U^{*F}\), the relation \(g^{-1}ugu_0\in U_\bullet^{*F}\) determines \(t\) uniquely modulo \(Z^F\) and leaves \(u'\) arbitrary. Thus the expression (10.5.2) becomes \[
q^l|U^{*F}|^{-1}|Z^F||U^{*F}|=|Z^F|q^l,
\] and (10.3) is proved.\par

We shall now prove the following\par

\noindent\textbf{LEMMA 10.6.} Let \(M\) be a virtual representation of \(G^F\) such that \[
\langle M,M\rangle=|Z^F|q^l
\qquad\text{and}\qquad
\langle M,R_T^\theta\rangle=\varepsilon_T^\theta=\pm1
\] for any \(F\)-stable maximal torus \(T\subset G\) and any \(\theta\in T^{\vee F}\). Then \[
M=\sum_{x\in\mathfrak S}\sigma_xM_x
\] where \(\sigma_x=\pm1\) and \(M_x\) are distinct irreducible representations of \(G^F\) such that \[
M_x
=
\sigma_x
\sum_{\substack{(T,\theta)\bmod G^F\\{}[\theta]=x}}
\frac{\varepsilon_T^\theta}{\langle R_T^\theta,R_T^\theta\rangle}R_T^\theta.
\]\par

Let \(\{M\}\) be the set of all \(\rho\in(G^F)^{\sim}\) such that \(\langle\rho,M\rangle\ne0\); \(\{M\}\) has at most \(|Z^F|q^l\) elements. By our assumption and by 5.7(i), the map \(\{M\}\to\mathfrak S\) (the restriction of (10.1.1)) is surjective. Since \(|\mathfrak S|=|Z^F|q^l\) (5.7), it is bijective and \(\{M\}\) has exactly \(|Z^F|q^l\) elements. Let \(M_x\in\{M\}\) be the element corresponding to \(x\in\mathfrak S\) under this bijection; we must have \(M=\sum_{x\in\mathfrak S}\sigma_xM_x\), \(\sigma_x=\pm1\), since \(\langle M,M\rangle=|\{M\}|\). If \((T,\theta)\) is such that \([\theta]=x\), we have \[
\langle M,R_T^\theta\rangle
=
\sigma_x\langle M_x,R_T^\theta\rangle
=
\varepsilon_T^\theta;
\] from the orthogonality of the \(R_T^\theta\)'s (6.8) we have \[
M_x
=
\sigma_x
\sum_{\substack{(T,\theta)\bmod G^F\\{}[\theta]=x}}
\frac{\varepsilon_T^\theta}{\langle R_T^\theta,R_T^\theta\rangle}R_T^\theta
+M_x'
\] where \(M_x'\) is orthogonal to all the \(R_T^\theta\)'s. It remains to prove that \(M_x'=0\). Since \(M_x\) is irreducible, this would follow if we prove that \(\langle M_x-M_x',M_x-M_x'\rangle=1\); we certainly have \(\langle M_x-M_x',M_x-M_x'\rangle\le1\), hence it is sufficient to prove\par
}

\EditionPage{56}{157}{%
that \[
\sum_{x\in\mathfrak S}\langle M_x-M_x',M_x-M_x'\rangle=|Z^F|q^l
\] or, in other words, that \[
\sum_{(T,\theta)\bmod G^F}
\frac1{\langle R_T^\theta,R_T^\theta\rangle}
=
|Z^F|q^l.
\] By (6.8), this is equivalent to the identity \[
\sum_{T\bmod G^F}\frac1{|W(T)^F|}|T^F|=|Z^F|q^l,
\] which is the same as \[
\frac1{|\mathbf W|}\sum_{w\in\mathbf W}\det(q-w\tau)=q^l
\] where \(w\) is regarded as an automorphism of the lattice \(X_0=X(T/Z)\) and \(\tau\) is defined by \(F^*=q\cdot\tau^{-1}\). The last identity (compare Steinberg [15, p. 91]) can be written as \[
\sum_{0\le i\le l}(-1)^iq^{l-i}\frac1{|\mathbf W|}
\sum_{w\in\mathbf W}\operatorname{Tr}(w\tau,\Lambda^i(X_0))
=q^l
\] and follows from the fact that \(\Lambda^i(X_0)^{\mathbf W}=0\) for \(i>0\) (5.8).\par

\noindent\textbf{THEOREM 10.7.} We recall that \(Z\) is connected. Let \(x\in\mathfrak S\), and \(x'\) be a representative of the corresponding semisimple conjugacy class in the dual group (5.24). We define \(\delta_x\) to be the \(\mathbf F_q\)-rank of the centralizer of \(x'\).\par

(i) The formulae \[
\tag{10.7.1}
\rho_x
=
\sum_{\substack{(T,\theta)\bmod G^F\\{}[\theta]=x}}
\frac{(-1)^{\sigma(G)-\sigma(T)}}{\langle R_T^\theta,R_T^\theta\rangle}R_T^\theta
\] and \[
\tag{10.7.2}
\rho_x'
=
(-1)^{\sigma(G)-\delta_x}
\sum_{\substack{(T,\theta)\bmod G^F\\{}[\theta]=x}}
\frac1{\langle R_T^\theta,R_T^\theta\rangle}R_T^\theta
\] define irreducible representations of \(G^F\). The \(|Z^F|q^l\) elements \(\rho_x\in(G^F)^{\sim}\), \(x\in\mathfrak S\) (resp. \(\rho_x'\in(G^F)^{\sim}\), \(x\in\mathfrak S\)) are distinct. The maps \(x\mapsto\rho_x\), \(x\mapsto\rho_x'\) are two cross-sections of the map (10.1.1).\par

(ii) One has \[
\tag{10.7.3}
\Gamma_G
=
\sum_{x\in\mathfrak S}\rho_x
=
\sum_{(T,\theta)\bmod G^F}
\frac{(-1)^{\sigma(G)-\sigma(T)}}{\langle R_T^\theta,R_T^\theta\rangle}R_T^\theta,
\]  \[
\tag{10.7.4}
\Delta_G
=
\sum_{x\in\mathfrak S}(-1)^{\sigma(G)-\delta_x}\rho_x'
=
\sum_{(T,\theta)\bmod G^F}
\frac1{\langle R_T^\theta,R_T^\theta\rangle}R_T^\theta.
\]\par
}

\EditionPage{57}{158}{%
(iii) Assume that all the roots have the same length. Fix \((T,\theta)\) in the geometric conjugacy class \(x\), let \(S\) be the connected centralizer of \((T,\theta)\) (5.19) and let \(\theta_S\) be the corresponding character of \(S^F\). We have \[
\tag{10.7.5}
\rho_x\otimes\operatorname{St}_G
=
\operatorname{Ind}_{S^F}^{G^F}(\theta_S\otimes\operatorname{St}_S\otimes\operatorname{St}_S)
\]  \[
\tag{10.7.6}
\rho_x'\otimes\operatorname{St}_G
=
\operatorname{Ind}_{S^F}^{G^F}(\theta_S\otimes\operatorname{St}_S).
\]\par

(iv) Let \(S^*\) be the centralizer of \(x'\). We still have \[
\tag{10.7.7}
\dim\rho_x
=
\frac{|G^F|/\operatorname{St}_G(e)}{|S^{*F}|/\operatorname{St}_{S^*}(e)}\cdot\operatorname{St}_{S^*}(e),
\]  \[
\tag{10.7.8}
\dim\rho_x'
=
\frac{|G^F|/\operatorname{St}_G(e)}{|S^{*F}|/\operatorname{St}_{S^*}(e)}.
\]\par

If we tensor the right hand side of (10.7.1) by \(\operatorname{St}_G\) we find (by 7.3): \[
\tag{10.7.9}
\sum_{\substack{(T,\theta)\bmod G^F\\{}[\theta]=x}}
\frac1{\langle R_T^\theta,R_T^\theta\rangle}
\operatorname{Ind}_{T^F}^{G^F}(\theta).
\]\par

Let us first assume that all roots have the same length. By 5.14(ii), the \(G^F\)-conjugacy classes of pairs \((T,\theta)\), \(T\subset G\), with \([\theta]=x\) are in one-to-one correspondence with \(S^F\)-conjugacy classes of pairs \((T',\theta_S|_{T'})\), with \(T'\) an \(F\)-stable maximal torus in \(S\). Moreover, by 5.14(i) and (6.8), we have \(\langle R_T^\theta,R_T^\theta\rangle=|W_S(T')^F|\). Thus the expression (10.7.9) equals \[
\begin{aligned}
&\sum_{\substack{T'\subset S\\\bmod S^F}}
\frac1{|W_S(T')^F|}\operatorname{Ind}_{T'^F}^{G^F}(\theta_S|_{T'})\\
&=\sum_{\substack{T'\subset S\\\bmod S^F}}
\frac1{|W_S(T')^F|}
\operatorname{Ind}_{S^F}^{G^F}\bigl(\operatorname{Ind}_{T'^F}^{S^F}(1)\otimes\theta_S\bigr)\\
&=\operatorname{Ind}_{S^F}^{G^F}\left(
\sum_{\substack{T'\subset S\\\bmod S^F}}
\frac1{|W_S(T')^F|}\operatorname{Ind}_{T'^F}^{S^F}(1)\otimes\theta_S
\right)\\
&=\operatorname{Ind}_{S^F}^{G^F}(\theta_S\otimes\operatorname{St}_S\otimes\operatorname{St}_S),
\qquad\text{by 7.14}.
\end{aligned}
\] Similarly, by tensoring the right hand side of (10.7.2) by \(\operatorname{St}_G\), we find \[
\operatorname{Ind}_{S^F}^{G^F}\left(
\sum_{\substack{T'\subset S\\\bmod S^F}}
\frac{(-1)^{\sigma(S)-\sigma(T')}}{|W_S(T')^F|}
\operatorname{Ind}_{T'^F}^{S^F}(1)\otimes\theta_S
\right)
=
\operatorname{Ind}_{S^F}^{G^F}(\theta_S\otimes\operatorname{St}_S),
\qquad\text{by 7.14}.
\]\par

In general, the \(G^F\)-conjugacy classes of pairs \((T,\theta)\), with \([\theta]=x\), are in one-to-one correspondence with the \(G^{*F}\)-conjugacy classes of pairs \((T',\theta')\), with \(\theta'\) conjugate to \(x'\) (5.24), i.e., with the \(S^{*F}\)-conjugacy classes of pairs \((T',x')\), \((T'\subset S^*)\). The dimension of (10.7.9) is that of \[
\sum_{T'\subset S^*}
\frac1{|W_{S^*}(T')^F|}
\operatorname{Ind}_{T'^F}^{G^{*F}}(1)
\] (because \(|G^{*F}|=|G^F|\)), while that of \((10.7.2)\otimes\operatorname{St}_G\) is \[
\sum_{T'\subset S^*}
\frac{(-1)^{\sigma(S)-\sigma(T')}}{|W_{S^*}(T')^F|}
\operatorname{Ind}_{T'^F}^{G^F}(1)
\]\par
}

\EditionPage{58}{159}{%
and the same proof as above gives (10.7.7), (10.7.8).\par

We now show that (10.6) can be applied with \(M=\Gamma_G\) or \(\Delta_G\). We have \[
\begin{aligned}
\langle\Delta_G,\Delta_G\rangle
&=\frac1{|G^F|}(|Z^F|q^l)^2\#\{\text{regular unipotents in }G^F\}\\
&=\frac1{|G^F|}(|Z^F|q^l)^2\frac{|G^F|}{|Z^F|q^l}
=|Z^F|q^l\qquad\text{(by 9.20)}.
\end{aligned}
\] The analogous identity \(\langle\Gamma_G,\Gamma_G\rangle=|Z^F|q^l\) follows from R. Steinberg ([14]). Now, for any \((T,\theta)\), \(\langle\Delta_G,R_T^\theta\rangle\) equals the average value of the character of \(R_T^\theta\) on the regular unipotents in \(G^F\) hence equals 1, by (9.16).\par

We now prove that \(\langle\Gamma_G,R_T^\theta\rangle=(-1)^{\sigma(G)-\sigma(T)}\) for any \((T,\theta)\). Using a result of Rodier (cf. T. A. Springer, Caractères de groupes de Chevalley finis, Sém. Bourbaki 429, Fév. 1973) and the induction formula 8.2, we are reduced to the case where \(T\) is not contained in any proper \(F\)-stable parabolic subgroup of \(G\). Using (10.3.1) and the already known fact that \(\langle\Delta_G,R_T^\theta\rangle=1\), we see that it is enough to prove \[
\tag{10.7.8}
\left\langle
\operatorname{Ind}_{P(J)^F}^{G^F}(\Gamma_{L(J)}),R_T^\theta
\right\rangle=0,
\qquad J\ne I.
\] (Note that in this case \((-1)^{\sigma(G)-\sigma(T)}=(-1)^{|I|}\).) Assuming that the theorem is already proved for groups of dimension strictly smaller than that of \(G\) and, in particular, for \(L(J)\), we see from (10.7.5) and (8.2) that \(\operatorname{Ind}_{P(J)^F}^{G^F}(\Gamma_{L(J)})\) is a linear combination of terms of the form \(R_{T'}^{\theta'}\), with \(T'\subset P(J)\), \(FT'=T'\), so that (10.7.8) follows from the orthogonality formula (6.8).\par

Now (10.7.1), (10.7.2), (10.7.5), (10.7.6) follow directly from (10.6) applied to \(M=\Gamma_G\) or \(\Delta_G\), except for an indeterminacy of sign. To remove this indeterminacy of sign, we only have to check that the right hand sides of (10.7.1) and (10.7.2) have strictly positive dimension; but they are both of the form \(|G^F|/|S^{*F}|\) up to a power of \(q\), by the first part of the proof. The theorem is proved.\par

\noindent\textbf{COROLLARY 10.8.} For any \(\rho\in(G^F)^{\sim}\), the average value of the character of \(\rho\) on the regular unipotents in \(G^F\) equals 0, 1, or \(-1\). This value is \(\pm1\) if and only if \(\rho=\rho_x'\) for some \(x\in\mathfrak S\).\par

Indeed, this average value is just \(\langle\rho,\Delta_G\rangle\).\par

Remark 10.9. If we assume that \(p\) is a good prime for \(G\), then all regular unipotents in \(G^F\) fall in a single \(G^F\)-conjugacy class so that the character of \(\rho\) at any regular unipotent in \(G^F\) equals 0, 1, or \(-1\). This was first proved in the article by J. A. Green, G. I. Lehrer, and G. Lusztig (On the degrees of certain group characters, to appear in The Quarterly Journal of Mathematics).\par

\noindent\textbf{PROPOSITION 10.10.} Let \(x\in\mathfrak S\). Let \((T,\theta)\) be a maximally split pair in \(x\)\par
}

\EditionPage{59}{160}{%
(see (5.17); \((T,\theta)\) is uniquely defined up to \(G^F\)-conjugacy by (5.18)). Then the virtual representation \((-1)^{\sigma(G)-\sigma(T)}R_T^\theta\) of \(G^F\) can be represented by an actual \(G^F\)-module, in which \(\rho_x\) and \(\rho_x'\) occur with multiplicity 1.\par

By (5.23), there exists an \(F\)-stable parabolic subgroup \(P\subset G\) such that \(T\) is contained in some \(F\)-stable Levi subgroup \(L\) of \(P\) and \((T,\theta)\) is non-singular in \(L\). \(L\) has connected centre, hence, by 5.20, \((T,\theta)\) is non-degenerate in \(L\). Let \(R_{T,L}^\theta\) be the virtual representation corresponding to \((T,\theta)\) with respect to \(L\) (see 1.20). By 7.4, \[
(-1)^{\sigma(L)-\sigma(T)}R_{T,L}^\theta
=
(-1)^{\sigma(G)-\sigma(T)}R_{T,L}^\theta
\] is irreducible. By 8.2, \[
(-1)^{\sigma(G)-\sigma(T)}R_T^\theta
=
\operatorname{Ind}_{P^F}^{G^F}\bigl((-1)^{\sigma(G)-\sigma(T)}R_{T,L}^\theta\bigr),
\] hence it can be realized as an actual \(G^F\)-module.\par

It remains to observe that for any \((T,\theta)\) in \(x\) (not necessarily maximally split) we have (by (10.7.1), (10.7.2), and (6.8)): \[
\langle\rho_x,R_T^\theta\rangle=(-1)^{\sigma(G)-\sigma(T)},
\]  \[
\langle\rho_x',R_T^\theta\rangle=(-1)^{\sigma(G)-\delta_x}.
\]\par

\EditionChapter{11. Suzuki and Ree groups}

Our methods apply also to the Suzuki and Ree groups \({}^2B_2(q)\), \({}^2G_2(q)\), \({}^2F_4(q)\) (with \(q\) an odd power of \(\sqrt2\), resp. \(\sqrt3\), \(\sqrt2\)).\par

These groups are the fixed point set of an exceptional isogeny \(F':G\to G\), with \(G\) respectively of type \(B_2\), \(G_2\) and \(F_4\), and \(F'^2\) being the Frobenius endomorphism relative to a rational structure over the field with \(q^2\) elements (R. Steinberg [15]). For the sake of induction, one should more generally consider groups of the form \(G^{F'}\), with \(G\) reductive and \(F'^2\) the Frobenius endomorphism relative to a rational structure over \(\mathbf F_{q^2}\) (\(q\) an odd power of \(\sqrt2\) or \(\sqrt3\)). Besides the Suzuki and Ree groups, this includes for instance groups of the form \(G=G_1\times G_1\), with \(F'(x,y)=(Fy,x)\); for such a group, \(G^{F'}=G_1^F\).\par

All definitions and results in Chapter 1 can be given with \(F\) replaced by \(F'\) throughout; in particular, for any \(F'\)-stable maximal torus \(T\) in \(G\) and any Borel subgroup \(B\) containing \(T\), the subscheme \(X_{T\subset B}\) of \(X_G\) and the \(G^{F'}\)-equivariant \(T^{F'}\)-torsor \(\widetilde X_{T\subset B}\) over \(X_{T\subset B}\) are well-defined (1.17). Thus for any \(\theta\in(T^{F'})^{\sim}\), the virtual representation \(R_{T\subset B}^\theta\) of \(G^{F'}\) is well-defined (1.20). For \(G=G_1\times G_1\), \(F'(x,y)=(Fy,x)\), this coincides with the similar representation of \(G_1^F=G^{F'}\).\par

All the results in Chapters 4, 5, 9 (except (9.18), (9.19)), the disjointness\par
}

\EditionPage{60}{161}{%
Theorem 6.2 (with its Corollary 6.3), the results (7.13) and (8.2) continue to be true. On the other hand, the proofs of the orthogonality relations (6.8), (6.9), and of the dimension formula (7.1) require the assumption \(|T^{F'}|>1\). If \(q>2\) this is automatically satisfied: \(|T^{F'}|\) is of the form \(\prod(q-\rho_i)\), where \(|\rho_i|=1\), hence \(|T^{F'}|>\prod(q-1)\). Thus, if we exclude the case where \(q\) equals \(\sqrt2\) or \(\sqrt3\), all the results in Chapters 6, 7, 8, 10 hold, as well as (9.18), (9.19). They can all be deduced from the orthogonality relations (6.8), (6.9). Even in the case \(q=\sqrt2\) or \(\sqrt3\), we can prove them for \({}^2B_2\) and \({}^2G_2\), but we cannot handle \({}^2F_4(\sqrt2)\).\par

\medskip\noindent\textbf{ADDED IN PROOF (2nd Dec. 1975). We can now handle \({}^2F_4(\sqrt2)\), too; this will be considered by one of us, in a future article.}\par

\medskip\noindent\textsc{INSTITUT DES HAUTES ÉTUDES SCIENTIFIQUES, FRANCE}\par

\medskip\noindent\textsc{MATHEMATICS INSTITUTE, UNIVERSITY OF WARWICK, ENGLAND}\par

\EditionHeading{REFERENCES}

\noindent\hangindent=1.7em\hangafter=1 [1] N. Bourbaki, Groupes et Algèbres de Lie, Chap. 4.5 et 6, Hermann, Paris 1968.\par

\noindent\hangindent=1.7em\hangafter=1 [2] B. Chang, R. Ree, The characters of \(G_2(q)\), Instituto Nazionale di Alta Mat., Symp. Mat. Vol. XIII (1974), 395--413.\par

\noindent\hangindent=1.7em\hangafter=1 [3] M. Demazure, Désingularisation des variétés de Schubert généralisées, Annales Sci. E.N.S., t. 7 (1974), 53--88.\par

\noindent\hangindent=1.7em\hangafter=1 [4] J. A. Green, The characters of the finite general linear groups, Trans. A.M.S. 80 (1955), 402--447.\par

\noindent\hangindent=1.7em\hangafter=1 [5] A. Grothendieck, Formule de Lefschetz et rationalité des Fonctions \(L\), Séminaire Bourbaki, 279, Décembre 1964 (Benjamin).\par

\noindent\hangindent=1.7em\hangafter=1 [6] H. C. Hansen, On Cycles on Flag Manifolds, Math. Scand. 33 (1973), 269--274.\par

\noindent\hangindent=1.7em\hangafter=1 [7] G. Lusztig, On the discrete series representations of the classical groups over finite fields. (Proc. of I.C.M., Vancouver 1974).\par

\noindent\hangindent=1.7em\hangafter=1 [8] ---------, Sur la conjecture de Macdonald (C. R. Acad. Sci. Paris, t. 280, (10 Fev. 1975), 317--320.\par

\noindent\hangindent=1.7em\hangafter=1 [9] I. G. Macdonald, Principal parts and the degeneracy rule (unpublished manuscript).\par

\noindent\hangindent=1.7em\hangafter=1 [10] Théorie des topos et cohomologie étale des schémas, Séminaire de Géométrie Algébrique du Bois Marie 4 (dirigé par M. Artin, A. Grothendieck, J. L. Verdier) I.H.E.S., 1963--64, Springer Lecture Notes 269, 270, 305.\par

\noindent\hangindent=1.7em\hangafter=1 [11] Cohomologie \(l\)-adique et fonctions \(L\), Séminaire de Géométrie Algébrique du Bois Marie 5; I.H.E.S., Notes polycopiées.\par

\noindent\hangindent=1.7em\hangafter=1 [12] T. A. Springer, Generalisation of Green's Polynomials, Proc. of Symp. in Pure Math. Vol. XXI (1971), A.M.S., 149--153.\par

\noindent\hangindent=1.7em\hangafter=1 [13] B. Srinivasan, The characters of the finite symplectic group \(\operatorname{Sp}(4,q)\), Trans. A.M.S. 131 (1968), 488--525.\par

\noindent\hangindent=1.7em\hangafter=1 [14] R. Steinberg, Lectures on Chevalley groups, Yale University, 1967.\par

\noindent\hangindent=1.7em\hangafter=1 [15] ---------, Endomorphisms of linear algebraic groups, Mem. of A.M.S. 80, 1968.\par

\medskip\noindent\itshape (Received June 6, 1975)\par\normalfont
}

\end{document}
